% Équations différentielles — Mathématiques Tle Tech — Cours signé OnBuch+
% Programme officiel MINESEC. Compiler avec : tectonic N15-equations-differentielles.tex
\newcommand{\DOCMATIERE}{Mathématiques}
\newcommand{\DOCNIVEAU}{Tle Tech}
\newcommand{\DOCMODULE}{Mathématiques – classe de Terminale technique}
\newcommand{\DOCLECON}{N15}
\newcommand{\DOCTITRE}{Équations différentielles}
% OnBuch+ — préambule « cours signés » ENRICHI (compilé avec tectonic / XeLaTeX)
% Système visuel « Pop » d'OnBuch+ : fond crème (jamais blanc), bordures encre
% épaisses, ombres « dures » décalées, titres Archivo Black en capitales.
% Version MATHÉMATIQUES : graphes (pgfplots/tikz), boîtes pédagogiques
% (définition, propriété, méthode, attention, le savais-tu, exercice résolu,
% exercice, TP, bilan), page de garde et signature OnBuch+ ancrée en bas.
%
% Macros attendues dans main.tex AVANT \input{preamble} :
%   \DOCMATIERE  \DOCNIVEAU  \DOCMODULE  \DOCLECON (numéro)  \DOCTITRE
\documentclass[11pt]{article}

\usepackage{fontspec}
\usepackage[french]{babel}
\usepackage[a4paper,margin=2cm,top=3.4cm,bottom=2.8cm,headheight=1.4cm,footskip=1.3cm]{geometry}
\usepackage{amsmath,amssymb,amsfonts}
\usepackage{mathtools} % \xmapsto, \coloneqq... souvent utilisés par les modèles
\usepackage{cancel} % simplifications barrées
% \abs{z} (module, valeur absolue) et \norm{u} (norme), utilisés par les modèles
\providecommand{\abs}{}\let\abs\relax\DeclarePairedDelimiter{\abs}{\lvert}{\rvert}
\providecommand{\norm}{}\let\norm\relax\DeclarePairedDelimiter{\norm}{\lVert}{\rVert}
\usepackage[table]{xcolor} % [table] : fournit aussi \rowcolors (alternance de lignes)
\usepackage{pagecolor}
\usepackage{tikz}
\usetikzlibrary{arrows.meta,positioning,calc,shapes.geometric,shapes.misc,shapes.arrows,decorations.pathmorphing,decorations.pathreplacing,decorations.markings,patterns,shadows,mindmap,trees,fit,backgrounds,matrix,intersections,angles,quotes}
\usepackage{pgfplots}
\usepackage{pgf-pie} % diagrammes circulaires \pie (statistiques)
\pgfplotsset{compat=1.18}
\usepgfplotslibrary{fillbetween,groupplots} % aires entre courbes (calcul intégral)
\usepackage{enumitem}
\usepackage{fancyhdr}
% [most] charge toutes les bibliothèques tcolorbox (listings, minted, raster,
% documentation...) et gonfle inutilement la mémoire TeX sur un document long
% (~300 boîtes) : on ne charge que ce qui est réellement utilisé.
\usepackage[skins,breakable]{tcolorbox}
\usepackage{graphicx}
\usepackage{colortbl}
\usepackage{booktabs}
\usepackage{tabularx}
\usepackage{diagbox} % case d'en-tête diagonale des tableaux à double entrée
% Colonnes extensibles centrée / gauche / droite (C, L, R), souvent utilisées par les modèles
\newcolumntype{C}{>{\centering\arraybackslash}X}
\newcolumntype{L}{>{\raggedright\arraybackslash}X}
\newcolumntype{R}{>{\raggedleft\arraybackslash}X}
\usepackage{array}
\usepackage{multicol}
\usepackage{siunitx}
% parse-numbers=false : accepte aussi \SI{2,5 \times 10^{-3}}{...}
\sisetup{locale=FR,output-decimal-marker={,},group-minimum-digits=4,parse-numbers=false}
\DeclareSIUnit\ohm{\ensuremath{\symup{\Omega}}} % U+2126 absent de la police
\usepackage{qrcode}
% \degree : commande fréquemment utilisée par les modèles pour un angle ou une
% température en dehors de \SI{}{} (non fournie par les paquets chargés).
\providecommand{\degree}{^{\circ}}
% \qed (fin de démonstration), utilisé par les modèles sans amsthm
\providecommand{\qed}{\hfill$\square$}
% \barre : double barre des valeurs interdites dans un tableau de variations (tkz-tab)
\providecommand{\barre}{\ensuremath{\|}}
% environnement proof (amsthm non chargé) utilisé par les modèles
\ifdefined\proof\else\newenvironment{proof}[1][Démonstration]{\par\noindent\textit{#1.}\ }{\qed\par}\fi
\usepackage{lastpage}
\usepackage{hyperref}
\hypersetup{hidelinks}

% ============================================================
% Palette « Pop » OnBuch+ — mêmes hex que la classe P (plus_ui.dart)
% ============================================================
\definecolor{poporange}{HTML}{F59321}
\definecolor{popdark}{HTML}{E07A0C}
\definecolor{popcream}{HTML}{FFF6E8}
\definecolor{popink}{HTML}{1C1714}
\definecolor{poppurple}{HTML}{7A5AE0}
\definecolor{popgreen}{HTML}{2E9E6A}
\definecolor{poppink}{HTML}{E0457B}
\definecolor{popblue}{HTML}{2F7DE1}
\definecolor{popgold}{HTML}{C98A12}
\definecolor{popmuted}{HTML}{8A7F76}
\definecolor{popline}{HTML}{EADFD2}
\definecolor{popgray}{HTML}{8A7F76}
\colorlet{popgrey}{popgray}
% Alias tolérants (le modèle nomme parfois les couleurs différemment)
\colorlet{popwhite}{white}
\colorlet{popblack}{popink}
\colorlet{popred}{poppink}
\colorlet{popyellow}{popgold}
% Teintes claires pour fonds de tableaux / zones
\colorlet{popblueL}{popblue!10!white}
\colorlet{poporangeL}{poporange!14!white}
\colorlet{popgreenL}{popgreen!12!white}
\colorlet{poppurpleL}{poppurple!12!white}
\colorlet{poppinkL}{poppink!12!white}
\colorlet{popgoldL}{popgold!14!white}

\pagecolor{popcream}

% ============================================================
% Typographie
% ============================================================
\defaultfontfeatures{Ligatures=TeX}
\setmainfont{PlusJakartaSans-Medium.ttf}[Path=../../../fonts/,BoldFont=PlusJakartaSans-Bold.ttf]
% Maths en sans-serif (Fira Math) avec chiffres et lettres droites de Plus
% Jakarta, pour que les formules se fondent dans le texte.
\usepackage[math-style=ISO,bold-style=ISO]{unicode-math}
\setmathfont{FiraMath-Regular.otf}[Path=../../../fonts/]
\setmathfont{PlusJakartaSans-Medium.ttf}[Path=../../../fonts/,range={up/num,up/{Latin,latin},\mathup}]
\usepackage{icomma} % virgule décimale sans espace en mode math
% Caractères mathématiques Unicode absents des polices de texte (Plus Jakarta,
% Archivo Black) : rendus par leur équivalent mathématique, en texte comme en
% formule — sinon ils s'affichent en carré vide (ex. « Arithmétique dans ℤ »).
\usepackage{newunicodechar}
\newunicodechar{ℝ}{\ensuremath{\mathbb{R}}}
\newunicodechar{ℤ}{\ensuremath{\mathbb{Z}}}
\newunicodechar{ℂ}{\ensuremath{\mathbb{C}}}
\newunicodechar{ℕ}{\ensuremath{\mathbb{N}}}
\newunicodechar{ℚ}{\ensuremath{\mathbb{Q}}}
\newunicodechar{𝔻}{\ensuremath{\mathbb{D}}}
\newunicodechar{α}{\ensuremath{\alpha}}
\newunicodechar{β}{\ensuremath{\beta}}
\newunicodechar{γ}{\ensuremath{\gamma}}
\newunicodechar{δ}{\ensuremath{\delta}}
\newunicodechar{ε}{\ensuremath{\varepsilon}}
\newunicodechar{θ}{\ensuremath{\theta}}
\newunicodechar{λ}{\ensuremath{\lambda}}
\newunicodechar{μ}{\ensuremath{\mu}}
\newunicodechar{σ}{\ensuremath{\sigma}}
\newunicodechar{τ}{\ensuremath{\tau}}
\newunicodechar{φ}{\ensuremath{\varphi}}
\newunicodechar{ψ}{\ensuremath{\psi}}
\newunicodechar{ω}{\ensuremath{\omega}}
\newunicodechar{Σ}{\ensuremath{\Sigma}}
% majuscules produites par \MakeUppercase dans les titres (α-aminés -> Α)
\newunicodechar{Α}{\ensuremath{\alpha}}
\newunicodechar{Β}{\ensuremath{\beta}}
\newunicodechar{Γ}{\ensuremath{\Gamma}}
\newunicodechar{Δ}{\ensuremath{\Delta}}
\newunicodechar{Ε}{\ensuremath{\varepsilon}}
\newunicodechar{Θ}{\ensuremath{\Theta}}
\newunicodechar{Λ}{\ensuremath{\Lambda}}
\newunicodechar{Μ}{\ensuremath{\mu}}
\newunicodechar{Τ}{\ensuremath{\tau}}
\newunicodechar{Φ}{\ensuremath{\Phi}}
\newunicodechar{Ψ}{\ensuremath{\Psi}}
\newunicodechar{Ω}{\ensuremath{\Omega}}
\newunicodechar{↦}{\ensuremath{\mapsto}}
\newunicodechar{∈}{\ensuremath{\in}}
\newunicodechar{∉}{\ensuremath{\notin}}
\newunicodechar{∧}{\ensuremath{\wedge}}
\newunicodechar{⇔}{\ensuremath{\Leftrightarrow}}
\newunicodechar{⟺}{\ensuremath{\Longleftrightarrow}}
\newunicodechar{⇒}{\ensuremath{\Rightarrow}}
\newunicodechar{⟹}{\ensuremath{\Longrightarrow}}
\newunicodechar{∘}{\ensuremath{\circ}}
\newunicodechar{∪}{\ensuremath{\cup}}
\newunicodechar{∩}{\ensuremath{\cap}}
\newunicodechar{≡}{\ensuremath{\equiv}}
\newunicodechar{⊂}{\ensuremath{\subset}}
\newunicodechar{⊥}{\ensuremath{\perp}}
\newunicodechar{∃}{\ensuremath{\exists}}
\newunicodechar{∀}{\ensuremath{\forall}}
\newunicodechar{∛}{\ensuremath{\sqrt[3]{\,}}}
\newunicodechar{∜}{\ensuremath{\sqrt[4]{\,}}}
\newunicodechar{⃗}{\ensuremath{{}^{\to}}}
\newunicodechar{ⁿ}{\ifmmode^{n}\else\textsuperscript{\ensuremath{n}}\fi}
\newunicodechar{ˣ}{\ifmmode^{x}\else\textsuperscript{\ensuremath{x}}\fi}
\newunicodechar{ᵇ}{\ifmmode^{b}\else\textsuperscript{\ensuremath{b}}\fi}
\newunicodechar{ʳ}{\ifmmode^{r}\else\textsuperscript{\ensuremath{r}}\fi}
\newunicodechar{ᵘ}{\ifmmode^{u}\else\textsuperscript{\ensuremath{u}}\fi}
\newunicodechar{ʸ}{\ifmmode^{y}\else\textsuperscript{\ensuremath{y}}\fi}
\newunicodechar{ᵃ}{\ifmmode^{a}\else\textsuperscript{\ensuremath{a}}\fi}
\newunicodechar{ᵉ}{\ifmmode^{e}\else\textsuperscript{\ensuremath{e}}\fi}
\newunicodechar{ᵏ}{\ifmmode^{k}\else\textsuperscript{\ensuremath{k}}\fi}
\newunicodechar{ᵖ}{\ifmmode^{p}\else\textsuperscript{\ensuremath{p}}\fi}
\newunicodechar{ᴺ}{\ifmmode^{N}\else\textsuperscript{\ensuremath{N}}\fi}
\newunicodechar{ᶜ}{\ifmmode^{c}\else\textsuperscript{\ensuremath{c}}\fi}
\newunicodechar{ʰ}{\ifmmode^{h}\else\textsuperscript{\ensuremath{h}}\fi}
\newunicodechar{ᵠ}{\ifmmode^{q}\else\textsuperscript{\ensuremath{q}}\fi}
\newunicodechar{ₙ}{\ifmmode_{n}\else\textsubscript{\ensuremath{n}}\fi}
\newunicodechar{ₐ}{\ifmmode_{a}\else\textsubscript{\ensuremath{a}}\fi}
\newunicodechar{ᵢ}{\ifmmode_{i}\else\textsubscript{\ensuremath{i}}\fi}
\newunicodechar{ᵣ}{\ifmmode_{r}\else\textsubscript{\ensuremath{r}}\fi}
\newunicodechar{ₖ}{\ifmmode_{k}\else\textsubscript{\ensuremath{k}}\fi}
\newunicodechar{ᵦ}{\ifmmode_{\beta}\else\textsubscript{\ensuremath{\beta}}\fi}
\newunicodechar{ₓ}{\ifmmode_{x}\else\textsubscript{\ensuremath{x}}\fi}
\newfontfamily\popdisplay{ArchivoBlack-Regular.ttf}[Path=../../../fonts/]
\newfontfamily\poplogo{SpaceGrotesk-Bold.ttf}[Path=../../../fonts/]
\newfontfamily\popsemi{PlusJakartaSans-SemiBold.ttf}[Path=../../../fonts/]
\newfontfamily\popxbold{PlusJakartaSans-ExtraBold.ttf}[Path=../../../fonts/]
\newfontfamily\popxboldls{PlusJakartaSans-ExtraBold.ttf}[Path=../../../fonts/,LetterSpace=3.0]

\setlist[enumerate]{leftmargin=1.7em,itemsep=5pt,topsep=4pt}
\setlist[itemize]{leftmargin=1.7em,itemsep=4pt,topsep=4pt,label={\color{poporange}\textbullet}}
\setlength{\parindent}{0pt}
\setlength{\parskip}{5pt}
\renewcommand{\arraystretch}{1.25}

% pgfplots : style Pop par défaut
\pgfplotsset{
  every axis/.append style={axis line style={popink,line width=1pt},tick style={popink},
    grid style={popline,line width=0.5pt},label style={font=\small},tick label style={font=\footnotesize}},
  popcurve/.style={poporange,line width=1.8pt,no markers,smooth},
}
% Même style disponible dans un \draw TikZ simple (hors pgfplots)
\tikzset{popcurve/.style={poporange,line width=1.8pt}}

% ============================================================
% Pill / squiggle
% ============================================================
\newcommand{\poppill}[3]{%
  \tikz[baseline=-0.55ex]\node[draw=popink,line width=1.2pt,rounded corners=2.6pt,%
    fill=#2,inner xsep=6.5pt,inner ysep=3pt]{%
    \popxboldls\fontsize{7}{7}\selectfont\color{#3}\MakeUppercase{#1}};%
}
\newcommand{\popsquiggle}[1]{%
  \tikz[baseline=-0.4ex]{
    \draw[#1,line width=2.6pt,line cap=round]
      plot [smooth] coordinates {(0,0.09) (0.34,0.01) (0.68,0.21) (1.02,0.01) (1.36,0.10)};
  }%
}
% Mot-clé surligné (marqueur orange sous le texte)
\newcommand{\cle}[1]{{\popxbold\color{popdark}#1}}

% ============================================================
% En-tête / pied
% ============================================================
\pagestyle{fancy}
\fancyhf{}
\renewcommand{\headrulewidth}{0pt}
\renewcommand{\footrulewidth}{0pt}
\lhead{\raisebox{-0.15ex}{\poplogo\fontsize{13}{13}\selectfont\color{popink}OnBuch\color{poporange}+}}
\rhead{\poppill{\DOCMATIERE\ \textbullet\ \DOCNIVEAU\ \textbullet\ Leçon \DOCLECON}{white}{popink}}
\lfoot{\popxbold\fontsize{7.6}{7.6}\selectfont\color{popmuted}\MakeUppercase{Cours signé OnBuch+}\ \textbullet\ {\popsemi\DOCTITRE}}
\rfoot{\tikz[baseline=-0.6ex]\node[rounded corners=8.5pt,fill=popink,minimum height=17pt,minimum width=17pt,inner xsep=5pt,inner sep=0]{\popxbold\fontsize{8}{8}\selectfont\color{popcream}\thepage\,/\,\pageref*{LastPage}};}

% ============================================================
% Titres
% ============================================================
\newcommand{\chaptitre}[2]{%
  \begin{tcolorbox}[enhanced,colback=poporange,colframe=popink,boxrule=2.6pt,arc=11pt,
      drop shadow=popink,left=15pt,right=15pt,top=12pt,bottom=11pt,before skip=3pt,after skip=18pt]
    \poppill{#2}{popink}{popcream}\par\vspace{9pt}
    {\raggedright\hyphenpenalty=10000\exhyphenpenalty=10000\popdisplay\fontsize{22}{24}\selectfont\color{popcream}\MakeUppercase{#1}\par}
  \end{tcolorbox}
}
% Section numérotée : pastille orange avec le numéro + titre Archivo
\newcounter{popsec}
\newcommand{\coursec}[1]{%
  \stepcounter{popsec}\setcounter{popsub}{0}%
  \par\vspace{16pt}\noindent
  \begin{minipage}{\linewidth}
  \tikz[baseline=-0.7ex]\node[draw=popink,line width=1.6pt,fill=poporange,rounded corners=4pt,minimum size=22pt,inner sep=0]{\popdisplay\fontsize{12}{12}\selectfont\color{popcream}\thepopsec};%
  \hspace{8pt}{\popdisplay\fontsize{15}{17}\selectfont\color{popink}\MakeUppercase{#1}}\par
  \vspace{4pt}\noindent{\color{poporange}\rule{36pt}{3.6pt}}
  \end{minipage}\par\nopagebreak\vspace{8pt}
}
\newcounter{popsub}
\newcommand{\courssub}[1]{%
  \stepcounter{popsub}%
  \par\vspace{9pt}\noindent{\popxbold\fontsize{11.5}{13}\selectfont\color{popink}{\color{poporange}\thepopsec.\thepopsub}\hspace{6pt}#1}\par\nopagebreak\vspace{4pt}
}

% ============================================================
% Boîtes pédagogiques — même recette : fond blanc, bordure épaisse colorée,
% ombre dure encre, badge plein. Titre optionnel [#1].
% ============================================================
\tcbset{popbase/.style={breakable,enhanced,colback=white,boxrule=2.3pt,arc=9pt,
  shadow={3pt}{-3pt}{0pt}{popink},left=14pt,right=14pt,top=11pt,bottom=11pt,before skip=11pt,after skip=15pt,
  fontupper=\popsemi\fontsize{10.6}{14.5}\selectfont\color{popink},
  fontlower=\popsemi\fontsize{10.6}{14.5}\selectfont\color{popink}}}
% Coche dessinée (le glyphe ✓ n'existe pas dans les polices du document)
\newcommand{\popcheck}[1]{\tikz[baseline=-0.5ex]\draw[#1,line width=1.6pt,line cap=round,line join=round](-0.13,0.0)--(-0.04,-0.09)--(0.13,0.1);}
\newcommand{\popboxhead}[3]{\poppill{#1}{#2}{white}\ifx\relax#3\relax\else\hspace{6pt}{\popxbold\fontsize{10.6}{12}\selectfont\color{popink}#3}\fi\par\vspace{7pt}}

\newtcolorbox{aretenir}[1][]{popbase,colframe=popgreen,before upper={\popboxhead{À retenir}{popgreen}{#1}}}
\newtcolorbox{exemplebox}[1][]{popbase,colframe=poporange,before upper={\popboxhead{Exemple}{poporange}{#1}}}
\newtcolorbox{definition}[1][]{popbase,colframe=popblue,before upper={\popboxhead{Définition}{popblue}{#1}}}
\newtcolorbox{propriete}[1][]{popbase,colframe=poppurple,before upper={\popboxhead{Propriété}{poppurple}{#1}}}
\newtcolorbox{methode}[1][]{popbase,colframe=popgold,before upper={\popboxhead{Méthode}{popgold}{#1}}}
\newtcolorbox{attention}[1][]{popbase,colframe=poppink,before upper={\popboxhead{Attention}{poppink}{#1}}}
\newtcolorbox{experience}[1][]{popbase,colframe=popblue,colback=popblueL,before upper={\popboxhead{Expérience}{popblue}{#1}}}
% « Le savais-tu ? » : carte violette pleine (contexte, culture, Cameroun)
\newtcolorbox{savaistu}[1][]{popbase,colback=poppurple,colframe=popink,
  fontupper=\popsemi\fontsize{10.4}{14.2}\selectfont\color{white},
  before upper={\poppill{Le savais-tu\,?}{white}{poppurple}\ifx\relax#1\relax\else\hspace{6pt}{\popxbold\color{white}#1}\fi\par\vspace{7pt}}}
% Exercice résolu : énoncé en haut, \tcblower puis la solution sur fond crème
\newcounter{popexo}\newcounter{popexr}
\newtcolorbox{exoresolu}[1][]{popbase,colframe=poporange,colbacklower=popcream,
  segmentation style={popink,line width=1.2pt,dashed},
  before upper={\stepcounter{popexr}\popboxhead{Exercice résolu \thepopexr}{poporange}{#1}},
  before lower={\poppill{Solution}{popink}{popcream}\par\vspace{6pt}}}
% Exercice (non résolu en place) — la correction est dans la partie Corrigés
\newtcolorbox{exercice}[1][]{popbase,colframe=popink,
  before upper={\stepcounter{popexo}\popboxhead{Exercice \thepopexo}{popink}{#1}}}
% Corrigé
\newtcolorbox{corrige}[1][]{popbase,colframe=popgreen,colback=popgreenL,
  before upper={\popboxhead{Corrigé}{popgreen}{#1}}}
% Objectifs / prérequis / bilan
\newtcolorbox{objectifs}{popbase,colframe=popink,colback=poporangeL,
  before upper={\popboxhead{Objectifs de la leçon}{popink}{}}}
\newtcolorbox{prerequis}{popbase,colframe=popink,colback=popblueL,
  before upper={\popboxhead{Prérequis}{popblue}{}}}
\newtcolorbox{bilan}{popbase,colframe=popink,colback=white,boxrule=2.8pt,
  before upper={\popboxhead{Fiche bilan}{poporange}{L'essentiel à connaître pour l'examen}}}
% Figure encadrée avec légende
\newtcolorbox{popfigure}[1][]{enhanced,breakable=false,colback=white,colframe=popline,boxrule=1.4pt,arc=8pt,
  left=10pt,right=10pt,top=10pt,bottom=8pt,before skip=10pt,after skip=12pt,halign=center,
  lower separated=false,halign lower=center,
  fontlower=\popsemi\fontsize{9}{11.5}\selectfont\color{popmuted},
  #1}
\newcommand{\legende}[1]{\par\vspace{6pt}{\popsemi\fontsize{9}{11.5}\selectfont\color{popmuted}#1}}

% ============================================================
% Signature OnBuch+
% ============================================================
\newcommand{\poplogoinline}[1]{{\poplogo\fontsize{#1}{#1}\selectfont\color{popink}OnBuch\color{poporange}+}}
\newcommand{\signatureonbuch}{%
  \begin{tcolorbox}[enhanced,colback=popink,colframe=popink,boxrule=2.2pt,arc=10pt,
      drop shadow=poporange,left=14pt,right=14pt,top=10pt,bottom=10pt,before skip=0pt,after skip=0pt]
    \tikz[baseline=-0.7ex]\node[circle,fill=popgreen,draw=popcream,line width=1.3pt,minimum size=22pt,inner sep=0]{\popcheck{white}};%
    \hspace{9pt}\begin{minipage}[c]{0.62\linewidth}
      {\popxbold\fontsize{12}{13}\selectfont\color{popcream}Signé\ \textbullet\ {\poplogo OnBuch\color{poporange}+}}\par\vspace{2pt}
      {\raggedright\popsemi\fontsize{8}{10.5}\selectfont\color{popline}\MakeUppercase{Équipe pédagogique OnBuch+}\\\MakeUppercase{Conforme au programme officiel MINESEC}\par}
    \end{minipage}\hfill
    {\poplogo\fontsize{22}{22}\selectfont\color{popcream}OnBuch\color{poporange}+}
  \end{tcolorbox}%
}

% ============================================================
% Page de garde
% #1 titre · #2 matière · #3 niveau · #4 module · #5 n° leçon · #6 durée ·
% #7 description / objectifs (texte ou itemize)
% ============================================================
\newcommand{\pagegarde}[7]{%
  \thispagestyle{empty}
  \newgeometry{a4paper,margin=1.7cm,top=1.6cm,bottom=1.4cm}%
  \begin{tikzpicture}[remember picture,overlay]
    \fill[poporange!18] ([xshift=-3.2cm,yshift=-3.6cm]current page.north east) circle (4.6cm);
    \fill[poppurple!14] ([xshift=2.4cm,yshift=7.6cm]current page.south west) circle (3.8cm);
  \end{tikzpicture}%
  {\poplogo\fontsize{30}{30}\selectfont\color{popink}OnBuch\color{poporange}+}\hfill
  \raisebox{0.6ex}{\poppill{Cours signé \textbullet\ Premium}{popink}{popcream}}\par
  \vspace{1pt}\popsquiggle{poppurple}\par
  \vspace{8pt}
  \begin{tcolorbox}[enhanced,colback=poporange,colframe=popink,boxrule=2.8pt,arc=13pt,
      drop shadow=popink,left=18pt,right=18pt,top=12pt,bottom=13pt,after skip=10pt]
    \poppill{#2 \textbullet\ #3}{popink}{popcream}\hspace{5pt}\poppill{#4}{white}{popink}\par\vspace{8pt}
    {\popdisplay\fontsize{14}{14}\selectfont\color{popink}LEÇON #5}\par\vspace{4pt}
    {\raggedright\hyphenpenalty=10000\exhyphenpenalty=10000\popdisplay\fontsize{25}{27}\selectfont\color{popcream}\MakeUppercase{#1}\par}
    \vspace{8pt}
    {\popxbold\fontsize{9.5}{9.5}\selectfont\color{popink}\MakeUppercase{Durée conseillée : #6}}
  \end{tcolorbox}
  \begin{tcolorbox}[enhanced,colback=white,colframe=popink,boxrule=2.2pt,arc=10pt,
      drop shadow=popink,left=16pt,right=16pt,top=10pt,bottom=10pt,after skip=8pt]
    \poppill{Dans cette leçon}{poppurple}{white}\par\vspace{5pt}
    \popsemi\fontsize{9.3}{12.6}\selectfont\color{popink}#7
  \end{tcolorbox}
  \noindent\begin{minipage}[t]{0.487\linewidth}
    \begin{tcolorbox}[enhanced,colback=popgreen,colframe=popink,boxrule=2.2pt,arc=10pt,
        drop shadow=popink,left=11pt,right=11pt,top=10pt,bottom=10pt,height=3.1cm,valign=center]
      \tcbox[colback=white,colframe=popink,boxrule=1.3pt,arc=5pt,boxsep=0pt,nobeforeafter,
        left=3pt,right=3pt,top=3pt,bottom=3pt,valign=center]{\qrcode[height=1.9cm]{https://play.google.com/store/apps/details?id=com.luvvix.onbuch&hl=fr}}%
      \hspace{8pt}\begin{minipage}[c]{0.5\linewidth}
        {\popdisplay\fontsize{11}{12.5}\selectfont\color{white}\MakeUppercase{Télécharge OnBuch}}\par\vspace{4pt}
        {\raggedright\popsemi\fontsize{8.4}{11}\selectfont\color{white}Gratuit \textbullet\ Google Play\\Tuteur IA, annales, concours, résultats\par}
      \end{minipage}
    \end{tcolorbox}
  \end{minipage}\hfill
  \begin{minipage}[t]{0.487\linewidth}
    \begin{tcolorbox}[enhanced,colback=poppurple,colframe=popink,boxrule=2.2pt,arc=10pt,
        drop shadow=popink,left=11pt,right=11pt,top=10pt,bottom=10pt,height=3.1cm,valign=center]
      \tikz[baseline=-0.7ex]\node[circle,fill=white,draw=popink,line width=1.3pt,minimum size=1.6cm,inner sep=0]{\popdisplay\fontsize{24}{24}\selectfont\color{poppurple}+};%
      \hspace{8pt}\begin{minipage}[c]{0.6\linewidth}
        {\popdisplay\fontsize{11}{12.5}\selectfont\color{white}\MakeUppercase{Passe à OnBuch+}}\par\vspace{4pt}
        {\raggedright\popsemi\fontsize{8.4}{11}\selectfont\color{white}Tous les cours signés, Tuteur IA illimité, fiches PDF — onglet OnBuch+ de l'app\par}
      \end{minipage}
    \end{tcolorbox}
  \end{minipage}
  \vspace{8pt}
  \popcontenu
  \vfill
  \signatureonbuch
  \restoregeometry
  \newpage
}

% Rangée « Ce cours contient » (page de garde)
\newcommand{\poptile}[3]{% #1 couleur #2 gros texte #3 légende
  \begin{tcolorbox}[enhanced,colback=#1,colframe=popink,boxrule=2pt,arc=9pt,shadow={2.5pt}{-2.5pt}{0pt}{popink},
      left=6pt,right=6pt,top=9pt,bottom=9pt,halign=center,valign=center,height=1.75cm,width=\linewidth,nobeforeafter]
    {\popdisplay\fontsize{12.5}{13}\selectfont\color{white}\MakeUppercase{#2}}\par\vspace{3pt}
    {\centering\popsemi\fontsize{7.8}{9.5}\selectfont\color{white}#3\par}
  \end{tcolorbox}}
\newcommand{\popcontenu}{%
  \par\noindent{\popxboldls\fontsize{7.5}{7.5}\selectfont\color{popmuted}CE COURS SIGNÉ CONTIENT}\par\vspace{7pt}
  \noindent\begin{minipage}[t]{0.235\linewidth}\poptile{popblue}{Cours}{complet, illustré, pas à pas}\end{minipage}\hfill
  \begin{minipage}[t]{0.235\linewidth}\poptile{poporange}{Méthodes}{et exercices résolus}\end{minipage}\hfill
  \begin{minipage}[t]{0.235\linewidth}\poptile{poppink}{Exercices}{type examen + corrigés}\end{minipage}\hfill
  \begin{minipage}[t]{0.235\linewidth}\poptile{popgreen}{Fiche bilan}{l'essentiel à retenir}\end{minipage}\par}

% Sommaire
\newtcolorbox{sommaire}{enhanced,colback=white,colframe=popink,boxrule=2.2pt,arc=10pt,
  drop shadow=popink,left=18pt,right=18pt,top=13pt,bottom=13pt,before skip=6pt,after skip=20pt,
  before upper={\poppill{Sommaire}{poppurple}{white}\par\vspace{9pt}\popsemi\fontsize{10.6}{14.5}\selectfont\color{popink}}}

% Signature de fin de cours — poussée tout en bas de la dernière page
\newcommand{\findecours}{%
  \par\vfill\nopagebreak
  \begin{center}\popsquiggle{poporange}\end{center}\vspace{4pt}
  \signatureonbuch
}
\AtBeginDocument{\def\llbracket{\mathopen{[\mkern-3.5mu[}}\def\rrbracket{\mathclose{]\mkern-3.5mu]}}} % intervalles d'entiers (glyphe ⟦ absent de la police)
% Glyphes absents de Fira Math : redéfinis à partir de glyphes disponibles
\AtBeginDocument{%
  \def\setminus{\mathbin{\backslash}}\let\smallsetminus\setminus
  \def\vdots{\mathord{\vbox{\baselineskip3.2pt\lineskiplimit0pt\kern4pt\hbox{.}\hbox{.}\hbox{.}}}}%
  \def\ddots{\mathinner{\mkern1mu\raise6.5pt\hbox{.}\mkern2mu\raise3.6pt\hbox{.}\mkern2mu\raise0.7pt\hbox{.}\mkern1mu}}%
  \def\triangle{\mathord{\scalebox{0.9}{$\symup{\Delta}$}}}%
  \def\nabla{\mathord{\raisebox{\depth}{\scalebox{1}[-1]{$\symup{\Delta}$}}}}%
  \def\checkmark{\popcheck{popdark}}%
  \def\star{\mathbin{\ast}}\def\bigstar{\mathord{\ast}}%
  \def\diamond{\mathbin{\scalebox{0.6}{$\Diamond$}}}%
  \def\bigcup{\mathop{\mathchoice{\raisebox{-0.3em}{\scalebox{1.7}{$\cup$}}}{\scalebox{1.25}{$\cup$}}{\cup}{\cup}}\displaylimits}%
  \def\bigcap{\mathop{\mathchoice{\raisebox{-0.3em}{\scalebox{1.7}{$\cap$}}}{\scalebox{1.25}{$\cap$}}{\cap}{\cap}}\displaylimits}%
  \def\lesseqqgtr{\mathrel{\lessgtr}}\def\bigoplus{\mathop{\oplus}\displaylimits}%
}
\providecommand{\ssi}{si et seulement si} % abréviation fréquente
\AtBeginDocument{\providecommand{\wideparen}{\overparen}} % arc de cercle

%% FIGFIX-BEGIN v5 — correctifs globaux des figures (docs/onbuch-generation/tools/figfix)
\usepackage{adjustbox}\usepackage{etoolbox}\tcbuselibrary{raster}
% toute figure TikZ/pgfplots plus large que la ligne est réduite à la largeur disponible
\BeforeBeginEnvironment{tikzpicture}{\begin{adjustbox}{max width=\linewidth}}
\AfterEndEnvironment{tikzpicture}{\end{adjustbox}}
% légendes pgfplots lisibles par-dessus les courbes
\pgfplotsset{every axis/.append style={legend style={fill=white,fill opacity=0.92,text opacity=1,draw=black!15,inner sep=3pt}}}
% étiquettes d'arêtes (syntaxe edge["..."]) avec fond blanc
\tikzset{every edge quotes/.append style={fill=white,inner sep=1.2pt}}
%% FIGFIX-END
\begin{document}

\pagegarde{Équations différentielles}{Mathématiques}{Tle Tech}{Mathématiques – classe de Terminale technique}{N15}{4 séances (fiche de progression harmonisée)}{Cette leçon introduit les équations différentielles, outil fondamental de modélisation des phénomènes continus (croissance, circuits électriques, amortissement). On y apprend à résoudre les équations y' = ay et les équations linéaires du second ordre à coefficients constants, avec ou sans second membre constant.
\vspace{4pt}
\begin{multicols}{2}\small
\begin{itemize}
\item À la fin de cette leçon, je sais reconnaître une équation différentielle et définir les notions de solution, d'ordre et de condition initiale
\item À la fin de cette leçon, je sais résoudre une équation différentielle du type y' = ay et déterminer la solution vérifiant une condition initiale
\item À la fin de cette leçon, je sais former l'équation caractéristique d'une équation ay'' + by' + cy = 0 et discuter selon son discriminant
\item À la fin de cette leçon, je sais donner la solution générale de ay'' + by' + cy = 0 dans les trois cas (Δ > 0, Δ = 0, Δ < 0)
\item À la fin de cette leçon, je sais déterminer une solution particulière constante de ay'' + by' + cy = d et en déduire la solution générale
\item À la fin de cette leçon, je sais déterminer l'unique solution d'un problème de Cauchy (conditions initiales y(t₀) et y'(t₀))
\end{itemize}\end{multicols}}

\begin{sommaire}
\begin{multicols}{2}
\begin{enumerate}
\item Présentation et vocabulaire
\item Équations du type y' = ay
\item Équations ay'' + by' + cy = 0 : équation caractéristique
\item Solutions générales de ay'' + by' + cy = 0
\item Équations ay'' + by' + cy = d (second membre constant)
\item Applications à des situations concrètes et synthèse
\item Activité d'intégration
\item Méthodes et exercices résolus
\item Exercices
\item Corrigés des exercices
\item Fiche bilan
\end{enumerate}
\end{multicols}
\end{sommaire}

\begin{objectifs}
\begin{itemize}
\item À la fin de cette leçon, je sais reconnaître une équation différentielle et définir les notions de solution, d'ordre et de condition initiale
\item À la fin de cette leçon, je sais résoudre une équation différentielle du type y' = ay et déterminer la solution vérifiant une condition initiale
\item À la fin de cette leçon, je sais former l'équation caractéristique d'une équation ay'' + by' + cy = 0 et discuter selon son discriminant
\item À la fin de cette leçon, je sais donner la solution générale de ay'' + by' + cy = 0 dans les trois cas (Δ > 0, Δ = 0, Δ < 0)
\item À la fin de cette leçon, je sais déterminer une solution particulière constante de ay'' + by' + cy = d et en déduire la solution générale
\item À la fin de cette leçon, je sais déterminer l'unique solution d'un problème de Cauchy (conditions initiales y(t₀) et y'(t₀))
\item À la fin de cette leçon, je sais modéliser une situation concrète (décroissance, circuit électrique, oscillation) par une équation différentielle et interpréter la solution
\item À la fin de cette leçon, je sais rédiger et justifier une démarche complète de résolution en citant les théorèmes utilisés
\end{itemize}
\end{objectifs}

\begin{prerequis}
\begin{itemize}
\item Dérivation des fonctions usuelles (polynômes, exponentielle, cosinus, sinus) et dérivation des composées
\item Fonction exponentielle : propriétés algébriques et dérivée de e\^{}{kx}
\item Résolution des équations du second degré dans ℝ et dans ℂ (discriminant, racines complexes conjuguées)
\item Notion de limite et comportement asymptotique de l'exponentielle
\item Lecture et interprétation graphique d'une courbe représentative
\end{itemize}
\end{prerequis}

\newpage

\coursec{Présentation et vocabulaire}

\textbf{Situation d'accroche.} À Yaoundé, un technicien en électricité met en service un circuit RL alimentant les machines d'un atelier de menuiserie : à l'instant où il ferme l'interrupteur, l'intensité du courant ne passe pas brutalement de $0$ à sa valeur maximale ; elle monte progressivement, et sa vitesse de montée dépend à chaque instant de l'intensité elle-même. Au CHU de Douala, un médecin suit la concentration d'un médicament dans le sang d'un patient : celle-ci décroît à une vitesse proportionnelle à la quantité restante. À la SONATREL, un ingénieur étudie les oscillations d'un circuit RLC : la tension aux bornes du condensateur oscille en s'amortissant. Dans ces trois situations, la grandeur inconnue (intensité, concentration, tension) vérifie une relation entre \emph{elle-même} et \emph{ses dérivées}. Une telle relation s'appelle une \cle{équation différentielle}. Cette leçon donne les outils pour les résoudre et interpréter leurs solutions.

\textbf{Problématique.} Comment modéliser ces phénomènes par des équations dont l'inconnue est une fonction, et comment reconnaître, vérifier et choisir les solutions de telles équations ?

\courssub{Des situations concrètes aux équations différentielles}

Dans les classes précédentes, résoudre une équation comme $2x+3=7$ ou $x^2-5x+6=0$ consistait à chercher un \emph{nombre} inconnu. Dans beaucoup de problèmes techniques, l'inconnue n'est pas un nombre mais une \emph{fonction du temps} : on cherche par exemple l'intensité $i(t)$, la charge $q(t)$ ou la position $x(t)$. Les lois de la physique relient alors cette fonction à sa vitesse de variation (la dérivée) ou à son accélération (la dérivée seconde). Voyons quatre exemples fondateurs.

\textbf{Exemple 1 : la décroissance radioactive.} Un échantillon contient $N(t)$ noyaux radioactifs à l'instant $t$. L'expérience montre que le nombre de désintégrations par unité de temps est proportionnel au nombre de noyaux présents :
\[ N'(t) = -\lambda\, N(t), \]
où $\lambda > 0$ est la constante radioactive. L'inconnue est la fonction $N$, et l'équation relie $N$ à sa dérivée $N'$.

\textbf{Exemple 2 : le courant dans un circuit RL (bobine en série).} Une bobine d'inductance $L$ est branchée en série avec une résistance $R$ sur une source de tension continue $E$. À la fermeture de l'interrupteur, l'intensité $i(t)$ vérifie, d'après la loi des mailles :
\[ L\,i'(t) + R\,i(t) = E. \]
L'inconnue $i$ apparaît avec sa dérivée $i'$ : c'est une équation d'ordre 1.

\begin{popfigure}
\begin{center}
\begin{tikzpicture}[line width=0.9pt, popdark]
% Interrupteur
\draw (0,2) -- (1,2);
\draw[fill=white] (1,2) circle (0.06);
\draw[fill=white] (1.8,2) circle (0.06);
\draw (1,2) -- (1.75,2.45);
\node[above] at (1.4,2.35) {$K$};
% Résistance (rectangle)
\draw (1.8,2) -- (2.6,2);
\draw (2.6,1.85) rectangle (4.4,2.15);
\node at (3.5,2.55) {$R$};
\draw (4.4,2) -- (5.5,2);
% Bobine (Inductance) - symbole standard : demi-cercles
\draw (5.5,2) -- (6.0,2);
\foreach \x in {6.2, 6.6, 7.0, 7.4} {
  \draw (\x, 2.3) arc (90:-90:0.2 and 0.3);
}
\node at (6.8, 2.7) {$L$};
\draw (7.8,2) -- (8.5,2);
% Générateur
\draw (0,2) -- (0,0.7);
\draw[line width=1.4pt] (-0.35,0.7) -- (0.35,0.7);
\draw[line width=0.9pt] (-0.18,0.4) -- (0.18,0.4);
\draw (0,0.4) -- (0,0);
\node[left] at (-0.4,0.55) {$E$};
% Retour
\draw (0,0) -- (8.5,0);
% Flèche courant i
\draw[->, poporange, line width=1pt] (4.0, 2.5) -- (4.0, 2.8);
\node[above, poporange] at (4.0, 2.95) {$i(t)$};
\end{tikzpicture}
\end{center}
\legende{Circuit RL série : à la fermeture de l'interrupteur $K$, l'intensité $i(t)$ vérifie $L\,i'(t)+R\,i(t)=E$.}
\end{popfigure}

\textbf{Exemple 3 : la chute avec frottements.} Un objet de masse $m$ tombe verticalement ; l'air exerce une force de frottement proportionnelle à la vitesse. Le principe fondamental de la dynamique donne :
\[ m\,v'(t) = mg - k\,v(t), \]
où $v(t)$ est la vitesse, $g$ l'accélération de la pesanteur et $k$ le coefficient de frottement.

\textbf{Exemple 4 : le circuit RLC.} Dans un circuit comportant une résistance $R$, une bobine d'inductance $L$ et un condensateur $C$, la charge $q(t)$ du condensateur vérifie :
\[ L\,q''(t) + R\,q'(t) + \frac{1}{C}\,q(t) = 0. \]
Cette fois, l'équation fait intervenir la dérivée \emph{seconde} $q''$ : c'est une équation d'ordre 2, qui décrit des oscillations amorties.

\begin{savaistu}[Newton, Leibniz et la naissance des équations différentielles]
C'est précisément pour résoudre des équations différentielles de mécanique qu'Isaac \cle{Newton} (1643--1727) et Gottfried Wilhelm \cle{Leibniz} (1646--1716) ont créé, indépendamment l'un de l'autre, le calcul différentiel au \textsc{xvii}\textsuperscript{e} siècle. Newton voulait décrire le mouvement des planètes : sa fameuse loi $F = ma$ est, en fait, une équation différentielle, car l'accélération $a$ est la dérivée seconde de la position. Aujourd'hui, les équations différentielles sont partout : prévisions météorologiques, conception des ponts et des antennes, réseaux électriques de la SONATREL, modèles épidémiologiques utilisés par le ministère de la Santé publique.
\end{savaistu}

\courssub{Définition d'une équation différentielle}

\begin{definition}[Équation différentielle]
Une \cle{équation différentielle} est une équation dont l'\emph{inconnue est une fonction} $y$ d'une variable réelle (souvent notée $x$ ou $t$), et qui relie cette fonction, la variable et une ou plusieurs de ses dérivées $y'$, $y''$, \dots

\emph{Exemples :} $y' = 2y$ ; \quad $y' + y = 5$ ; \quad $2y'' + 3y' - y = 0$.
\end{definition}

L'écriture $y' = 2y$ est un raccourci : elle signifie que l'on cherche toutes les fonctions $f$, définies et dérivables sur un intervalle, telles que pour tout $x$ de cet intervalle, $f'(x) = 2f(x)$.

\begin{definition}[Ordre d'une équation différentielle]
L'\cle{ordre} d'une équation différentielle est l'ordre de la dérivée la plus élevée qui y apparaît.
\begin{itemize}
\item Une équation ne faisant intervenir que $y$ et $y'$ est d'\cle{ordre 1} : par exemple $y' = 2y$ ou $L\,i' + R\,i = E$.
\item Une équation faisant intervenir $y''$ est d'\cle{ordre 2} : par exemple $Lq'' + Rq' + \dfrac{1}{C}q = 0$.
\end{itemize}
\end{definition}

\begin{exemplebox}[Reconnaître l'ordre]
\begin{itemize}
\item $y' - 3y = 0$ est une équation différentielle d'ordre 1.
\item $y'' + 4y' + 4y = 0$ est une équation différentielle d'ordre 2.
\item $m\,v' = mg - k\,v$ (chute avec frottements) est d'ordre 1 en l'inconnue $v$.
\end{itemize}
\end{exemplebox}

\courssub{Solution d'une équation différentielle sur un intervalle}

\begin{definition}[Solution sur un intervalle]
Soit $I$ un intervalle de $\mathbb{R}$. Une \cle{solution} d'une équation différentielle sur $I$ est une fonction $f$, dérivable au moins autant de fois que l'ordre de l'équation sur $I$, qui \emph{vérifie l'équation pour tout $x$ de $I$} lorsqu'on remplace $y$ par $f(x)$, $y'$ par $f'(x)$, etc.
\end{definition}

Pour savoir si une fonction donnée est solution, on procède par \cle{substitution} : on calcule ses dérivées, on les injecte dans l'équation, et on vérifie que l'égalité est vraie pour tout $x$.

\begin{methode}[Vérifier qu'une fonction est solution d'une équation différentielle]
Pour vérifier qu'une fonction $f$ est solution d'une équation différentielle sur un intervalle $I$ (la variable pouvant être $x$, $t$, ou une autre lettre) :
\begin{enumerate}
\item Calculer la (ou les) dérivée(s) de $f$ qui interviennent dans l'équation ($f'$, éventuellement $f''$).
\item \emph{Substituer} : remplacer $y$ par $f(x)$, $y'$ par $f'(x)$, $y''$ par $f''(x)$ dans l'équation.
\item Simplifier et vérifier que l'égalité obtenue est vraie \emph{pour tout $x$ de $I$}.
\item Conclure par une phrase : « $f$ est solution de l'équation sur $I$ ».
\end{enumerate}
\end{methode}

\begin{exoresolu}[Vérifier une solution avec condition initiale]
\textbf{Énoncé.} On considère l'équation différentielle $(E) : y' = 2y$ sur $\mathbb{R}$. Montrer que la fonction $f$ définie par $f(t) = 3e^{2t}$ est solution de $(E)$ sur $\mathbb{R}$, et vérifier que $f(0) = 3$.
\tcblower
\textbf{Solution détaillée.}
\begin{enumerate}
\item \emph{Dérivée.} $f(t) = 3e^{2t}$, donc $f'(t) = 3 \times 2e^{2t} = 6e^{2t}$.
\item \emph{Substitution.} D'une part $f'(t) = 6e^{2t}$ ; d'autre part $2f(t) = 2 \times 3e^{2t} = 6e^{2t}$.
\item \emph{Comparaison.} Pour tout réel $t$, $f'(t) = 2f(t)$ : l'égalité est vraie.
\item \emph{Conclusion.} $f$ est solution de $(E)$ sur $\mathbb{R}$. De plus $f(0) = 3e^{0} = 3 \times 1 = 3$ : la condition $f(0)=3$ est bien vérifiée.
\end{enumerate}
\end{exoresolu}

\begin{attention}[Solution d'une équation ou d'une équation différentielle ?]
Erreur fréquente : confondre « solution d'une équation » et « solution d'une équation différentielle ».
\begin{itemize}
\item La solution de $2x+3=7$ est un \emph{nombre} : $x = 2$.
\item Une solution de $y' = 2y$ est une \emph{fonction} : par exemple $f(t) = 3e^{2t}$.
\end{itemize}
On ne répond donc jamais « $y = 3$ » à une équation différentielle ! Autre piège : vérifier l'égalité seulement en \emph{un} point ne suffit pas — l'égalité doit être vraie \emph{pour tout} $x$ de l'intervalle. Enfin, une fonction peut vérifier l'équation sans vérifier la condition initiale : $g(t) = 5e^{2t}$ est solution de $y' = 2y$, mais $g(0) = 5 \neq 3$.
\end{attention}

\courssub{Une infinité de solutions ; le rôle des conditions initiales}

Reprenons l'équation $y' = y$. La fonction $x \mapsto e^{x}$ est solution, car sa dérivée est elle-même. Mais $x \mapsto 2e^{x}$ est aussi solution (sa dérivée est $2e^{x}$, égale à elle-même), et plus généralement, pour \emph{tout} réel $k$, la fonction $f_k(x) = k\,e^{x}$ vérifie :
\[ f_k'(x) = k\,e^{x} = f_k(x). \]
L'équation $y' = y$ admet donc une \emph{infinité} de solutions, une pour chaque valeur de $k$. C'est naturel physiquement : l'équation $N' = -\lambda N$ décrit la loi de décroissance, mais ne dit pas \emph{combien} de noyaux on avait au départ. Pour décrire une situation précise, il faut une information supplémentaire : la valeur de la fonction à un instant donné.

\begin{definition}[Condition initiale, problème de Cauchy]
Une \cle{condition initiale} est la donnée de la valeur de la fonction (et, pour une équation d'ordre 2, de sa dérivée) en un point $x_0$ : par exemple $y(x_0) = y_0$.

Un \cle{problème de Cauchy} est la donnée d'une équation différentielle accompagnée de condition(s) initiale(s) :
\begin{itemize}
\item ordre 1 : $\begin{cases} y' = ay \\ y(x_0) = y_0 \end{cases}$ \quad (une condition) ;
\item ordre 2 : $\begin{cases} ay'' + by' + cy = 0 \\ y(x_0) = y_0 \\ y'(x_0) = y_1 \end{cases}$ \quad (deux conditions).
\end{itemize}
On admet qu'un tel problème de Cauchy admet, sur l'intervalle considéré, une \emph{unique} solution : les conditions initiales \emph{sélectionnent} une seule courbe parmi la famille infinie des solutions.
\end{definition}

\begin{popfigure}
\begin{center}
\begin{tikzpicture}
\begin{axis}[
  width=\linewidth, height=6.2cm,
  axis lines=middle,
  xlabel={$x$}, ylabel={$y$},
  xmin=-2.2, xmax=2.2, ymin=-1, ymax=8.5,
  xtick={-2,-1,1,2}, ytick={1,2,4,6,8},
  tick label style={font=\small},
  legend style={font=\small, at={(0.03,0.97)}, anchor=north west, draw=popline}
]
\addplot[popcurve, domain=-2.2:2.1, samples=120] {0.5*exp(x)};
\addplot[popcurve, poporange, domain=-2.2:2.1, samples=120] {exp(x)};
\addplot[popcurve, popgreen!80!black, domain=-2.2:1.75, samples=120] {2*exp(x)};
\addplot[popcurve, poppurple, domain=-2.2:1.45, samples=120] {4*exp(x)};
\addplot[popcurve, popmuted, dashed, domain=-2.2:2.1, samples=120] {-exp(x)};
\legend{$k=0.5$, $k=1$, $k=2$, $k=4$, $k=-1$}
\addplot[only marks, mark=*, mark size=2pt, poporange] coordinates {(0,1)};
\node[poporange, font=\small, anchor=west] at (axis cs:0.08,0.55) {$(0\,;\,1)$};
\end{axis}
\end{tikzpicture}
\end{center}
\legende{Famille de solutions de $y' = y$ : les courbes $y = k\,e^{x}$. La condition initiale $y(0)=1$ sélectionne une unique courbe ($k=1$, en orange).}
\end{popfigure}

Sur la figure, toutes les courbes $y = k\,e^{x}$ sont solutions de $y' = y$ : elles ne se croisent pas et « remplissent » le plan. Imposer $y(0) = 1$ revient à exiger que la courbe passe par le point $(0\,;\,1)$ : une seule courbe convient, celle de $k = 1$.

\begin{exoresolu}[Choisir la solution grâce à une condition initiale]
\textbf{Énoncé.} On admet que toutes les solutions de $y' = y$ sur $\mathbb{R}$ sont les fonctions $f_k(x) = k\,e^{x}$, $k \in \mathbb{R}$. Déterminer la solution du problème de Cauchy : $y' = y$ et $y(0) = 5$.
\tcblower
\textbf{Solution détaillée.}
\begin{enumerate}
\item Toute solution s'écrit $f_k(x) = k\,e^{x}$.
\item La condition initiale impose $f_k(0) = 5$, c'est-à-dire $k\,e^{0} = k = 5$.
\item Donc $k = 5$ et l'unique solution du problème de Cauchy est $f(x) = 5e^{x}$.
\end{enumerate}
\emph{Vérification :} $f'(x) = 5e^{x} = f(x)$ et $f(0) = 5$. $\square$
\end{exoresolu}

\begin{exercice}[À vous de jouer]
\begin{enumerate}
\item Montrer que la fonction $g(t) = 7e^{-3t}$ est solution sur $\mathbb{R}$ de l'équation $y' = -3y$. Vérifie-t-elle la condition $y(0) = 7$ ?
\item La fonction $h(x) = x^2 + 1$ est-elle solution de l'équation $y' = 2x$ sur $\mathbb{R}$ ? Justifier.
\item Parmi les fonctions $f_k(x) = k\,e^{2x}$ ($k \in \mathbb{R}$), toutes solutions de $y' = 2y$, déterminer celle qui vérifie $y(0) = -4$.
\end{enumerate}
\end{exercice}

\begin{corrige}[Exercice « À vous de jouer »]
\begin{enumerate}
\item $g'(t) = 7 \times (-3)e^{-3t} = -21e^{-3t}$ et $-3g(t) = -21e^{-3t}$. Pour tout réel $t$, $g'(t) = -3g(t)$ : $g$ est solution sur $\mathbb{R}$. De plus $g(0) = 7e^{0} = 7$ : la condition est vérifiée.
\item $h'(x) = 2x$ : l'égalité $y' = 2x$ est vérifiée pour tout $x$, donc $h$ est solution sur $\mathbb{R}$ (remarque : $x \mapsto x^2 + 5$ est aussi solution, ce qui confirme l'infinité de solutions).
\item $f_k(0) = k = -4$, donc la solution cherchée est $f(x) = -4e^{2x}$. Vérification : $f'(x) = -8e^{2x} = 2f(x)$ et $f(0) = -4$.
\end{enumerate}
\end{corrige}

\begin{aretenir}[L'essentiel de la section]
\begin{itemize}
\item Une \cle{équation différentielle} est une équation dont l'inconnue est une \emph{fonction}, reliant $y$, $y'$, $y''$, \dots et la variable.
\item Son \cle{ordre} est celui de la dérivée la plus élevée : $y' = 2y$ est d'ordre 1 ; $ay'' + by' + cy = 0$ est d'ordre 2.
\item Une \cle{solution} sur un intervalle $I$ est une fonction qui vérifie l'équation pour tout $x$ de $I$ ; on la teste par \emph{calcul des dérivées puis substitution}.
\item Une équation différentielle admet en général une \emph{infinité} de solutions ; un \cle{problème de Cauchy} (équation + conditions initiales) en sélectionne une unique.
\end{itemize}
\end{aretenir}

\coursec{Équations du type $y' = ay$}

Dans la section précédente, nous avons découvert ce qu'est une équation différentielle : une égalité qui relie une fonction inconnue $y$ à ses dérivées, et dont l'inconnue n'est pas un nombre mais une \cle{fonction}. Nous avons aussi appris le vocabulaire : ordre, solution, solution générale, condition initiale. Nous allons maintenant résoudre complètement la famille d'équations différentielles la plus simple et la plus utile de tout le programme : les équations du type $y' = ay$, où $a$ est un réel fixé. Ces équations modélisent toutes les situations où la \cle{vitesse de variation} d'une grandeur est \cle{proportionnelle} à la grandeur elle-même : élimination d'un médicament dans le sang, désintégration radioactive, croissance d'une population, intérêts composés en continu.

\courssub{Position du problème : une vitesse proportionnelle à la quantité}

\begin{definition}[Équation différentielle $y' = ay$]
Soit $a$ un réel fixé. On appelle \cle{équation différentielle linéaire homogène du premier ordre à coefficient constant} l'équation
\[ (E) : \quad y' = ay \]
c'est-à-dire l'ensemble des fonctions $f$, définies et dérivables sur $\mathbb{R}$, telles que pour tout réel $t$ :
\[ f'(t) = a\,f(t). \]
\end{definition}

Lisons cette équation avec des mots de la vie courante. La dérivée $f'(t)$ est la \emph{vitesse de variation} de la grandeur $f(t)$ à l'instant $t$. Dire que $f'(t) = a\,f(t)$ signifie donc : \og à chaque instant, la grandeur varie à une vitesse proportionnelle à sa propre valeur, avec le coefficient de proportionnalité $a$ \fg{}.

\begin{itemize}
\item Plus il y a de médicament dans le sang, plus l'organisme en élimine par heure.
\item Plus un capital est grand, plus il produit d'intérêts.
\item Plus une population de bactéries est nombreuse, plus elle se reproduit vite.
\end{itemize}

C'est exactement cette idée que traduit l'équation $y' = ay$.

\courssub{Recherche de solutions : l'idée de l'exponentielle}

Cherchons des solutions \emph{simples}. Nous connaissons une fonction remarquable qui est presque égale à sa propre dérivée : la fonction exponentielle. En effet, pour tout réel $t$ :
\[ \left(e^{at}\right)' = a\,e^{at}. \]
Autrement dit, la fonction $t \mapsto e^{at}$ vérifie exactement l'équation $y' = ay$ : c'est une \cle{solution} de $(E)$.

Mieux : si l'on multiplie cette fonction par une constante $k$, la propriété est conservée, car
\[ \left(k\,e^{at}\right)' = k \times a\,e^{at} = a\left(k\,e^{at}\right). \]
Donc \textbf{toutes} les fonctions $t \mapsto k\,e^{at}$, pour $k$ réel quelconque, sont des solutions de $(E)$. La question naturelle est alors : y en a-t-il d'autres ? Le théorème suivant répond que non.

\begin{propriete}[Théorème fondamental — démonstration exigible]
Soit $a$ un réel fixé. Les solutions de l'équation différentielle $y' = ay$ sont exactement les fonctions définies sur $\mathbb{R}$ par
\[ f(t) = k\,e^{at}, \quad \text{où } k \text{ est un réel quelconque.} \]
Cette démonstration est au programme et peut être demandée au Probatoire et au Baccalauréat technique.
\end{propriete}

\begin{experience}[Démonstration guidée du théorème]
L'astuce consiste à \og neutraliser \fg{} l'exponentielle en étudiant une fonction auxiliaire. Suivons les étapes.

\textbf{Étape 1.} Soit $f$ une solution quelconque de $(E)$ : pour tout réel $t$, on a $f'(t) = a\,f(t)$. On pose
\[ g(t) = f(t)\,e^{-at}. \]
La fonction $g$ est dérivable sur $\mathbb{R}$ comme produit de fonctions dérivables.

\textbf{Étape 2.} Dérivons $g$ avec la formule du produit $(uv)' = u'v + uv'$ :
\begin{align*}
g'(t) &= f'(t)\,e^{-at} + f(t)\times\left(-a\,e^{-at}\right) \\
&= e^{-at}\left(f'(t) - a\,f(t)\right).
\end{align*}

\textbf{Étape 3.} Or $f$ est solution de $(E)$, donc $f'(t) - a\,f(t) = 0$ pour tout $t$. Par conséquent
\[ g'(t) = 0 \quad \text{pour tout réel } t. \]

\textbf{Étape 4.} Une fonction de dérivée nulle sur $\mathbb{R}$ est \cle{constante} : il existe un réel $k$ tel que $g(t) = k$ pour tout $t$, c'est-à-dire
\[ f(t)\,e^{-at} = k \quad \Longleftrightarrow \quad f(t) = k\,e^{at}. \]

\textbf{Étape 5 (réciproque).} Nous avons déjà vérifié que toute fonction de la forme $t \mapsto k\,e^{at}$ est bien solution de $(E)$.

\textbf{Conclusion.} Les solutions de $y' = ay$ sont exactement les fonctions $f(t) = k\,e^{at}$, $k \in \mathbb{R}$. \hfill$\blacksquare$
\end{experience}

\begin{attention}[La constante est multiplicative, pas additive]
L'erreur la plus fréquente est d'écrire la solution générale sous la forme $f(t) = e^{at} + k$. C'est \textbf{faux} : la constante $k$ \emph{multiplie} l'exponentielle, elle ne s'additionne pas. Vérifions : si $f(t) = e^{at} + k$, alors $f'(t) = a\,e^{at}$, qui n'est pas égal à $a\,f(t) = a\,e^{at} + ak$ (sauf si $k = 0$). Retenez : la solution générale de $y' = ay$ est $\boxed{f(t) = k\,e^{at}}$.
\end{attention}

\courssub{Condition initiale : existence et unicité de la solution}

L'équation $y' = ay$ possède une infinité de solutions (une pour chaque valeur de $k$). Mais dès que l'on connaît la valeur de la grandeur à un instant donné, une seule solution convient.

\begin{propriete}[Existence et unicité]
Soit $a$ un réel, $t_0$ un réel et $y_0$ un réel. Il existe une \cle{unique} solution $f$ de l'équation $y' = ay$ vérifiant la condition initiale $f(t_0) = y_0$. Elle est donnée par
\[ f(t) = y_0\,e^{a(t - t_0)}. \]
\end{propriete}

En effet, toute solution s'écrit $f(t) = k\,e^{at}$. La condition $f(t_0) = y_0$ donne $k\,e^{at_0} = y_0$, d'où $k = y_0\,e^{-at_0}$. En reportant :
\[ f(t) = y_0\,e^{-at_0}\,e^{at} = y_0\,e^{a(t-t_0)}. \]
Dans le cas très fréquent où la condition est donnée à l'instant $t_0 = 0$, la formule se simplifie joliment : $f(0) = y_0$ donne directement $k = y_0$, et
\[ f(t) = y_0\,e^{at}. \]

\begin{methode}[Résoudre $y' = ay$ avec une condition initiale]
\begin{enumerate}
\item Identifier le coefficient $a$ dans l'équation (attention à son signe).
\item Écrire la solution générale : $f(t) = k\,e^{at}$.
\item Remplacer $t$ par $t_0$ et $f(t_0)$ par $y_0$ pour obtenir une équation d'inconnue $k$.
\item Résoudre cette équation : $k = y_0\,e^{-at_0}$.
\item Conclure en écrivant la solution complète, puis répondre à la question posée (valeur à un instant donné, seuil, etc.).
\end{enumerate}
\end{methode}

\begin{exoresolu}[Concentration d'un médicament dans le sang]
Un technicien de laboratoire injecte un médicament à un patient. La concentration $C(t)$ du produit dans le sang (en mg/L), $t$ heures après l'injection, vérifie l'équation différentielle
\[ C' = -0{,}2\,C, \qquad C(0) = 10 \text{ mg/L}. \]
\begin{enumerate}
\item Déterminer l'expression de $C(t)$.
\item Calculer la concentration au bout de $5$ heures (arrondir à $0{,}01$ mg/L).
\end{enumerate}
\tcblower
\textbf{1.} L'équation est du type $y' = ay$ avec $a = -0{,}2$. La solution générale est
\[ C(t) = k\,e^{-0{,}2\,t}. \]
La condition initiale $C(0) = 10$ donne $k\,e^{0} = 10$, donc $k = 10$. D'où
\[ \boxed{C(t) = 10\,e^{-0{,}2\,t}}. \]

\textbf{2.} Au bout de $5$ heures :
\[ C(5) = 10\,e^{-0{,}2 \times 5} = 10\,e^{-1} \approx 10 \times 0{,}3679 \approx 3{,}68 \text{ mg/L}. \]
La concentration est tombée à environ $3{,}68$ mg/L au bout de $5$ heures : le médicament s'élimine de plus en plus lentement, mais ne s'annule jamais complètement (l'exponentielle reste strictement positive).
\end{exoresolu}

\courssub{Interprétation du signe de $a$}

Le signe du coefficient $a$ commande tout le comportement des solutions.

\begin{aretenir}[Trois régimes selon le signe de $a$]
Pour une solution $f(t) = k\,e^{at}$ avec $k > 0$ :
\begin{itemize}
\item si $a > 0$ : la fonction est strictement \cle{croissante} et tend vers $+\infty$ quand $t \to +\infty$ (\emph{croissance exponentielle}) ;
\item si $a < 0$ : la fonction est strictement \cle{décroissante} et tend vers $0$ quand $t \to +\infty$ (\emph{décroissance exponentielle}) ; la droite $y = 0$ est asymptote horizontale ;
\item si $a = 0$ : l'équation devient $y' = 0$, les solutions sont les fonctions \cle{constantes} $f(t) = k$.
\end{itemize}
C'est cohérent avec le cours sur la dérivation : $f'(t) = a\,f(t)$, donc $f'$ a le signe de $a \times f(t)$.
\end{aretenir}

\begin{popfigure}
\begin{center}
\begin{tikzpicture}
\begin{axis}[width=0.85\linewidth,height=6.5cm,
axis lines=middle, xlabel={$t$}, ylabel={$y$},
xmin=-0.5, xmax=8.5, ymin=-0.5, ymax=11,
xtick={1,2,3,4,5,6,7,8}, ytick={2,4,6,8,10},
legend style={at={(0.03,0.97)},anchor=north west,font=\small}]
\addplot[popcurve,popblue,domain=0:8,samples=100]{10*exp(-0.5*x)};
\addlegendentry{$k=10$}
\addplot[popcurve,poporange,domain=0:8,samples=100]{6*exp(-0.5*x)};
\addlegendentry{$k=6$}
\addplot[popcurve,popgreen,domain=0:8,samples=100]{3*exp(-0.5*x)};
\addlegendentry{$k=3$}
\addplot[popcurve,popmuted,domain=0:8,samples=2]{0};
\end{axis}
\end{tikzpicture}
\end{center}
\legende{Solutions de $y' = -0{,}5\,y$ pour plusieurs valeurs de $k$ : décroissances exponentielles vers l'asymptote horizontale $y = 0$.}
\end{popfigure}

\begin{popfigure}
\begin{center}
\begin{tikzpicture}
\begin{axis}[width=0.85\linewidth,height=6.5cm,
axis lines=middle, xlabel={$t$}, ylabel={$y$},
xmin=-0.5, xmax=6.5, ymin=-0.5, ymax=13,
xtick={1,2,3,4,5,6}, ytick={2,4,6,8,10,12},
legend style={at={(0.03,0.97)},anchor=north west,font=\small}]
\addplot[popcurve,popblue,domain=0:6,samples=100]{2*exp(0.3*x)};
\addlegendentry{$f(t)=2e^{0{,}3t}$}
\addplot[popcurve,poporange,domain=0:6,samples=2]{0.6*x+2};
\addlegendentry{tangente en $t=0$}
\addplot[mark=*,popdark] coordinates {(0,2)};
\end{axis}
\end{tikzpicture}
\end{center}
\legende{Solution de $y' = 0{,}3\,y$ avec $y(0) = 2$ : croissance exponentielle. La tangente à l'origine a pour pente $f'(0) = 0{,}3 \times 2 = 0{,}6$.}
\end{popfigure}

\courssub{Temps de demi-vie et temps de doublement}

Dans les applications, deux questions reviennent constamment : \og en combien de temps la quantité est-elle réduite de moitié ? \fg{} (décroissance) ou \og en combien de temps double-t-elle ? \fg{} (croissance).

\begin{definition}[Demi-vie et temps de doublement]
Soit $f(t) = y_0\,e^{at}$.
\begin{itemize}
\item Si $a < 0$, le \cle{temps de demi-vie} $T$ est la durée telle que $f(T) = \dfrac{y_0}{2}$, soit $e^{aT} = \dfrac{1}{2}$, d'où $T = \dfrac{-\ln 2}{a} = \dfrac{\ln 2}{|a|}$.
\item Si $a > 0$, le \cle{temps de doublement} $T$ est la durée telle que $f(T) = 2y_0$, soit $e^{aT} = 2$, d'où $T = \dfrac{\ln 2}{a}$.
\end{itemize}
Ces durées ne dépendent \textbf{pas} de la valeur initiale $y_0$ : c'est une propriété remarquable des phénomènes exponentiels.
\end{definition}

\begin{exemplebox}[Demi-vie du médicament]
Pour le médicament de l'exercice résolu ($a = -0{,}2$), le temps de demi-vie est
\[ T = \frac{\ln 2}{0{,}2} = 5\ln 2 \approx 3{,}47 \text{ h} \approx 3 \text{ h } 28 \text{ min}. \]
Toutes les $3$ h $28$ min environ, la concentration est divisée par deux, quelle que soit la concentration de départ.
\end{exemplebox}

\begin{savaistu}[La datation au carbone 14]
Le carbone 14 est un isotope radioactif présent dans tous les êtres vivants, en proportion constante tant qu'ils sont en vie. À la mort de l'organisme, le carbone 14 n'est plus renouvelé et sa quantité $N(t)$ vérifie l'équation $N' = aN$ avec $a < 0$ : sa demi-vie est d'environ \num{5730} ans. En mesurant la proportion de carbone 14 restant dans un ossement, un tissu ou un objet en bois, les archéologues remontent à l'équation $N(t) = N_0\,e^{at}$ et calculent l'âge de l'échantillon. Cette méthode, mise au point par Willard Libby (prix Nobel de chimie en 1960), a permis de dater des vestiges du monde entier, y compris des sites archéologiques africains.
\end{savaistu}

\begin{exercice}[Intérêts composés en continu]
Un capital de \num{500000} FCFA est placé à intérêts composés continus au taux annuel de $6\,\%$. On admet que le capital $K(t)$, en francs, vérifie $K'(t) = 0{,}06\,K(t)$ avec $t$ en années.
\begin{enumerate}
\item Déterminer $K(t)$.
\item Calculer le capital au bout de $10$ ans (arrondir au franc).
\item Déterminer le temps de doublement du capital.
\end{enumerate}
\end{exercice}

\begin{corrige}[Exercice : intérêts composés en continu]
\textbf{1.} Solution générale : $K(t) = k\,e^{0{,}06t}$. La condition $K(0) = \num{500000}$ donne $k = \num{500000}$, donc $K(t) = \num{500000}\,e^{0{,}06t}$.

\textbf{2.} $K(10) = \num{500000}\,e^{0{,}6} \approx \num{500000} \times 1{,}8221 \approx \num{911059}$ FCFA.

\textbf{3.} Temps de doublement : $T = \dfrac{\ln 2}{0{,}06} \approx 11{,}55$ ans, soit environ $11$ ans et $7$ mois.
\end{corrige}

\coursec{Équations ay'' + by' + cy = 0 : équation caractéristique}

Dans la section précédente, nous avons résolu les équations différentielles du premier ordre $y' = ay$ : leurs solutions sont les fonctions exponentielles $t \mapsto Ke^{at}$. Nous allons maintenant monter d'un cran et étudier des équations où interviennent à la fois une fonction, sa dérivée première \emph{et} sa dérivée seconde. Ces équations modélisent des phénomènes fondamentaux de la vie technique : les oscillations d'un circuit électrique, les vibrations d'une suspension de véhicule, le mouvement d'un piston amorti. La bonne nouvelle, c'est que toute la résolution repose sur une simple équation du second degré, que vous savez résoudre depuis longtemps : l'\cle{équation caractéristique}.

\courssub{Présentation et idée directrice}

\begin{definition}[Équation différentielle linéaire homogène du second ordre]
Soient $a$, $b$ et $c$ trois réels, avec $a \neq 0$. On appelle \cle{équation différentielle linéaire homogène du second ordre à coefficients constants} toute équation de la forme
\[ (E) : \quad ay'' + by' + cy = 0 \]
où $y$ est une fonction inconnue de la variable réelle $t$, deux fois dérivable sur $\mathbb{R}$.
\end{definition}

Décortiquons ce vocabulaire :
\begin{itemize}
\item \emph{Second ordre} : la dérivée la plus « haute » qui apparaît est la dérivée seconde $y''$.
\item \emph{Linéaire} : $y$, $y'$ et $y''$ apparaissent à la puissance $1$, sans produit entre eux.
\item \emph{Homogène} : le second membre est nul (égal à $0$).
\item \emph{Coefficients constants} : $a$, $b$ et $c$ sont des nombres fixés, qui ne dépendent pas de $t$.
\end{itemize}

\textbf{Idée directrice.} Pour l'équation $y' = ay$, les solutions étaient des exponentielles. L'exponentielle possède une propriété remarquable : dériver $e^{rt}$ revient simplement à la multiplier par $r$. Cherchons donc, par analogie, des solutions de $(E)$ de la forme
\[ y(t) = e^{rt} \]
où $r$ est un nombre (réel ou complexe) à déterminer. On a alors :
\[ y'(t) = re^{rt} \qquad \text{et} \qquad y''(t) = r^2 e^{rt}. \]
Substituons dans l'équation $(E)$ :
\[ a r^2 e^{rt} + b r e^{rt} + c e^{rt} = 0. \]
En factorisant par $e^{rt}$, qui n'est jamais nul :
\[ e^{rt}\left(ar^2 + br + c\right) = 0 \quad \Longleftrightarrow \quad ar^2 + br + c = 0. \]
Ainsi, $e^{rt}$ est solution de $(E)$ exactement lorsque $r$ est solution d'une équation du second degré. C'est ce résultat qui motive la définition suivante.

\courssub{L'équation caractéristique}

\begin{definition}[Équation caractéristique]
À l'équation différentielle $ay'' + by' + cy = 0$ (avec $a \neq 0$), on associe l'\cle{équation caractéristique} :
\[ ar^2 + br + c = 0 \]
d'inconnue $r$. C'est une équation du second degré à coefficients réels.
\end{definition}

\begin{attention}[Bien écrire l'équation caractéristique]
Deux erreurs reviennent très souvent dans les copies :
\begin{itemize}
\item Écrire l'équation caractéristique avec des $x$ au lieu de $r$ : on écrit $ar^2 + br + c = 0$, et l'inconnue s'appelle $r$ (c'est une convention liée au taux de l'exponentielle $e^{rt}$). Surtout, ne confondez pas l'inconnue $r$ de l'équation caractéristique avec la variable $t$ de la fonction $y$.
\item Oublier la condition $a \neq 0$ : si $a = 0$, l'équation différentielle n'est plus du second ordre et l'équation caractéristique n'est plus du second degré ; toute la méthode s'effondre. Vérifiez toujours que le coefficient de $y''$ est non nul avant de commencer.
\end{itemize}
\end{attention}

\begin{propriete}[Lien entre solutions exponentielles et équation caractéristique]
La fonction $t \mapsto e^{rt}$ est solution de l'équation différentielle $ay'' + by' + cy = 0$ \textbf{si et seulement si} $r$ est racine de l'équation caractéristique $ar^2 + br + c = 0$.
\end{propriete}

\begin{experience}[Démonstration guidée de la propriété]
Justifions rigoureusement cette propriété par substitution, en suivant les étapes.
\begin{enumerate}
\item Soit $y(t) = e^{rt}$. Calculez $y'(t)$ et $y''(t)$ : on obtient $y'(t) = re^{rt}$ et $y''(t) = r^2e^{rt}$.
\item Remplacez dans le premier membre de l'équation :
\[ ay'' + by' + cy = ar^2e^{rt} + bre^{rt} + ce^{rt} = e^{rt}\left(ar^2 + br + c\right). \]
\item Une exponentielle ne s'annule jamais : pour tout réel $t$, $e^{rt} \neq 0$. Donc le produit $e^{rt}(ar^2+br+c)$ est nul pour tout $t$ si et seulement si $ar^2 + br + c = 0$.
\item Conclusion : $e^{rt}$ est solution de l'équation différentielle si et seulement si $r$ est racine de l'équation caractéristique. $\square$
\end{enumerate}
Cette démonstration est formatrice : elle montre que la propriété n'est pas une formule tombée du ciel, mais une conséquence directe de la règle de dérivation de l'exponentielle.
\end{experience}

\courssub{Résolution de l'équation caractéristique dans C}

L'équation caractéristique $ar^2 + br + c = 0$ est une équation du second degré à coefficients réels. On calcule son \cle{discriminant} :
\[ \Delta = b^2 - 4ac. \]
Trois cas se présentent, selon le signe de $\Delta$.

\textbf{Premier cas : $\Delta > 0$.} L'équation caractéristique admet deux racines réelles distinctes :
\[ r_1 = \frac{-b - \sqrt{\Delta}}{2a} \qquad \text{et} \qquad r_2 = \frac{-b + \sqrt{\Delta}}{2a}. \]
On obtient alors deux solutions exponentielles réelles $e^{r_1 t}$ et $e^{r_2 t}$.

\textbf{Deuxième cas : $\Delta = 0$.} L'équation caractéristique admet une racine réelle double :
\[ r_0 = -\frac{b}{2a}. \]
On obtient la solution exponentielle $e^{r_0 t}$, mais il faudra une deuxième solution indépendante (ce sera $t\,e^{r_0 t}$, comme nous le verrons dans la section suivante).

\textbf{Troisième cas : $\Delta < 0$.} L'équation caractéristique admet deux racines complexes conjuguées :
\[ r_1 = \alpha - i\beta \qquad \text{et} \qquad r_2 = \alpha + i\beta, \quad \text{avec } \alpha = -\frac{b}{2a} \text{ et } \beta = \frac{\sqrt{-\Delta}}{2a} > 0. \]
Ce cas correspond physiquement aux phénomènes \emph{oscillants} (vibrations, courants alternatifs amortis).

\begin{aretenir}[Les trois cas selon le signe du discriminant]
\begin{center}
\begin{tabularx}{\linewidth}{|l|X|X|}
\hline
\rowcolor{popblueL} \textbf{Signe de $\Delta$} & \textbf{Racines de l'équation caractéristique} & \textbf{Nature des solutions} \\
\hline
$\Delta > 0$ & Deux racines réelles distinctes $r_1$ et $r_2$ & Exponentielles réelles $e^{r_1 t}$, $e^{r_2 t}$ \\
\hline
$\Delta = 0$ & Une racine réelle double $r_0 = -\dfrac{b}{2a}$ & Exponentielle $e^{r_0 t}$ (et $t e^{r_0 t}$) \\
\hline
$\Delta < 0$ & Deux racines complexes conjuguées $\alpha \pm i\beta$ & Oscillations (sinus et cosinus amortis) \\
\hline
\end{tabularx}
\end{center}
\end{aretenir}

\begin{popfigure}
\begin{center}
\begin{tikzpicture}[
  node distance=0.9cm and 1.1cm,
  start/.style={rectangle, rounded corners, draw=popblue, fill=popblueL, thick, text width=4.6cm, align=center, font=\small},
  decision/.style={diamond, draw=poporange, fill=poporangeL, thick, aspect=2.2, align=center, font=\small, inner sep=1pt},
  cas/.style={rectangle, rounded corners, draw=popgreen, fill=popgreenL, thick, text width=3.4cm, align=center, font=\small},
  sol/.style={rectangle, rounded corners, draw=poppurple, fill=poppurpleL, thick, text width=3.4cm, align=center, font=\small},
  fleche/.style={-{Stealth[length=2.5mm]}, thick, popdark}
]
\node[start, align=center] (eq){Équation $ay''+by'+cy=0$\\ Écrire l'équation caractéristique $ar^2+br+c=0$};
\node[decision, below=of eq] (delta) {Signe de $\Delta = b^2-4ac$ ?};
\node[cas, below left=1.4cm and 0.2cm of delta, align=center] (pos){$\Delta > 0$\\ deux racines réelles $r_1 \neq r_2$};
\node[cas, below=1.4cm of delta, align=center] (nul){$\Delta = 0$\\ racine double $r_0$};
\node[cas, below right=1.4cm and 0.2cm of delta, align=center] (neg){$\Delta < 0$\\ racines $\alpha \pm i\beta$};
\node[sol, below=of pos] (s1) {$y = \lambda e^{r_1 t} + \mu e^{r_2 t}$};
\node[sol, below=of nul] (s2) {$y = (\lambda t + \mu)e^{r_0 t}$};
\node[sol, below=of neg] (s3) {$y = e^{\alpha t}\big(\lambda\cos\beta t + \mu\sin\beta t\big)$};
\draw[fleche] (eq) -- (delta);
\draw[fleche] (delta) -| (pos);
\draw[fleche] (delta) -- (nul);
\draw[fleche] (delta) -| (neg);
\draw[fleche] (pos) -- (s1);
\draw[fleche] (nul) -- (s2);
\draw[fleche] (neg) -- (s3);
\end{tikzpicture}
\end{center}
\legende{Arbre de décision : on calcule $\Delta$, puis le signe de $\Delta$ détermine la forme de la solution générale ($\lambda$ et $\mu$ sont des constantes réelles).}
\end{popfigure}

\courssub{Annonce du théorème de résolution}

Nous pouvons maintenant annoncer le résultat central, qui sera démontré et exploité dans la section suivante.

\begin{propriete}[Solution générale de $ay''+by'+cy=0$ — admise pour l'instant]
Soit $\Delta = b^2 - 4ac$ le discriminant de l'équation caractéristique $ar^2 + br + c = 0$. La solution générale de l'équation différentielle $ay'' + by' + cy = 0$ est :
\begin{itemize}
\item si $\Delta > 0$ : \quad $y(t) = \lambda e^{r_1 t} + \mu e^{r_2 t}$, où $r_1$ et $r_2$ sont les deux racines réelles ;
\item si $\Delta = 0$ : \quad $y(t) = (\lambda t + \mu)e^{r_0 t}$, où $r_0$ est la racine double ;
\item si $\Delta < 0$ : \quad $y(t) = e^{\alpha t}\big(\lambda \cos(\beta t) + \mu \sin(\beta t)\big)$, où $\alpha \pm i\beta$ sont les racines complexes conjuguées ;
\end{itemize}
dans tous les cas, $\lambda$ et $\mu$ désignent des constantes réelles arbitraires.
\end{propriete}

Remarquez la logique d'ensemble : résoudre l'équation différentielle revient à résoudre une équation du second degré, puis à « traduire » ses racines en fonctions. Tout le travail algébrique se concentre donc sur l'équation caractéristique.

\begin{methode}[Méthode complète : de l'équation différentielle aux solutions]
\begin{enumerate}
\item Vérifier que l'équation est bien de la forme $ay'' + by' + cy = 0$ avec $a \neq 0$, et identifier les coefficients $a$, $b$, $c$ (attention aux signes !).
\item Écrire l'équation caractéristique $ar^2 + br + c = 0$.
\item Calculer le discriminant $\Delta = b^2 - 4ac$ et déterminer son signe.
\item Résoudre l'équation caractéristique selon le cas : deux racines réelles, racine double, ou racines complexes conjuguées.
\item Conclure en écrivant la solution générale correspondant au signe de $\Delta$.
\end{enumerate}
\end{methode}

\begin{exemplebox}[Un exemple chiffré de bout en bout]
Considérons l'équation différentielle
\[ y'' - 5y' + 6y = 0. \]
\begin{enumerate}
\item Coefficients : $a = 1$, $b = -5$, $c = 6$ (et $a = 1 \neq 0$).
\item Équation caractéristique : $r^2 - 5r + 6 = 0$.
\item Discriminant : $\Delta = (-5)^2 - 4 \times 1 \times 6 = 25 - 24 = 1 > 0$.
\item Deux racines réelles distinctes :
\[ r_1 = \frac{5 - 1}{2} = 2 \qquad \text{et} \qquad r_2 = \frac{5 + 1}{2} = 3. \]
\item Solution générale : $y(t) = \lambda e^{2t} + \mu e^{3t}$, avec $\lambda, \mu \in \mathbb{R}$.
\end{enumerate}
\end{exemplebox}

\begin{exoresolu}[Vérifier qu'une fonction est bien solution]
Énoncé. On considère l'équation différentielle $(E) : y'' - 5y' + 6y = 0$.
\begin{enumerate}
\item Montrer que la fonction $f(t) = e^{2t}$ est solution de $(E)$.
\item Montrer que $g(t) = 3e^{2t} - 2e^{3t}$ est aussi solution de $(E)$.
\end{enumerate}
\tcblower
Solution détaillée.
\begin{enumerate}
\item On calcule les dérivées de $f(t) = e^{2t}$ : $f'(t) = 2e^{2t}$ et $f''(t) = 4e^{2t}$. Substituons :
\begin{align*}
f''(t) - 5f'(t) + 6f(t) &= 4e^{2t} - 5 \times 2e^{2t} + 6e^{2t} \\
&= (4 - 10 + 6)e^{2t} = 0 \times e^{2t} = 0.
\end{align*}
Donc $f$ est bien solution de $(E)$. C'est cohérent avec la propriété du cours : $r = 2$ est racine de l'équation caractéristique $r^2 - 5r + 6 = 0$ car $4 - 10 + 6 = 0$.
\item On calcule : $g'(t) = 6e^{2t} - 6e^{3t}$ et $g''(t) = 12e^{2t} - 18e^{3t}$. Alors :
\begin{align*}
g'' - 5g' + 6g &= \left(12e^{2t} - 18e^{3t}\right) - 5\left(6e^{2t} - 6e^{3t}\right) + 6\left(3e^{2t} - 2e^{3t}\right) \\
&= (12 - 30 + 18)e^{2t} + (-18 + 30 - 12)e^{3t} \\
&= 0 \times e^{2t} + 0 \times e^{3t} = 0.
\end{align*}
Donc $g$ est solution de $(E)$. On observe au passage qu'une combinaison de solutions est encore une solution : c'est la \emph{linéarité} de l'équation.
\end{enumerate}
\end{exoresolu}

\courssub{Pourquoi ces équations sont partout en technique}

Les équations du type $ay'' + by' + cy = 0$ décrivent tous les systèmes physiques où une « inertie » (terme en $y''$), un « frottement » (terme en $y'$) et un « rappel » (terme en $y$) se combinent. L'exemple emblématique est le circuit électrique RLC série.

\begin{popfigure}
\begin{center}
\begin{tikzpicture}[thick, popdark, scale=1.05]
% Cadre du circuit
\draw (0,0) -- (0,2);
\draw (0,2) -- (1,2);
% Résistance R (zigzag)
\draw[decorate, decoration={zigzag, segment length=6pt, amplitude=3pt}] (1,2) -- (3,2);
\draw (3,2) -- (4,2);
\node[above] at (2,2.25) {$R$};
% Bobine L (arcs)
\draw (4,2) -- (4.4,2);
\draw (4.4,2) arc[start angle=180, end angle=0, radius=0.3];
\draw (5.0,2) arc[start angle=180, end angle=0, radius=0.3];
\draw (5.6,2) arc[start angle=180, end angle=0, radius=0.3];
\draw (6.2,2) -- (7,2);
\node[above] at (5.3,2.35) {$L$};
% Condensateur C
\draw (7,2) -- (7,1.4);
\draw (6.6,1.4) -- (7.4,1.4);
\draw (6.6,1.15) -- (7.4,1.15);
\draw (7,1.15) -- (7,0);
\node[right] at (7.5,1.3) {$C$};
% Bas du circuit
\draw (7,0) -- (0,0);
% Générateur
\draw[fill=popcream] (0,1) circle (0.35);
\node at (0,1) {$e$};
% Annotations
\node[below, popmuted, font=\small] at (3.5,-0.25) {Circuit RLC série : $R$ résistance, $L$ bobine, $C$ condensateur};
\end{tikzpicture}
\end{center}
\legende{Circuit RLC série. En notant $q(t)$ la charge du condensateur, la loi des mailles donne $Lq'' + Rq' + \dfrac{q}{C} = 0$ : une équation $ay''+by'+cy=0$ avec $a = L$, $b = R$, $c = \dfrac{1}{C}$.}
\end{popfigure}

Appliquons la loi des mailles à ce circuit parcouru par un courant : la somme des tensions aux bornes de la bobine ($Lq''$), de la résistance ($Rq'$) et du condensateur ($\dfrac{q}{C}$) est nulle en l'absence de générateur. On obtient :
\[ Lq'' + Rq' + \frac{1}{C}\,q = 0. \]
L'équation caractéristique associée est $Lr^2 + Rr + \dfrac{1}{C} = 0$, de discriminant $\Delta = R^2 - \dfrac{4L}{C}$. Selon la valeur de la résistance $R$ :
\begin{itemize}
\item si $R$ est grande ($\Delta > 0$), le régime est \emph{apériodique} : le courant s'éteint sans osciller ;
\item si $R$ a la valeur critique $R = 2\sqrt{L/C}$ ($\Delta = 0$), le régime est \emph{critique} : retour à l'équilibre le plus rapide possible ;
\item si $R$ est petite ($\Delta < 0$), le régime est \emph{pseudo-périodique} : le courant oscille en s'amortissant.
\end{itemize}
Le signe de $\Delta$ a donc une interprétation physique directe : c'est lui qui décide si le système oscille ou non.

\begin{savaistu}[L'électricien et la résonance]
L'ingénieur électricien utilise précisément ces équations pour éviter les \cle{résonances} dangereuses dans les réseaux. Lorsque les racines de l'équation caractéristique sont complexes ($\Delta < 0$), le circuit oscille à une fréquence propre liée à $\beta$. Si une source extérieure excite le circuit à cette même fréquence, les amplitudes peuvent croître dangereusement : surtensions, échauffements, destruction de matériel. C'est le même phénomène qui a fait s'effondrer le pont de Tacoma aux États-Unis en 1940 sous l'effet du vent. Au Cameroun, le dimensionnement des installations électriques des usines et des postes de transformation tient compte de ces calculs pour garantir la stabilité du réseau.
\end{savaistu}

\begin{exercice}[À vous de jouer]
Pour chacune des équations différentielles suivantes, écrire l'équation caractéristique, calculer son discriminant et préciser dans quel cas ($\Delta > 0$, $\Delta = 0$ ou $\Delta < 0$) on se trouve :
\begin{enumerate}
\item $y'' + 3y' + 2y = 0$ ;
\item $y'' - 4y' + 4y = 0$ ;
\item $2y'' + y' + y = 0$.
\end{enumerate}
\end{exercice}

\begin{corrige}[Exercice : équations caractéristiques]
\begin{enumerate}
\item Équation caractéristique : $r^2 + 3r + 2 = 0$. Discriminant : $\Delta = 9 - 8 = 1 > 0$. Cas de deux racines réelles distinctes ($r_1 = -2$, $r_2 = -1$).
\item Équation caractéristique : $r^2 - 4r + 4 = 0$. Discriminant : $\Delta = 16 - 16 = 0$. Cas de la racine double $r_0 = 2$.
\item Équation caractéristique : $2r^2 + r + 1 = 0$. Discriminant : $\Delta = 1 - 8 = -7 < 0$. Cas de deux racines complexes conjuguées : $r = \dfrac{-1 \pm i\sqrt{7}}{4}$.
\end{enumerate}
\end{corrige}

\begin{aretenir}[L'essentiel de la section]
\begin{itemize}
\item À toute équation $ay'' + by' + cy = 0$ (avec $a \neq 0$), on associe l'\cle{équation caractéristique} $ar^2 + br + c = 0$.
\item $e^{rt}$ est solution si et seulement si $r$ est racine de l'équation caractéristique (justification par substitution).
\item Le discriminant $\Delta = b^2 - 4ac$ commande tout : $\Delta > 0$ donne deux racines réelles, $\Delta = 0$ une racine double, $\Delta < 0$ deux racines complexes conjuguées $\alpha \pm i\beta$.
\item La forme de la solution générale dépend du cas rencontré ; elle sera établie en détail dans la section suivante.
\end{itemize}
\end{aretenir}

\coursec{Solutions générales de $ay'' + by' + cy = 0$}

Dans la section précédente, nous avons associé à toute équation différentielle $ay'' + by' + cy = 0$ son \cle{équation caractéristique} $ar^2 + br + c = 0$, et nous avons observé que les fonctions exponentielles $t \mapsto e^{rt}$ sont solutions lorsque $r$ est racine de cette équation. Reste une question essentielle : comment décrire \emph{toutes} les solutions ? La réponse dépend du signe du discriminant $\Delta = b^2 - 4ac$ de l'équation caractéristique. C'est l'objet de cette section, qui constitue le cœur de la leçon.

\courssub{Le théorème fondamental : trois cas selon le signe de $\Delta$}

L'idée intuitive est la suivante. Une équation du second ordre « contient » deux degrés de liberté : pour démarrer un mouvement (penser à une masse suspendue à un ressort), il faut connaître la position initiale \emph{et} la vitesse initiale. Il est donc naturel que la solution générale dépende de \textbf{deux constantes arbitraires}, que nous noterons $\lambda$ et $\mu$.

\begin{propriete}[Théorème fondamental — forme des solutions générales]
Soit l'équation différentielle $(E) : ay'' + by' + cy = 0$ où $a$, $b$, $c$ sont des réels avec $a \neq 0$, et soit $\Delta = b^2 - 4ac$ le discriminant de son équation caractéristique $ar^2 + br + c = 0$.
\begin{itemize}
\item \textbf{Cas $\Delta > 0$ :} l'équation caractéristique admet deux racines réelles distinctes $r_1$ et $r_2$. La solution générale de $(E)$ est
\[ f(t) = \lambda e^{r_1 t} + \mu e^{r_2 t}, \qquad (\lambda, \mu) \in \mathbb{R}^2. \]
\item \textbf{Cas $\Delta = 0$ :} l'équation caractéristique admet une racine double $r_0 = -\dfrac{b}{2a}$. La solution générale de $(E)$ est
\[ f(t) = (\lambda t + \mu)\, e^{r_0 t}, \qquad (\lambda, \mu) \in \mathbb{R}^2. \]
\item \textbf{Cas $\Delta < 0$ :} l'équation caractéristique admet deux racines complexes conjuguées $\alpha + i\beta$ et $\alpha - i\beta$ (avec $\beta \neq 0$). La solution générale \emph{réelle} de $(E)$ est
\[ f(t) = e^{\alpha t}\big(\lambda \cos(\beta t) + \mu \sin(\beta t)\big), \qquad (\lambda, \mu) \in \mathbb{R}^2. \]
\end{itemize}
Le cas $\Delta > 0$ est démontré ci-dessous (démonstration guidée). Le cas $\Delta = 0$ est justifié de façon guidée. Le cas $\Delta < 0$ est admis pour l'existence ; la vérification par substitution sur un exemple est exigible au Probatoire.
\end{propriete}

\begin{experience}[Démonstration guidée : le cas $\Delta > 0$]
Supposons $\Delta > 0$ et notons $r_1$, $r_2$ les deux racines réelles distinctes de $ar^2 + br + c = 0$.
\begin{enumerate}
\item \textbf{Deux solutions exponentielles.} Posons $f_1(t) = e^{r_1 t}$. Alors $f_1'(t) = r_1 e^{r_1 t}$ et $f_1''(t) = r_1^2 e^{r_1 t}$, donc
\[ a f_1'' + b f_1' + c f_1 = (a r_1^2 + b r_1 + c)\, e^{r_1 t} = 0 \]
car $r_1$ est racine de l'équation caractéristique. Donc $f_1$ est solution. De même, $f_2(t) = e^{r_2 t}$ est solution.
\item \textbf{Toute combinaison linéaire est solution.} Soit $f = \lambda f_1 + \mu f_2$. Par linéarité de la dérivation :
\[ a f'' + b f' + c f = \lambda\,(a f_1'' + b f_1' + c f_1) + \mu\,(a f_2'' + b f_2' + c f_2) = 0. \]
Donc toute fonction de la forme $\lambda e^{r_1 t} + \mu e^{r_2 t}$ est solution.
\item \textbf{Il n'y en a pas d'autres.} On admet (résultat d'unicité évoqué plus bas) que toute solution s'écrit sous cette forme : les deux constantes $\lambda$ et $\mu$ suffisent à décrire l'ensemble complet des solutions.
\end{enumerate}
\end{experience}

\begin{experience}[Justification guidée : le cas $\Delta = 0$]
Supposons $\Delta = 0$ : la racine double est $r_0 = -\dfrac{b}{2a}$, et l'on a la relation $2a r_0 + b = 0$. La fonction $f_1(t) = e^{r_0 t}$ est solution (même calcul que ci-dessus). Mais il nous faut une \emph{deuxième} solution, indépendante de la première. Vérifions que $f_2(t) = t\, e^{r_0 t}$ convient.
\begin{align*}
f_2'(t) &= e^{r_0 t} + r_0 t\, e^{r_0 t} = (1 + r_0 t)\, e^{r_0 t}, \\
f_2''(t) &= r_0 e^{r_0 t} + r_0(1 + r_0 t)\, e^{r_0 t} = (2 r_0 + r_0^2 t)\, e^{r_0 t}.
\end{align*}
Alors :
\begin{align*}
a f_2'' + b f_2' + c f_2 &= \big[a(2 r_0 + r_0^2 t) + b(1 + r_0 t) + c t\big] e^{r_0 t} \\
&= \big[\underbrace{(2 a r_0 + b)}_{=\,0} + \underbrace{(a r_0^2 + b r_0 + c)}_{=\,0}\, t\big] e^{r_0 t} = 0.
\end{align*}
Donc $t\,e^{r_0 t}$ est bien solution, et toute combinaison $(\lambda t + \mu)e^{r_0 t}$ l'est aussi.
\end{experience}

\begin{attention}[Erreur fréquente : le facteur $t$ oublié]
Dans le cas $\Delta = 0$, la solution générale est $(\lambda t + \mu)e^{r_0 t}$ et \textbf{non} pas $\lambda e^{r_0 t} + \mu e^{r_0 t}$ (ce qui ne dépendrait en réalité que d'une seule constante $\lambda + \mu$ !). Le facteur $t$ devant $\lambda$ est indispensable : c'est lui qui fournit la deuxième solution indépendante. Au Probatoire comme au Baccalauréat technique, oublier ce $t$ fausse toute la suite de l'exercice.
\end{attention}

\begin{exemplebox}[Vérification par substitution dans le cas $\Delta < 0$]
Considérons l'équation $y'' + 2y' + 5y = 0$. L'équation caractéristique $r^2 + 2r + 5 = 0$ a pour discriminant $\Delta = 4 - 20 = -16 < 0$, d'où les racines $r = \dfrac{-2 \pm 4i}{2} = -1 \pm 2i$. Ici $\alpha = -1$ et $\beta = 2$. Vérifions que $f(t) = e^{-t}\cos(2t)$ est solution.
\begin{align*}
f'(t) &= e^{-t}\big(-\cos(2t) - 2\sin(2t)\big), \\
f''(t) &= e^{-t}\big(-3\cos(2t) + 4\sin(2t)\big).
\end{align*}
Alors :
\[ f'' + 2f' + 5f = e^{-t}\big[(-3 - 2 + 5)\cos(2t) + (4 - 4)\sin(2t)\big] = 0. \]
La fonction $f$ est bien solution. On vérifierait de même que $e^{-t}\sin(2t)$ est solution.
\end{exemplebox}

\courssub{Structure de l'ensemble des solutions}

Dans les trois cas, la solution générale s'écrit $f = \lambda f_1 + \mu f_2$ où $f_1$ et $f_2$ sont deux solutions \cle{indépendantes} (aucune n'est un multiple constant de l'autre) :

\begin{center}
\begin{tabularx}{\linewidth}{|l|X|X|}
\hline
\rowcolor{popblueL} \textbf{Cas} & \textbf{Solution $f_1$} & \textbf{Solution $f_2$} \\
\hline
$\Delta > 0$ & $e^{r_1 t}$ & $e^{r_2 t}$ \\
\hline
$\Delta = 0$ & $e^{r_0 t}$ & $t\,e^{r_0 t}$ \\
\hline
$\Delta < 0$ & $e^{\alpha t}\cos(\beta t)$ & $e^{\alpha t}\sin(\beta t)$ \\
\hline
\end{tabularx}
\end{center}

\begin{aretenir}[Espace des solutions]
L'ensemble des solutions de $ay'' + by' + cy = 0$ est l'ensemble de toutes les combinaisons $\lambda f_1 + \mu f_2$ de deux solutions indépendantes. On dit que c'est un \cle{espace vectoriel de dimension 2} : comme un plan est décrit par deux vecteurs directeurs, l'ensemble des solutions est décrit par deux solutions de base, et il faut exactement \textbf{deux constantes} pour désigner une solution particulière.
\end{aretenir}

\courssub{Conditions initiales et problème de Cauchy}

Les constantes $\lambda$ et $\mu$ se déterminent lorsqu'on connaît deux informations sur la solution, typiquement la valeur de $f$ et de $f'$ en un instant $t_0$ :
\[ f(t_0) = y_0 \qquad \text{et} \qquad f'(t_0) = y_1. \]
Ces deux conditions conduisent à un \textbf{système de deux équations à deux inconnues} $\lambda$ et $\mu$, que l'on résout par substitution ou combinaison.

\begin{propriete}[Existence et unicité — admise]
Pour tous réels $t_0$, $y_0$, $y_1$, l'équation $ay'' + by' + cy = 0$ admet une \textbf{unique} solution $f$ vérifiant $f(t_0) = y_0$ et $f'(t_0) = y_1$. Ce problème s'appelle un \cle{problème de Cauchy} du second ordre. Concrètement : la position et la vitesse initiales déterminent entièrement le mouvement.
\end{propriete}

\begin{methode}[Déterminer $\lambda$ et $\mu$ à partir des conditions initiales]
\begin{enumerate}
\item Résoudre l'équation caractéristique et écrire la solution générale $f(t)$ selon le signe de $\Delta$.
\item Calculer la dérivée $f'(t)$ (attention à la dérivation d'un produit dans les cas $\Delta = 0$ et $\Delta < 0$).
\item Écrire les deux équations $f(t_0) = y_0$ et $f'(t_0) = y_1$ : on obtient un système en $\lambda$ et $\mu$.
\item Résoudre le système, remplacer $\lambda$ et $\mu$ dans $f(t)$, puis \textbf{vérifier} que les conditions initiales sont satisfaites.
\end{enumerate}
\end{methode}

\begin{exoresolu}[Résolution complète avec conditions initiales]
\textbf{Énoncé.} Résoudre l'équation différentielle $y'' - 3y' + 2y = 0$ avec les conditions initiales $y(0) = 1$ et $y'(0) = 0$.
\tcblower
\textbf{Solution détaillée.}
\textbf{Étape 1.} Équation caractéristique : $r^2 - 3r + 2 = 0$. Discriminant : $\Delta = 9 - 8 = 1 > 0$. Racines : $r_1 = \dfrac{3 - 1}{2} = 1$ et $r_2 = \dfrac{3 + 1}{2} = 2$. La solution générale est
\[ f(t) = \lambda e^{t} + \mu e^{2t}. \]
\textbf{Étape 2.} Dérivée : $f'(t) = \lambda e^{t} + 2\mu e^{2t}$.
\textbf{Étape 3.} Conditions initiales en $t_0 = 0$ (avec $e^0 = 1$) :
\[ \begin{cases} f(0) = \lambda + \mu = 1 \\ f'(0) = \lambda + 2\mu = 0 \end{cases} \]
\textbf{Étape 4.} En soustrayant la première équation de la deuxième : $\mu = -1$, puis $\lambda = 1 - \mu = 2$. La solution cherchée est donc
\[ \boxed{f(t) = 2e^{t} - e^{2t}}. \]
\textbf{Vérification :} $f(0) = 2 - 1 = 1$ et $f'(t) = 2e^t - 2e^{2t}$ donne $f'(0) = 2 - 2 = 0$. Les conditions sont bien satisfaites.
\end{exoresolu}

\coursec{Équations $ay'' + by' + cy = d$ : second membre constant}

Nous abordons maintenant le cas non homogène où le second membre est une constante $d \in \mathbb{R}$. La méthode repose sur le principe de \cle{superposition} : la solution générale est la somme de la solution générale de l'équation homogène associée et d'une \cle{solution particulière} de l'équation non homogène.

\begin{propriete}[Structure de la solution générale]
Soit $(E_d) : ay'' + by' + cy = d$ avec $a \neq 0$. Soit $f_h$ la solution générale de l'équation homogène associée $(E_0) : ay'' + by' + cy = 0$ et $f_p$ une solution particulière de $(E_d)$. Alors la solution générale de $(E_d)$ est
\[ f(t) = f_h(t) + f_p(t). \]
\end{propriete}

\courssub{Recherche d'une solution particulière constante}

On cherche une solution constante $f_p(t) = k$. Alors $f_p' = 0$ et $f_p'' = 0$. L'équation devient $c k = d$. Trois cas se présentent :

\begin{itemize}
\item \textbf{Cas $c \neq 0$ :} On prend $k = \dfrac{d}{c}$. La solution particulière est $f_p(t) = \dfrac{d}{c}$.
\item \textbf{Cas $c = 0$ et $b \neq 0$ :} L'équation est $ay'' + by' = d$. On cherche une solution de la forme $f_p(t) = k t$. Alors $f_p' = k$, $f_p'' = 0$, ce qui donne $b k = d$, soit $k = \dfrac{d}{b}$. Solution particulière : $f_p(t) = \dfrac{d}{b}\,t$.
\item \textbf{Cas $c = 0$ et $b = 0$ :} L'équation est $a y'' = d$. On cherche $f_p(t) = k t^2$. Alors $f_p'' = 2k$, ce qui donne $2a k = d$, soit $k = \dfrac{d}{2a}$. Solution particulière : $f_p(t) = \dfrac{d}{2a}\,t^2$.
\end{itemize}

\begin{attention}[Cas $c=0$ : ne pas oublier le facteur $t$ ou $t^2$]
Si $c=0$, une solution \emph{constante} ne convient pas (elle donnerait $0=d$). Il faut multiplier par $t$ (si $b \neq 0$) ou $t^2$ (si $b=0$) pour annuler la dérivée d'ordre le plus bas non nul. C'est la méthode de l'\cle{annulation de l'opérateur différentiel}.
\end{attention}

\begin{methode}[Résolution de $ay'' + by' + cy = d$]
\begin{enumerate}
\item Résoudre l'équation homogène $ay'' + by' + cy = 0$ (solution $f_h$ avec deux constantes $\lambda, \mu$).
\item Trouver une solution particulière $f_p$ selon les cas ci-dessus ($c \neq 0$, $c=0/b\neq0$, $c=b=0$).
\item Écrire la solution générale $f = f_h + f_p$.
\item Si des conditions initiales sont données, déterminer $\lambda$ et $\mu$ en utilisant $f$ et $f'$ (attention : $f_p$ contribue aux conditions initiales !).
\end{enumerate}
\end{methode}

\begin{exoresolu}[Circuit $RLC$ sous tension constante]
\textbf{Énoncé.} Dans un circuit série $RLC$ alimenté par une tension continue $E$, la charge $q(t)$ du condensateur vérifie $L q'' + R q' + \frac{1}{C} q = E$. Données : $L = 1\,\text{H}$, $R = 4\,\Omega$, $C = 0,25\,\text{F}$, $E = 10\,\text{V}$. Conditions initiales : $q(0) = 0$, $q'(0) = 0$ (circuit au repos).
\tcblower
\textbf{Solution.}
L'équation est $q'' + 4q' + 4q = 10$ (on a divisé par $L=1$, et $1/C = 4$).
\textbf{1. Homogène :} $r^2 + 4r + 4 = 0 \Rightarrow (r+2)^2=0$, racine double $r_0 = -2$ ($\Delta=0$).
$q_h(t) = (\lambda t + \mu)e^{-2t}$.
\textbf{2. Particulière :} $c=4 \neq 0$, donc $q_p = \frac{d}{c} = \frac{10}{4} = 2,5$.
\textbf{3. Générale :} $q(t) = (\lambda t + \mu)e^{-2t} + 2,5$.
\textbf{4. Conditions initiales :}
$q(0) = \mu + 2,5 = 0 \Rightarrow \mu = -2,5$.
$q'(t) = \lambda e^{-2t} - 2(\lambda t + \mu)e^{-2t}$.
$q'(0) = \lambda - 2\mu = 0 \Rightarrow \lambda = 2\mu = -5$.
\textbf{Résultat :} $\boxed{q(t) = (-5t - 2,5)e^{-2t} + 2,5}$.
\textbf{Interprétation :} La charge tend vers $2,5\,\text{C}$ (régime permanent) sans oscillation (régime critique $\Delta=0$).
\end{exoresolu}

\coursec{Applications concrètes et synthèse}

\courssub{Modélisation mécanique : oscillateur amorti}

Un système masse-ressort-amortisseur de masse $m$, raideur $k$, coefficient d'amortissement $c$, soumis à une force extérieure $F(t)$, vérifie l'équation différentielle :
\[ m y'' + c y' + k y = F(t) \]
où $y(t)$ est le déplacement par rapport à l'équilibre.
\begin{itemize}
\item \textbf{Libre ($F=0$) :} $m y'' + c y' + k y = 0$. C'est l'équation homogène étudiée précédemment.
\item \textbf{Forcé constant ($F = F_0$) :} $m y'' + c y' + k y = F_0$. On retrouve le cas $d = F_0$. La solution particulière $y_p = F_0/k$ est l'allongement statique sous la charge $F_0$.
\end{itemize}

\courssub{Interprétation physique : les trois régimes (rappel)}

Pour l'équation libre $m y'' + c y' + k y = 0$, le discriminant est $\Delta = c^2 - 4mk$. Le signe de $\Delta$ détermine l'allure du mouvement :

\begin{itemize}
\item \textbf{$\Delta > 0$ : régime amorti (apériodique).} $c^2 > 4mk$ (amortissement fort). La solution est une somme de deux exponentielles décroissantes : le système revient vers sa position d'équilibre \emph{sans osciller}, mais lentement.
\item \textbf{$\Delta = 0$ : régime critique.} $c^2 = 4mk$. Le retour à l'équilibre est le \emph{plus rapide possible sans oscillation}. C'est le réglage idéal recherché en mécanique (amortisseurs automobiles).
\item \textbf{$\Delta < 0$ : régime pseudo-périodique.} $c^2 < 4mk$ (amortissement faible). La solution est $e^{-\frac{c}{2m}t}(\lambda \cos \omega t + \mu \sin \omega t)$ avec $\omega = \frac{\sqrt{4mk-c^2}}{2m}$ : le système oscille en s'affaiblissant.
\end{itemize}

\begin{popfigure}
\begin{center}
\begin{tikzpicture}
\begin{axis}[width=0.95\linewidth, height=6.2cm, axis lines=middle, xlabel={$t$}, ylabel={$y$}, xmin=0, xmax=6.4, ymin=-0.6, ymax=1.3, xtick=\empty, ytick=\empty, legend style={at={(0.98,0.95)}, anchor=north east, font=\small}]
\addplot[popcurve, domain=0:6.2, samples=150] {exp(-1.5*x)};
\addlegendentry{$\Delta>0$ : amorti}
\addplot[popcurve, poporange, domain=0:6.2, samples=150] {(1+2*x)*exp(-2*x)};
\addlegendentry{$\Delta=0$ : critique}
\addplot[popcurve, popgreen, domain=0:6.2, samples=200] {exp(-0.5*x)*cos(deg(3*x))};
\addlegendentry{$\Delta<0$ : oscillant}
\end{axis}
\end{tikzpicture}
\end{center}
\legende{Trois solutions partant de la même position initiale $y(0)=1$ : retour lent sans oscillation ($\Delta>0$), retour le plus rapide ($\Delta=0$), oscillations amorties ($\Delta<0$).}
\end{popfigure}

Dans le régime pseudo-périodique, la courbe $f(t) = e^{-t}\cos(3t)$ reste « emprisonnée » entre les deux courbes enveloppes $t \mapsto e^{-t}$ et $t \mapsto -e^{-t}$, car $-1 \leq \cos(3t) \leq 1$.

\begin{popfigure}
\begin{center}
\begin{tikzpicture}
\begin{axis}[width=0.95\linewidth, height=6.2cm, axis lines=middle, xlabel={$t$}, ylabel={$y$}, xmin=0, xmax=6.4, ymin=-1.2, ymax=1.3, xtick=\empty, ytick=\empty]
\addplot[popcurve, popblue, domain=0:6.2, samples=250] {exp(-x)*cos(deg(3*x))};
\addplot[popmuted, dashed, thick, domain=0:6.2, samples=100] {exp(-x)};
\addplot[popmuted, dashed, thick, domain=0:6.2, samples=100] {-exp(-x)};
\end{axis}
\end{tikzpicture}
\end{center}
\legende{Oscillation amortie $f(t)=e^{-t}\cos(3t)$ (bleu) et ses enveloppes $t\mapsto e^{-t}$ et $t\mapsto -e^{-t}$ (pointillés).}
\end{popfigure}

\begin{savaistu}[Les amortisseurs et le régime critique]
Les amortisseurs des voitures sont réglés au plus près du \textbf{régime critique} ($\Delta = 0$). Si l'amortissement est trop faible ($\Delta < 0$), la voiture rebondit plusieurs fois après chaque nid-de-poule — inconfortable et dangereux sur les routes dégradées. S'il est trop fort ($\Delta > 0$), la suspension réagit trop lentement et la roue perd le contact avec la chaussée. Le régime critique garantit le retour à l'équilibre le plus rapide sans oscillation : c'est le compromis idéal entre confort et tenue de route, sur les avenues de Yaoundé comme sur les pistes de l'Extrême-Nord.
\end{savaistu}

\courssub{Synthèse : fiche méthode « Équations différentielles du 2nd ordre »}

\begin{center}
\begin{tabularx}{\linewidth}{|l|X|}
\hline
\rowcolor{popblueL} \textbf{Type d'équation} & \textbf{Démarche de résolution} \\
\hline
$y' = a y$ & Solution générale : $y = C e^{a t}$. Condition initiale $\Rightarrow$ valeur de $C$. \\
\hline
$a y'' + b y' + c y = 0$ & 1. Équation caractéristique $a r^2 + b r + c = 0$, $\Delta = b^2 - 4ac$. \\
& 2. Selon $\Delta$ : \\
& \quad $\Delta > 0$ : $r_1, r_2$ réels distincts $\to y = \lambda e^{r_1 t} + \mu e^{r_2 t}$ \\
& \quad $\Delta = 0$ : $r_0$ double $\to y = (\lambda t + \mu) e^{r_0 t}$ \\
& \quad $\Delta < 0$ : $\alpha \pm i\beta$ $\to y = e^{\alpha t}(\lambda \cos \beta t + \mu \sin \beta t)$ \\
& 3. Conditions initiales $\Rightarrow$ système $2\times 2$ pour $\lambda, \mu$. \\
\hline
$a y'' + b y' + c y = d$ & 1. Résoudre l'homogène $\to y_h$. \\
($d$ constant) & 2. Solution particulière $y_p$ : \\
& \quad $c \neq 0 \to y_p = d/c$ \\
& \quad $c=0, b\neq 0 \to y_p = (d/b)t$ \\
& \quad $c=b=0 \to y_p = (d/2a)t^2$ \\
& 3. Générale $y = y_h + y_p$. \\
& 4. Conditions initiales sur $y$ et $y'$ $\Rightarrow$ $\lambda, \mu$. \\
\hline
\end{tabularx}
\end{center}

\begin{exercice}[Entraînement : régimes et conditions initiales]
On considère l'équation différentielle $y'' + 6y' + 13y = 0$.
\begin{enumerate}
\item Déterminer la nature des solutions (signe de $\Delta$) et écrire la solution générale.
\item Trouver la solution particulière vérifiant $y(0) = 2$ et $y'(0) = -1$.
\item Préciser le régime physique correspondant et décrire l'allure de la courbe.
\end{enumerate}
\end{exercice}

\begin{corrige}[Exercice]
\begin{enumerate}
\item $\Delta = 36 - 52 = -16 < 0$. Racines : $r = \frac{-6 \pm 4i}{2} = -3 \pm 2i$. $\alpha = -3$, $\beta = 2$.
Solution générale : $y(t) = e^{-3t}(\lambda \cos 2t + \mu \sin 2t)$.
\item $y(0) = \lambda = 2$.
$y'(t) = -3e^{-3t}(\lambda \cos 2t + \mu \sin 2t) + e^{-3t}(-2\lambda \sin 2t + 2\mu \cos 2t)$.
$y'(0) = -3\lambda + 2\mu = -1 \Rightarrow -6 + 2\mu = -1 \Rightarrow \mu = 2,5$.
Solution : $\boxed{y(t) = e^{-3t}(2\cos 2t + 2,5\sin 2t)}$.
\item $\Delta < 0$ : régime \textbf{pseudo-périodique}. Oscillations de pulsation $2$ rad/s, amplitude décroissante comme $e^{-3t}$. Retour à l'équilibre rapide avec oscillations.
\end{enumerate}
\end{corrige}

\begin{exercice}[Application : circuit $RLC$ série]
Un circuit $RLC$ série a pour paramètres $L = 0,5\,\text{H}$, $R = 2\,\Omega$, $C = 0,1\,\text{F}$. Il est alimenté par une tension continue $E = 12\,\text{V}$. Le circuit est au repos pour $t=0$ ($q(0)=0, i(0)=0$).
\begin{enumerate}
\item Écrire l'équation différentielle vérifiée par la charge $q(t)$.
\item Résoudre cette équation avec les conditions initiales.
\item Quelle est la charge finale du condensateur ? Interpréter.
\end{enumerate}
\end{exercice}

\begin{corrige}[Exercice]
\begin{enumerate}
\item $L q'' + R q' + \frac{1}{C}q = E \Rightarrow 0,5 q'' + 2 q' + 10 q = 12$.
Multiplier par 2 : $q'' + 4 q' + 20 q = 24$.
\item Homogène : $r^2 + 4r + 20 = 0$, $\Delta = 16 - 80 = -64 < 0$.
$r = -2 \pm 4i$. $q_h(t) = e^{-2t}(\lambda \cos 4t + \mu \sin 4t)$.
Particulière : $c=20 \neq 0$, $q_p = 24/20 = 1,2$.
Générale : $q(t) = e^{-2t}(\lambda \cos 4t + \mu \sin 4t) + 1,2$.
Conditions : $q(0) = \lambda + 1,2 = 0 \Rightarrow \lambda = -1,2$.
$i(t) = q'(t) = e^{-2t}[(-2\lambda + 4\mu)\cos 4t + (-2\mu - 4\lambda)\sin 4t]$.
$i(0) = -2\lambda + 4\mu = 0 \Rightarrow 2,4 + 4\mu = 0 \Rightarrow \mu = -0,6$.
Solution : $\boxed{q(t) = e^{-2t}(-1,2\cos 4t - 0,6\sin 4t) + 1,2}$.
\item Charge finale : $\lim_{t\to+\infty} q(t) = 1,2\,\text{C}$. C'est la charge statique $C \times E = 0,1 \times 12 = 1,2\,\text{C}$. Le régime transitoire (oscillations amorties) s'éteint, le condensateur se charge à la tension source.
\end{enumerate}
\end{corrige}

\coursec{Équations $ay'' + by' + cy = d$ (second membre constant)}

Dans la section précédente, nous avons appris à résoudre complètement l'équation \emph{homogène} (ou \emph{sans second membre}) $ay''+by'+cy=0$ : l'équation caractéristique $ar^2+br+c=0$ fournit, selon le signe de son discriminant, la solution générale de cette équation. Mais dans les situations concrètes de l'atelier ou du chantier, les systèmes physiques sont rarement laissés à eux-mêmes : un générateur impose une tension continue dans un circuit, un poids suspendu à un ressort exerce une force permanente, une machine applique un couple constant. Mathématiquement, cela se traduit par un \cle{second membre} non nul. Cette section est consacrée au cas le plus fréquent et le plus simple : celui où ce second membre est une \emph{constante} $d$.

\courssub{Présentation et situations concrètes}

\begin{definition}[Équation différentielle du second ordre avec second membre constant]
Soient $a$, $b$, $c$ et $d$ quatre réels, avec $a\neq 0$. On appelle \cle{équation différentielle linéaire du second ordre à coefficients constants avec second membre constant} toute équation de la forme
\[ (E)\ :\ ay''+by'+cy=d \]
d'inconnue la fonction $y$, deux fois dérivable sur $\mathbb{R}$. L'équation $(E_0)\ :\ ay''+by'+cy=0$ est appelée \cle{équation homogène associée} à $(E)$.
\end{definition}

\begin{exemplebox}[Deux situations de la vie technique]
\begin{itemize}
\item \textbf{Circuit électrique.} Un circuit comportant une résistance, une bobine et un condensateur, alimenté par un générateur délivrant une tension continue $U$, a une charge électrique $q(t)$ qui vérifie une équation du type $Lq''+Rq'+\dfrac{1}{C}q=U$. Le second membre $U$ est constant.
\item \textbf{Masse suspendue à un ressort.} Une masse accrochée à un ressort vertical subit son propre poids, force constante. Son élongation $x(t)$ vérifie $mx''+fx'+kx=mg$ : le second membre $d=mg$ est le poids, permanent.
\end{itemize}
\end{exemplebox}

L'idée directrice de toute la section est la suivante : un second membre constant « décale » le système vers une \emph{position d'équilibre} constante. On va donc chercher d'abord cette position d'équilibre (une solution constante), puis lui ajouter les solutions de l'équation homogène déjà étudiées.

\courssub{Recherche d'une solution particulière constante}

Cherchons une solution de $(E)$ sous la forme la plus simple possible : une fonction \textbf{constante}, $f_p(t)=k$ pour tout réel $t$. Alors $f_p'(t)=0$ et $f_p''(t)=0$. En reportant dans $(E)$ :
\[ a\times 0+b\times 0+c\times k=d \quad\Longleftrightarrow\quad ck=d. \]

\begin{propriete}[Solution particulière constante, cas $c\neq 0$]
Si $c\neq 0$, la fonction constante définie par
\[ f_p(t)=\dfrac{d}{c} \]
est une solution particulière de l'équation $(E)\ :\ ay''+by'+cy=d$.
\end{propriete}

\textbf{Vérification immédiate.} $f_p$ est constante donc $f_p'=f_p''=0$, et $af_p''+bf_p'+cf_p=c\times\dfrac{d}{c}=d$. La fonction $f_p$ convient bien. $\blacksquare$

\begin{attention}[Cas $c=0$ : la solution particulière n'est pas constante]
Si $c=0$, écrire $\dfrac{d}{c}$ est une \textbf{division par zéro} : c'est interdit. L'équation devient $ay''+by'=d$, et aucune fonction constante ne peut convenir (le membre de gauche serait nul). On cherche alors une solution particulière \emph{polynomiale} :
\begin{itemize}
\item si $c=0$ et $b\neq 0$ : on cherche $f_p$ de la forme $f_p(t)=\alpha t$ (degré 1). En reportant, $b\alpha=d$, d'où $f_p(t)=\dfrac{d}{b}\,t$ ;
\item si $c=0$ et $b=0$ : l'équation est $ay''=d$ ; on cherche $f_p$ de la forme $\alpha t^2$ (degré 2). En reportant, $2a\alpha=d$, d'où $f_p(t)=\dfrac{d}{2a}\,t^2$.
\end{itemize}
Ce cas est un complément utile ; au Probatoire et au Baccalauréat technique, le cas $c\neq 0$ est le cas de référence.
\end{attention}

\courssub{Théorème fondamental : structure de la solution générale}

\begin{propriete}[Théorème — solution générale de $(E)$]
La solution générale de l'équation $(E)\ :\ ay''+by'+cy=d$ est la somme :
\begin{itemize}
\item d'une \textbf{solution particulière} $f_p$ de $(E)$ (par exemple $\dfrac{d}{c}$ si $c\neq 0$) ;
\item et de la \textbf{solution générale} $f_0$ de l'équation homogène associée $(E_0)\ :\ ay''+by'+cy=0$.
\end{itemize}
Autrement dit : $f=f_p+f_0$. Cette démonstration est exigible au Probatoire et au Baccalauréat technique.
\end{propriete}

\begin{experience}[Démonstration guidée du théorème]
On raisonne en deux temps.
\begin{enumerate}
\item \textbf{Toute somme $f_p+f_0$ est solution de $(E)$.} Soit $f_0$ une solution quelconque de $(E_0)$ et posons $f=f_p+f_0$. Par linéarité de la dérivation :
\begin{align*}
af''+bf'+cf &= a(f_p''+f_0'')+b(f_p'+f_0')+c(f_p+f_0)\\
&= \underbrace{(af_p''+bf_p'+cf_p)}_{=\,d\ \text{car } f_p \text{ solution de } (E)} + \underbrace{(af_0''+bf_0'+cf_0)}_{=\,0\ \text{car } f_0 \text{ solution de } (E_0)} = d.
\end{align*}
Donc $f$ est bien solution de $(E)$.
\item \textbf{Toute solution de $(E)$ est de cette forme.} Soit $g$ une solution quelconque de $(E)$. Posons $h=g-f_p$ (la \emph{différence} de deux solutions de $(E)$). Alors :
\[ ah''+bh'+ch=(ag''+bg'+cg)-(af_p''+bf_p'+cf_p)=d-d=0. \]
Donc $h$ est solution de l'équation homogène $(E_0)$, et $g=f_p+h$ : toute solution de $(E)$ s'écrit bien comme somme d'une solution particulière et d'une solution de $(E_0)$. $\blacksquare$
\end{enumerate}
\end{experience}

\begin{aretenir}[À retenir]
Pour résoudre $ay''+by'+cy=d$ avec $c\neq 0$ :
\[ \boxed{\ \text{Solution générale de } (E)\ =\ \underbrace{\dfrac{d}{c}}_{\text{régime permanent}}\ +\ \underbrace{\text{solution générale de } (E_0)}_{\text{régime transitoire}}\ } \]
Les constantes $\lambda$ et $\mu$ contenues dans la solution de $(E_0)$ se déterminent \textbf{ensuite}, à l'aide des conditions initiales, sur la solution \emph{complète}.
\end{aretenir}

\courssub{Méthode de résolution complète}

\begin{methode}[Résoudre $ay''+by'+cy=d$ en 4 étapes]
\begin{enumerate}
\item \textbf{Équation homogène.} Résoudre $(E_0)\ :\ ay''+by'+cy=0$ grâce à l'équation caractéristique $ar^2+br+c=0$ (selon le signe de $\Delta$ : deux exponentielles, exponentielle et polynôme, ou exponentielle amortie oscillante).
\item \textbf{Solution particulière.} Si $c\neq 0$, prendre la constante $f_p=\dfrac{d}{c}$.
\item \textbf{Somme.} Écrire la solution générale complète $f(t)=\dfrac{d}{c}+f_0(t)$.
\item \textbf{Conditions initiales.} Déterminer les constantes $\lambda$ et $\mu$ en appliquant les conditions initiales (par exemple $f(0)$ et $f'(0)$) à la solution \emph{complète} $f$, jamais à $f_0$ seule.
\end{enumerate}
\end{methode}

\begin{exoresolu}[Résolution complète avec conditions initiales]
\textbf{Énoncé.} Résoudre l'équation différentielle
\[ (E)\ :\ y''+y'-2y=6 \]
avec les conditions initiales $y(0)=0$ et $y'(0)=0$.
\tcblower
\textbf{Solution détaillée.}

\textbf{Étape 1 : équation homogène.} L'équation caractéristique de $(E_0)\ :\ y''+y'-2y=0$ est
\[ r^2+r-2=0. \]
$\Delta=1+8=9>0$, donc deux racines réelles : $r_1=\dfrac{-1-3}{2}=-2$ et $r_2=\dfrac{-1+3}{2}=1$. La solution générale de $(E_0)$ est
\[ f_0(t)=\lambda e^{-2t}+\mu e^{t},\quad (\lambda,\mu)\in\mathbb{R}^2. \]

\textbf{Étape 2 : solution particulière.} Ici $c=-2\neq 0$ et $d=6$, donc
\[ f_p(t)=\dfrac{d}{c}=\dfrac{6}{-2}=-3. \]

\textbf{Étape 3 : solution générale de $(E)$.}
\[ f(t)=-3+\lambda e^{-2t}+\mu e^{t}. \]

\textbf{Étape 4 : conditions initiales.} On calcule $f'(t)=-2\lambda e^{-2t}+\mu e^{t}$. Alors :
\[ f(0)=-3+\lambda+\mu=0 \quad\text{et}\quad f'(0)=-2\lambda+\mu=0. \]
On résout le système
\[ \begin{cases} \lambda+\mu=3\\ -2\lambda+\mu=0 \end{cases} \]
En soustrayant la deuxième équation de la première : $3\lambda=3$, donc $\lambda=1$, puis $\mu=2$.

\textbf{Conclusion.} La solution cherchée est
\[ f(t)=-3+e^{-2t}+2e^{t}. \]
\textbf{Vérification.} $f(0)=-3+1+2=0$ ; $f'(t)=-2e^{-2t}+2e^t$ donc $f'(0)=0$ ; et $f''(t)=4e^{-2t}+2e^t$, d'où $f''+f'-2f=(4e^{-2t}+2e^t)+(-2e^{-2t}+2e^t)-2(-3+e^{-2t}+2e^t)=6$. La solution convient.

\textbf{Remarque.} Dans cet exemple, la racine $r_2=1$ est positive : le terme $2e^{t}$ diverge quand $t\to+\infty$. Ce système est \emph{instable} (pas d'amortissement suffisant). Les figures ci-dessous illustrent le cas stable (parties réelles négatives), typique des systèmes techniques réels (frottements, résistance).
\end{exoresolu}

\begin{attention}[Erreur fréquente : conditions initiales appliquées trop tôt]
Une erreur classique consiste à appliquer les conditions initiales à la solution de l'équation homogène $f_0$ \emph{seule}, avant d'ajouter la solution particulière. Dans l'exemple précédent, écrire $\lambda+\mu=0$ (au lieu de $-3+\lambda+\mu=0$) donnerait des constantes fausses. \textbf{Règle d'or :} on détermine $\lambda$ et $\mu$ \emph{après} avoir écrit la solution générale complète $f=f_p+f_0$, en utilisant $f$ et $f'$.
\end{attention}

\courssub{Interprétation physique : régime transitoire et régime permanent}

Dans les situations technologiques, les racines de l'équation caractéristique ont le plus souvent des \textbf{parties réelles négatives} (frottements, résistances : le système s'amortit). La partie $f_0$ de la solution, formée d'exponentielles décroissantes (éventuellement oscillantes), tend alors vers $0$ : c'est le \cle{régime transitoire}, qui disparaît avec le temps. Il reste la solution particulière constante $\dfrac{d}{c}$ : c'est le \cle{régime permanent}, la valeur vers laquelle le système se stabilise.

\begin{popfigure}
\begin{center}
\begin{tikzpicture}
\begin{axis}[
width=0.9\linewidth, height=6.2cm,
xlabel={$t$}, ylabel={$y(t)$},
xmin=0, xmax=6, ymin=-0.5, ymax=5,
axis lines=left,
xtick={0,1,2,3,4,5,6},
ytick={4},
yticklabels={$y=4$},
grid=both, grid style={popline!40},
legend style={at={(0.97,0.35)},anchor=east,draw=popline}
]
\addplot[popcurve,domain=0:6,samples=200]{4-4*exp(-x)-4*x*exp(-x)};
\addlegendentry{solution $y(t)=4-4e^{-t}-4t\,e^{-t}$}
\addplot[poporange,dashed,thick,domain=0:6]{4};
\addlegendentry{régime permanent $y=\dfrac{d}{c}=4$}
\end{axis}
\end{tikzpicture}
\end{center}
\legende{Solution de $y''+2y'+y=4$ avec $y(0)=0$ et $y'(0)=0$. L'équation caractéristique $r^2+2r+1=0$ admet la racine double $r=-1$ ; la solution complète $y(t)=4+\big(\lambda+\mu t\big)e^{-t}$ avec $\lambda=-4$, $\mu=-4$ converge vers le régime permanent $\dfrac{d}{c}=4$ (asymptote en pointillés) : la partie homogène $-4e^{-t}-4te^{-t}$ est le régime transitoire, qui s'amortit.}
\end{popfigure}

\begin{popfigure}
\begin{center}
\begin{tikzpicture}[scale=1.05]
% plafond
\fill[pattern=north east lines, pattern color=popmuted] (-2,0) rectangle (2,0.25);
\draw[popdark,thick] (-2,0) -- (2,0);
% ressort (zigzag)
\draw[popblue,thick,decorate,decoration={coil,segment length=5pt,amplitude=4pt}] (0,0) -- (0,-2.4);
% masse
\fill[poporangeL,draw=poporange,thick] (-0.55,-3.1) rectangle (0.55,-2.4);
\node[popdark] at (0,-2.75) {\small masse $m$};
% position sans charge
\draw[popmuted,dashed] (-1.6,-1.2) -- (1.6,-1.2);
\node[popmuted,right] at (1.6,-1.2) {\footnotesize repos à vide};
% position d'équilibre avec charge
\draw[popgreen,thick,dashed] (-1.6,-2.75) -- (1.6,-2.75);
\node[popgreen,right] at (1.6,-2.75) {\footnotesize équilibre : $x=\dfrac{d}{c}=\dfrac{mg}{k}$};
% flèche décalage
\draw[<->,popink,thick] (1.1,-1.2) -- (1.1,-2.75) node[midway,left,popink]{\footnotesize $\dfrac{d}{c}$};
% poids
\draw[->,popdark,thick] (0,-3.1) -- (0,-3.9) node[below]{\footnotesize poids $\vec{P}=m\vec{g}$ (force constante)};
\end{tikzpicture}
\end{center}
\legende{Masse suspendue à un ressort vertical. Le poids, force constante $d=mg$, décale la position d'équilibre de la quantité $\dfrac{d}{c}=\dfrac{mg}{k}$ : c'est exactement la solution particulière constante de $mx''+fx'+kx=mg$. Autour de cette position, les oscillations amorties forment le régime transitoire.}
\end{popfigure}

\begin{savaistu}[Ce que mesure l'appareil après stabilisation]
En électricité, lorsqu'on branche un circuit sur un générateur de tension continue, le courant et les tensions « bougent » pendant un court instant : c'est le régime transitoire. Puis tout se stabilise, et le multimètre affiche une valeur fixe. Cette valeur mesurée est précisément le régime permanent, c'est-à-dire la solution particulière $\dfrac{d}{c}$ de l'équation différentielle du circuit. Les techniciens savent qu'il faut attendre la fin du transitoire avant de relever une mesure fiable : les mathématiques de cette leçon expliquent pourquoi, et permettent même de prévoir la durée de cette attente grâce aux exponentielles de la solution homogène.
\end{savaistu}

\begin{exercice}[À vous de jouer]
On considère l'équation différentielle $(E)\ :\ y''+3y'+2y=6$.
\begin{enumerate}
\item Résoudre l'équation homogène associée $(E_0)$.
\item Déterminer une solution particulière constante de $(E)$.
\item En déduire la solution générale de $(E)$.
\item Déterminer la solution $f$ de $(E)$ vérifiant $f(0)=0$ et $f'(0)=0$. Vers quelle valeur $f(t)$ tend-elle lorsque $t$ devient grand ? Interpréter.
\end{enumerate}
\end{exercice}

\begin{corrige}[Exercice : régime transitoire et permanent]
\begin{enumerate}
\item Équation caractéristique : $r^2+3r+2=0$, $\Delta=9-8=1$, racines $r_1=-2$ et $r_2=-1$. Donc $f_0(t)=\lambda e^{-2t}+\mu e^{-t}$.
\item $c=2\neq 0$, donc $f_p=\dfrac{d}{c}=\dfrac{6}{2}=3$.
\item Solution générale : $f(t)=3+\lambda e^{-2t}+\mu e^{-t}$.
\item $f'(t)=-2\lambda e^{-2t}-\mu e^{-t}$. Conditions : $f(0)=3+\lambda+\mu=0$ et $f'(0)=-2\lambda-\mu=0$. De la deuxième équation $\mu=-2\lambda$ ; en reportant : $3+\lambda-2\lambda=0$, donc $\lambda=3$ et $\mu=-6$. D'où
\[ f(t)=3+3e^{-2t}-6e^{-t}. \]
Lorsque $t\to+\infty$, les deux exponentielles tendent vers $0$ (régime transitoire amorti) : $f(t)$ tend vers $3$, le régime permanent $\dfrac{d}{c}$.
\end{enumerate}
\end{corrige}

\coursec{Applications à des situations concrètes et synthèse}

Dans la section précédente, nous avons appris à résoudre complètement les équations du type $ay''+by'+cy=d$, où $d$ est une constante : on résout d'abord l'équation sans second membre, puis on ajoute une solution particulière constante. Nous disposons donc désormais de toute la « boîte à outils » de la leçon : les équations $y'=ay$, les équations $ay''+by'+cy=0$ et les équations $ay''+by'+cy=d$. Mais à quoi servent ces outils dans la vie d'un technicien ? C'est toute l'ambition de cette dernière section : montrer que ces équations décrivent des phénomènes concrets de l'atelier, du chantier ou du laboratoire, et apprendre à rédiger une démarche complète, de l'énoncé physique jusqu'à l'interprétation du résultat.

\courssub{La démarche de modélisation en quatre temps}

Lorsqu'un technicien étudie un phénomène qui évolue dans le temps (une température qui baisse, un courant qui s'établit, une suspension qui oscille), il ne connaît pas directement la grandeur cherchée : il connaît une \emph{loi} qui relie cette grandeur à sa vitesse de variation. Traduire cette loi, c'est écrire une \cle{équation différentielle}.

\begin{methode}[Fiche de modélisation : les quatre temps]
Pour résoudre un problème concret conduisant à une équation différentielle, on suit toujours la même démarche :
\begin{enumerate}
\item \textbf{Choisir la grandeur inconnue.} On nomme la grandeur étudiée (par exemple $\theta(t)$, $i(t)$, $q(t)$), on précise son unité et la variable (le temps $t$, en général).
\item \textbf{Traduire la loi en équation différentielle.} L'énoncé donne une relation de proportionnalité ou une loi physique (loi de Newton, loi d'Ohm, principe de la dynamique) : on la traduit en une équation portant sur la fonction et ses dérivées.
\item \textbf{Résoudre.} On identifie le type d'équation ($y'=ay$, $ay''+by'+cy=0$ ou $ay''+by'+cy=d$), on écrit la solution générale, puis on détermine les constantes grâce aux \cle{conditions initiales}.
\item \textbf{Interpréter et critiquer.} On revient au problème concret : que vaut la grandeur au temps demandé ? Le résultat est-il plausible ? Que se passe-t-il quand $t$ devient grand ? Le modèle a-t-il des limites ?
\end{enumerate}
\end{methode}

\begin{attention}[Erreur fréquente : mal traduire l'énoncé]
La difficulté n'est presque jamais la résolution, mais la \textbf{traduction}. Deux pièges reviennent constamment :
\begin{itemize}
\item Confondre « la vitesse de variation est proportionnelle à $y$ » (qui donne $y'=ay$) avec « la variation est proportionnelle à $y'$ ». Le mot \emph{vitesse} ou \emph{taux de variation} désigne toujours la \textbf{dérivée} : c'est elle qui est proportionnelle à la grandeur.
\item Oublier le signe : si la grandeur \emph{décroît}, le coefficient est \textbf{négatif}. Un corps qui se refroidit vérifie $\theta'=-k(\theta-\theta_{\text{amb}})$ avec $k>0$, et non $+k$.
\end{itemize}
Soulignez dans l'énoncé les mots « proportionnel à », « varie », « décroît » avant d'écrire la moindre équation.
\end{attention}

\courssub{Application 1 : le refroidissement d'un objet (loi de Newton)}

\textbf{Intuition.} Une tasse de thé brûlant posée sur la table d'un atelier ne se refroidit pas à vitesse constante : elle refroidit vite au début (quand l'écart avec la pièce est grand), puis de plus en plus lentement. C'est la \cle{loi de refroidissement de Newton} : la vitesse de variation de la température est proportionnelle à l'écart entre la température de l'objet et la température ambiante.

\begin{propriete}[Loi de Newton et équation associée]
Si $\theta(t)$ désigne la température d'un corps à l'instant $t$, plongé dans un milieu à température constante $\theta_a$, alors il existe un réel $k>0$ tel que :
\[ \theta'(t)=-k\bigl(\theta(t)-\theta_a\bigr). \]
En posant $y=\theta-\theta_a$, on obtient $y'=-ky$, équation du type $y'=ay$ avec $a=-k$. La solution générale est donc :
\[ \theta(t)=\theta_a+C\,e^{-kt}, \qquad C\in\mathbb{R}. \]
Si $\theta(0)=\theta_0$, alors $C=\theta_0-\theta_a$ et $\theta(t)=\theta_a+(\theta_0-\theta_a)e^{-kt}$.
\end{propriete}

L'interprétation est immédiate : quand $t\to+\infty$, $e^{-kt}\to 0$, donc $\theta(t)\to\theta_a$ : le corps finit par prendre la température ambiante, sans jamais la dépasser. Le modèle est cohérent.

\begin{exoresolu}[Refroidissement d'une pièce métallique]
Un corps à $80$ \textdegree C est placé dans une pièce à $25$ \textdegree C. Sa température $\theta(t)$ (en degrés Celsius, $t$ en minutes) vérifie $\theta'=-0{,}1(\theta-25)$ et $\theta(0)=80$.
\begin{enumerate}
\item Déterminer $\theta(t)$.
\item À quel instant la température vaut-elle $40$ \textdegree C ? (On donne $\ln 3\simeq 1{,}10$ et $\ln(11/3)\simeq 1{,}30$.)
\end{enumerate}
\tcblower
\textbf{Solution.}
\begin{enumerate}
\item On pose $y=\theta-25$. Alors $y'=\theta'$, et l'équation devient $y'=-0{,}1\,y$ : c'est une équation du type $y'=ay$ avec $a=-0{,}1$. D'après le théorème du cours, les solutions sont les fonctions $y(t)=Ce^{-0{,}1t}$, $C\in\mathbb{R}$. Donc :
\[ \theta(t)=25+Ce^{-0{,}1t}. \]
La condition initiale donne $\theta(0)=25+C=80$, d'où $C=55$. Ainsi :
\[ \boxed{\theta(t)=25+55\,e^{-0{,}1t}}. \]
\item On résout $\theta(t)=40$ :
\begin{align*}
25+55e^{-0{,}1t}&=40\\
55e^{-0{,}1t}&=15\\
e^{-0{,}1t}&=\frac{15}{55}=\frac{3}{11}\\
-0{,}1\,t&=\ln\frac{3}{11}=-\ln\frac{11}{3}\\
t&=10\ln\frac{11}{3}\simeq 10\times 1{,}30=13.
\end{align*}
Le corps atteint $40$ \textdegree C au bout d'environ \textbf{13 minutes}. Vérification de cohérence : $40$ \textdegree C est bien compris entre $25$ et $80$, et le temps est positif.
\end{enumerate}
\end{exoresolu}

\begin{popfigure}
\begin{center}
\begin{tikzpicture}
\begin{axis}[
width=0.9\linewidth, height=6.5cm,
xlabel={Temps $t$ (min)}, ylabel={Température $\theta$ (\textdegree C)},
xmin=0, xmax=45, ymin=20, ymax=85,
grid=both, grid style={popline!40},
axis lines=left, legend style={at={(0.97,0.6)},anchor=east,draw=popline}
]
\addplot[popcurve,domain=0:45,samples=200]{25+55*exp(-0.1*x)};
\addlegendentry{$\theta(t)=25+55e^{-0{,}1t}$}
\addplot[dashed,poporange,thick,domain=0:45]{25};
\addlegendentry{Température ambiante $\theta_a=25$}
\addplot[dashed,popmuted] coordinates {(13,20) (13,40) (0,40)};
\addplot[only marks,mark=*,popink] coordinates {(13,40)};
\node[popink,anchor=south west] at (axis cs:13,40) {\small $\theta=40$ \textdegree C, $t\simeq13$ min};
\end{axis}
\end{tikzpicture}
\end{center}
\legende{Refroidissement d'un corps de $80$ \textdegree C vers la température ambiante de $25$ \textdegree C : la décroissance est rapide au début, puis ralentit ; la courbe se rapproche de l'asymptote $\theta=25$ sans jamais l'atteindre.}
\end{popfigure}

\courssub{Application 2 : circuit RC ou RL (premier ordre)}

\textbf{Intuition.} Quand on branche un condensateur ou une bobine dans un circuit électrique, le courant ou la tension ne s'établit pas instantanément : il y a un \emph{régime transitoire}. Les lois de l'électricité (loi d'Ohm, loi des mailles) conduisent naturellement à une équation différentielle du premier ordre.

\begin{exemplebox}[Charge d'un condensateur (circuit RC)]
Un condensateur de capacité $C$ se charge à travers une résistance $R$ sous une tension constante $E$. La tension $u(t)$ aux bornes du condensateur vérifie la loi des mailles :
\[ RC\,u'(t)+u(t)=E, \qquad u(0)=0. \]
En posant $y=u-E$, on obtient $RC\,y'+y=0$, c'est-à-dire $y'=-\dfrac{1}{RC}y$ : équation du type $y'=ay$ avec $a=-\dfrac{1}{RC}$. Donc $y(t)=K e^{-t/(RC)}$, puis :
\[ u(t)=E+Ke^{-t/(RC)}. \]
La condition $u(0)=0$ donne $K=-E$, d'où :
\[ u(t)=E\bigl(1-e^{-t/(RC)}\bigr). \]
\textbf{Interprétation :} la tension monte progressivement de $0$ vers $E$ ; le produit $\tau=RC$, appelé \cle{constante de temps}, mesure la rapidité de la charge. Plus $R$ ou $C$ est grand, plus la charge est lente. C'est exactement ce qui se passe quand on charge un téléphone.
\end{exemplebox}

De même, dans un circuit RL (résistance + bobine), l'intensité vérifie $L\,i'+Ri=E$, qui est une équation du premier ordre à second membre constant : on la résout par le même changement de fonction $y=i-\dfrac{E}{R}$.

\courssub{Application 3 : circuit RLC et système masse-ressort amorti (second ordre)}

\textbf{Intuition.} Une voiture passe sur un nid-de-poule : sa suspension oscille puis se stabilise. Selon l'état des amortisseurs, le retour à l'équilibre peut être mou, rapide, ou accompagné d'oscillations. Mathématiquement, la position $x(t)$ de la masse vérifie une équation du type :
\[ m\,x''+f\,x'+k\,x=0, \]
où $m$ est la masse, $f$ le coefficient de frottement (amortisseur) et $k$ la raideur du ressort. De même, dans un circuit RLC série, la charge $q(t)$ du condensateur vérifie $Lq''+Rq'+\dfrac{1}{C}q=0$. \textbf{C'est la même équation} : l'équation caractéristique $ar^2+br+c=0$ et le signe de son discriminant $\Delta=b^2-4ac$ commandent trois régimes.

\begin{aretenir}[Les trois régimes d'un système du second ordre]
Pour $ay''+by'+cy=0$ avec $a>0$, $b>0$, $c>0$ (système amorti), d'équation caractéristique $ar^2+br+c=0$ de discriminant $\Delta=b^2-4ac$ :
\begin{itemize}
\item $\Delta>0$ : deux racines réelles $r_1,r_2<0$ ; solution $y=C_1e^{r_1t}+C_2e^{r_2t}$. \textbf{Régime apériodique} : retour à l'équilibre sans oscillation (amortissement fort, amortisseurs très durs).
\item $\Delta=0$ : racine double $r_0<0$ ; solution $y=(C_1t+C_2)e^{r_0t}$. \textbf{Régime critique} : retour à l'équilibre le plus rapide possible sans oscillation.
\item $\Delta<0$ : racines complexes $\alpha\pm i\beta$ avec $\alpha<0$ ; solution $y=e^{\alpha t}\bigl(C_1\cos(\beta t)+C_2\sin(\beta t)\bigr)$. \textbf{Régime pseudo-périodique} : oscillations dont l'amplitude décroît exponentiellement (amortisseurs usés : la voiture « rebondit »).
\end{itemize}
Dans les trois cas, $y(t)\to 0$ quand $t\to+\infty$ : le système revient à l'équilibre.
\end{aretenir}

\begin{exemplebox}[Suspension d'un véhicule]
La suspension d'un véhicule est modélisée par $x''+4x'+4x=0$, avec $x(0)=5$ (cm) et $x'(0)=0$.
Équation caractéristique : $r^2+4r+4=0$, soit $(r+2)^2=0$ : racine double $r_0=-2$. On est en \textbf{régime critique} :
\[ x(t)=(C_1t+C_2)e^{-2t}. \]
$x(0)=C_2=5$ ; puis $x'(t)=C_1e^{-2t}-2(C_1t+C_2)e^{-2t}$, donc $x'(0)=C_1-2C_2=0$, d'où $C_1=10$. Finalement :
\[ x(t)=(10t+5)e^{-2t}. \]
Interprétation : la suspension revient vers sa position d'équilibre sans osciller, c'est le réglage idéal recherché par le mécanicien.
\end{exemplebox}

\courssub{Application 4 : concentration d'un médicament avec apport constant}

\textbf{Intuition.} Un patient reçoit un médicament par perfusion à débit constant, tandis que l'organisme élimine le médicament proportionnellement à la quantité présente. Deux phénomènes s'opposent : l'apport (constant) et l'élimination (proportionnelle). La quantité $q(t)$ de médicament dans le sang vérifie une équation \textbf{avec second membre constant} :
\[ q'(t)=-k\,q(t)+d, \qquad k>0,\ d>0. \]

\begin{exoresolu}[Perfusion à débit constant]
La quantité de médicament (en mg) dans le sang d'un patient vérifie $q'=-0{,}2\,q+4$, avec $q(0)=0$.
\begin{enumerate}
\item Déterminer $q(t)$.
\item Vers quelle valeur la quantité se stabilise-t-elle ?
\end{enumerate}
\tcblower
\textbf{Solution.}
\begin{enumerate}
\item C'est une équation du type $y'=ay+d$ avec $a=-0{,}2$ et $d=4$.
\textbf{Solution particulière constante :} on cherche $q_p$ telle que $0=-0{,}2\,q_p+4$, d'où $q_p=\dfrac{4}{0{,}2}=20$.
\textbf{Équation sans second membre :} $q'=-0{,}2\,q$, de solutions $Ce^{-0{,}2t}$.
\textbf{Solution générale :} $q(t)=20+Ce^{-0{,}2t}$.
La condition $q(0)=0$ donne $20+C=0$, soit $C=-20$ :
\[ \boxed{q(t)=20\bigl(1-e^{-0{,}2t}\bigr)}. \]
\item Quand $t\to+\infty$, $e^{-0{,}2t}\to 0$, donc $q(t)\to 20$ mg : c'est la \textbf{concentration d'équilibre}, où l'élimination compense exactement l'apport. Le médecin sait ainsi quelle dose sera maintenue en régime permanent.
\end{enumerate}
\end{exoresolu}

\courssub{Rédiger une démarche complète et justifiée au Probatoire}

Au Probatoire comme au Baccalauréat technique, la qualité de la rédaction compte autant que le résultat. Une solution « tombée du ciel » sans justification perd des points, même si elle est exacte.

\begin{methode}[Rédiger une justification complète]
\begin{enumerate}
\item \textbf{Annoncer le théorème utilisé.} « D'après le cours, les solutions de l'équation $y'=ay$ sont les fonctions $t\mapsto Ce^{at}$, $C\in\mathbb{R}$ » (ou : « l'équation caractéristique associée à $ay''+by'+cy=0$ est $ar^2+br+c=0$ »).
\item \textbf{Vérifier les hypothèses.} Identifier le type de l'équation, calculer le discriminant si nécessaire, et justifier le cas dans lequel on se trouve ($\Delta>0$, $\Delta=0$ ou $\Delta<0$).
\item \textbf{Exploiter les conditions initiales} pour déterminer les constantes, en montrant les calculs.
\item \textbf{Conclure proprement} : encadrer ou souligner le résultat, donner l'unité, et répondre à la question posée en français (« le corps atteint $40$ \textdegree C au bout de 13 minutes »).
\end{enumerate}
\end{methode}

\begin{attention}[Pièges de rédaction fréquents]
\begin{itemize}
\item Écrire la solution générale \textbf{avant} d'avoir vérifié le signe de $\Delta$ : le discriminant commande la forme de la solution, il doit être calculé et commenté.
\item Oublier la solution particulière constante dans une équation avec second membre : la solution générale est la \emph{somme} des deux.
\item Donner un temps négatif ou une température inférieure à la température ambiante sans s'alarmer : toujours contrôler la \textbf{plausibilité physique} du résultat.
\end{itemize}
\end{attention}

\courssub{Bilan synthétique de la leçon}

\begin{popfigure}
\begin{center}
\begin{tikzpicture}[
box/.style={draw=popblue, fill=popblueL, rounded corners, align=center, text width=3.1cm, minimum height=1.6cm, font=\small},
boxb/.style={draw=poporange, fill=poporangeL, rounded corners, align=center, text width=3.4cm, minimum height=1.6cm, font=\small},
boxc/.style={draw=poppurple, fill=poppurpleL, rounded corners, align=center, text width=3.6cm, minimum height=1.6cm, font=\small},
boxd/.style={draw=popgreen, fill=popgreenL, rounded corners, align=center, text width=3.6cm, minimum height=1.6cm, font=\small},
arr/.style={-{Stealth[length=2.5mm]}, thick, popdark}
]
\node[box, align=center] (t1) at (0,0){\textbf{Type 1}\\[2pt] $y'=ay$};
\node[boxb, align=center] (t2) at (4.6,0){\textbf{Type 2}\\[2pt] $ay''+by'+cy=0$};
\node[boxc, align=center] (t3) at (9.4,0){\textbf{Type 3}\\[2pt] $ay''+by'+cy=d$};

\node[box, fill=white] (s1) at (0,-2.6) {Pas d'équation caractéristique : lecture directe de $a$};
\node[boxb, fill=white] (s2) at (4.6,-2.6) {Équation caractéristique $ar^2+br+c=0$, discriminant $\Delta$};
\node[boxc, fill=white] (s3) at (9.4,-2.6) {Même équation caractéristique + solution particulière constante $\dfrac{d}{c}$};

\node[box, fill=white] (r1) at (0,-5.4) {$y=Ce^{at}$, $C\in\mathbb{R}$};
\node[boxb, fill=white, align=center] (r2) at (4.6,-5.4){$\Delta>0$ : $C_1e^{r_1t}+C_2e^{r_2t}$\\ $\Delta=0$ : $(C_1t+C_2)e^{r_0t}$\\ $\Delta<0$ : $e^{\alpha t}(C_1\cos\beta t+C_2\sin\beta t)$};
\node[boxc, fill=white] (r3) at (9.4,-5.4) {Solution de l'équation sans second membre $+\ \dfrac{d}{c}$};

\node[boxd] (i1) at (0,-8.2) {Refroidissement, circuit RC, RL : retour exponentiel à l'équilibre};
\node[boxd] (i2) at (4.6,-8.2) {Suspension, circuit RLC : trois régimes selon l'amortissement};
\node[boxd] (i3) at (9.4,-8.2) {Perfusion, charge avec source : convergence vers un régime permanent};

\draw[arr] (t1) -- (s1); \draw[arr] (s1) -- (r1); \draw[arr] (r1) -- (i1);
\draw[arr] (t2) -- (s2); \draw[arr] (s2) -- (r2); \draw[arr] (r2) -- (i2);
\draw[arr] (t3) -- (s3); \draw[arr] (s3) -- (r3); \draw[arr] (r3) -- (i3);
\end{tikzpicture}
\end{center}
\legende{Tableau-synthèse de la leçon : pour chaque type d'équation, la méthode de résolution, la forme de la solution générale et l'interprétation physique typique.}
\end{popfigure}

\begin{aretenir}[L'essentiel à retenir]
\begin{itemize}
\item Toute situation concrète se traite en quatre temps : \textbf{grandeur} $\to$ \textbf{équation} $\to$ \textbf{résolution} $\to$ \textbf{interprétation et critique}.
\item $y'=ay$ : solutions $Ce^{at}$ ; croissance si $a>0$, décroissance si $a<0$.
\item $ay''+by'+cy=0$ : on résout l'équation caractéristique $ar^2+br+c=0$ et le signe de $\Delta$ donne la forme de la solution.
\item $ay''+by'+cy=d$ : solution générale $=$ solution sans second membre $+$ solution particulière constante $\dfrac{d}{c}$.
\item Les conditions initiales déterminent les constantes ; la conclusion doit être rédigée en français, avec l'unité, et contrôlée physiquement.
\end{itemize}
\end{aretenir}

\begin{savaistu}[Des équations partout autour de nous]
Les équations différentielles de cette leçon sont de véritables « couteaux suisses » des sciences et des techniques. La même équation $y'=ay$ décrit la charge de la batterie d'un téléphone, la décroissance radioactive utilisée pour dater les objets anciens, et les premiers instants de la diffusion d'une épidémie simple. L'équation $ay''+by'+cy=0$ décrit à la fois les suspensions d'un véhicule sur les pistes du Cameroun, les oscillations d'un circuit électrique RLC et les vibrations d'un bâtiment soumis au vent. Comprendre ces équations, c'est comprendre le fonctionnement d'objets très différents avec un seul et même outil mathématique : c'est toute la puissance de l'abstraction.
\end{savaistu}

\begin{exercice}[Pour s'entraîner]
Un moteur électrique atteint sa vitesse de régime progressivement : sa vitesse angulaire $\omega(t)$ (en rad/s) vérifie $\omega'=-2\omega+100$, avec $\omega(0)=0$.
\begin{enumerate}
\item Déterminer $\omega(t)$.
\item Vers quelle vitesse le moteur se stabilise-t-il ?
\item Au bout de combien de temps atteint-il $90$ \% de cette vitesse de régime ? (On donne $\ln 10\simeq 2{,}30$.)
\end{enumerate}
\end{exercice}

\begin{corrige}[Exercice : moteur électrique]
\begin{enumerate}
\item Équation du type $y'=ay+d$ avec $a=-2$, $d=100$. Solution particulière constante : $\omega_p=\dfrac{100}{2}=50$. Équation sans second membre : $\omega'=-2\omega$, solutions $Ce^{-2t}$. Solution générale : $\omega(t)=50+Ce^{-2t}$. Condition $\omega(0)=0$ : $50+C=0$, donc $C=-50$ et $\omega(t)=50\bigl(1-e^{-2t}\bigr)$.
\item Quand $t\to+\infty$, $\omega(t)\to 50$ rad/s : c'est la vitesse de régime.
\item On résout $\omega(t)=0{,}9\times 50=45$ : $50(1-e^{-2t})=45$ donne $e^{-2t}=0{,}1$, soit $-2t=\ln 0{,}1=-\ln 10$, d'où $t=\dfrac{\ln 10}{2}\simeq 1{,}15$ s.
\end{enumerate}
\end{corrige}

\newpage

\coursec{Activité d'intégration}

\coursec{Activité d'intégration : Modélisation d'un transitoire électrique (Régime critique)}

\courssub{Objectifs de l'activité}
Cette activité d'intégration vise à mobiliser l'ensemble des savoirs et savoir-faire de la leçon sur les équations différentielles (premier et second ordre à coefficients constants) dans un contexte professionnel authentique. À l'issue de cette séance, l'élève doit être capable de :
\begin{itemize}
    \item Modéliser un phénomène physique réel (transitoire électrique) par une équation différentielle linéaire du second ordre à coefficients constants.
    \item Identifier le type de régime (critique, pseudo-périodique, apériodique) à partir de l'équation caractéristique.
    \item Déterminer la solution générale de l'équation homogène et une solution particulière adaptée au second membre constant.
    \item Exploiter les conditions initiales (état initial du système) pour déterminer les constantes d'intégration.
    \item Interpréter graphiquement et numériquement la solution : temps de réponse, dépassement, régime permanent.
    \item Rédiger un compte-rendu technique structuré, argumenté et accessible à un non-mathématicien (technicien de maintenance).
\end{itemize}

\courssub{Situation : Stabilisation de la tension à l'atelier « Bois \& Précision » de Nkongsamba}
L'atelier de menuiserie \emph{Bois \& Précision}, situé dans la zone industrielle de Nkongsamba, est alimenté par le réseau triphasé 380\,V / 220\,V de la société ENEO. Suite à une coupure intempestive suivie d'un rétablissement brutal du courant (manœuvre de réenclenchement du disjoncteur principal), le responsable de l'atelier observe des scintillements sur les lampes témoins du tableau électrique et craint pour l'intégrité des variateurs de vitesse des machines à commande numérique (scie à panneau, centre d'usinage).

Le technicien de maintenance, M. Fouda, modélise la tension $u(t)$ (en volts) aux bornes du bus principal basse tension après le rétablissement ($t=0$) par l'équation différentielle suivante :
\[
u''(t) + 6\,u'(t) + 9\,u(t) = 1980 \qquad (t \ge 0)
\]
avec les conditions initiales : $u(0) = 0$ (tension nulle avant rétablissement) et $u'(0) = 0$ (pas de variation brutale instantanée du courant dans les inductances du transformateur).

\medskip
\textbf{Données techniques :} La tension nominale de service est $U_n = 220\,\text{V}$. Les équipements sensibles supportent une surtension maximale de $10\,\%$ (soit $242\,\text{V}$) mais doivent atteindre $90\,\%$ de la tension nominale (soit $198\,\text{V}$) dans un délai compatible avec le redémarrage des automates (généralement $< 2\,\text{s}$). Le coefficient d'amortissement du circuit équivalent (résistances des câbles + pertes fer du transformateur) conduit au terme $6u'$, et la résonance propre du circuit $RLC$ équivalent donne le terme $9u$.

\medskip
M. Fouda demande à son équipe de valider ce modèle et de répondre aux questions suivantes :
\begin{enumerate}
    \item Le régime transitoire est-il \emph{critique} (amorti critique) ? Justifier.
    \item Quelle est l'expression exacte de la tension $u(t)$ pour $t \ge 0$ ?
    \item La tension dépasse-t-elle $200\,\text{V}$ ? Si oui, à partir de quel instant $t_1$ (en secondes) ?
    \item Y a-t-il un risque de surtension dommageable (dépassement de $242\,\text{V}$) ?
    \item Rédiger une note technique pour le chef d'atelier concluant sur la conformité du rétablissement.
\end{enumerate}

\courssub{Consignes}
Travailler par groupes de 3 ou 4 élèves. Chaque groupe rendra une fiche de résolution structurée et préparera une présentation orale de 5 minutes au tableau.
\begin{enumerate}
    \item \textbf{Analyse de l'équation caractéristique :} Écrire l'équation caractéristique associée à la partie homogène. Calculer son discriminant $\Delta$ et ses racines. En déduire la nature du régime (critique, apériodique, pseudo-périodique) et écrire la forme de la solution générale de l'équation homogène $u_h(t)$.
    \item \textbf{Recherche d'une solution particulière :} Proposer une forme de solution particulière $u_p$ adaptée au second membre constant. Déterminer sa valeur. Interpréter physiquement cette valeur (régime permanent).
    \item \textbf{Solution générale et conditions initiales :} Écrire la solution générale $u(t) = u_h(t) + u_p$. Utiliser les conditions initiales $u(0)=0$ et $u'(0)=0$ pour déterminer les constantes d'intégration. Donner l'expression finale de $u(t)$.
    \item \textbf{Étude graphique et numérique :} Tracer la courbe représentative de $u(t)$ sur l'intervalle $[0\,;\,3]$ (sur calculatrice, logiciel ou papier millimétré). Repérer graphiquement l'instant $t_1$ où $u(t)$ franchit la valeur $200\,\text{V}$. Affiner $t_1$ par calcul numérique (méthode de dichotomie ou solveur de la calculatrice). Vérifier si $u(t)$ dépasse $242\,\text{V}$.
    \item \textbf{Synthèse et communication :} Rédiger la conclusion technique : le régime critique est-il favorable ici ? Pourquoi ? Quels seraient les avantages/inconvénients d'un régime pseudo-périodique (sous-amorti) ou apériodique (suramorti) pour ce type d'installation ?
\end{enumerate}

\courssub{Ressources à mobiliser}
\begin{propriete}[Équation différentielle linéaire du second ordre à coefficients constants]
Soit l'équation $a y'' + b y' + c y = f(x)$ avec $a,b,c \in \mathbb{R}$, $a \neq 0$.
\begin{itemize}
    \item \textbf{Équation caractéristique :} $a r^2 + b r + c = 0$. Discriminant $\Delta = b^2 - 4ac$.
    \item \textbf{Solution générale de l'homogène ($f=0$) :}
    \begin{itemize}
        \item $\Delta > 0$ : deux racines réelles distinctes $r_1, r_2$ $\rightarrow$ $y_h(x) = A e^{r_1 x} + B e^{r_2 x}$ (Régime \emph{apériodique} / suramorti).
        \item $\Delta = 0$ : une racine réelle double $r_0$ $\rightarrow$ $y_h(x) = (A + B x) e^{r_0 x}$ (Régime \emph{critique}).
        \item $\Delta < 0$ : deux racines complexes conjuguées $\alpha \pm i\beta$ $\rightarrow$ $y_h(x) = e^{\alpha x}(A \cos \beta x + B \sin \beta x)$ (Régime \emph{pseudo-périodique} / sous-amorti).
    \end{itemize}
    \item \textbf{Solution particulière pour $f(x) = k$ (constante) :} On cherche $y_p = C$ (constante). $c C = k \Rightarrow C = k/c$ (si $c \neq 0$).
    \item \textbf{Solution générale :} $y = y_h + y_p$.
\end{itemize}
\end{propriete}

\begin{methode}[Résolution d'une EDL d'ordre 2 avec second membre constant]
\begin{enumerate}
    \item Écrire l'équation caractéristique $a r^2 + b r + c = 0$ et calculer $\Delta$.
    \item Écrire $y_h$ selon le signe de $\Delta$ (cas $\Delta=0$ : racine double $\rightarrow (A+Bx)e^{rx}$).
    \item Trouver $y_p = \frac{k}{c}$ (si $c \neq 0$).
    \item Écrire $y = y_h + y_p$.
    \item Calculer $y'$.
    \item Appliquer les conditions initiales pour obtenir un système linéaire en $A, B$.
    \item Résoudre le système et écrire la solution finale.
\end{enumerate}
\end{methode}

\begin{attention}[Pièges fréquents]
\begin{itemize}
    \item Oublier le terme $B x e^{rx}$ dans le cas $\Delta = 0$ (racine double) : la solution $A e^{rx}$ seule ne contient qu'une constante arbitraire, insuffisante pour satisfaire deux conditions initiales.
    \item Confondre la racine $r$ de l'équation caractéristique et la solution particulière $y_p$.
    \item Erreur de signe dans le calcul de $y'$ pour le cas critique : dériver $(A+Bx)e^{rx}$ demande la règle du produit : $y' = (B + r(A+Bx))e^{rx}$.
    \item Ne pas vérifier que la solution particulière trouvée convient bien (remplacer dans l'équation).
\end{itemize}
\end{attention}

\begin{savaistu}[Contexte technique : Le régime critique en électricité]
En électrotechnique, un circuit $RLC$ série est décrit par $L u'' + R u' + \frac{1}{C}u = E$. L'équation caractéristique est $L r^2 + R r + 1/C = 0$.
Le régime critique correspond à $R^2 = 4L/C$ (amortissement optimal). C'est le régime qui permet de rejoindre le régime permanent \textbf{sans oscillation} (donc pas de surtension) et \textbf{le plus rapidement possible}. C'est le choix idéal pour les alimentations d'équipements sensibles (informatiques, variateurs, automates) car il évite les pics de tension destructeurs tout en minimisant le temps de mise en service. Au Cameroun, la norme NF C 15-100 impose des temps de rétablissement courts pour les circuits de sécurité ; le dimensionnement des câbles (résistance $R$) et des protections vise souvent ce comportement critique.
\end{savaistu}

\courssub{Démarche et corrigé commenté}
\begin{exoresolu}[Résolution détaillée de l'activité : Tension de l'atelier de Nkongsamba]
\noindent \textbf{1. Équation caractéristique et nature du régime}
L'équation homogène associée est $u'' + 6u' + 9u = 0$.
L'équation caractéristique est :
\[
r^2 + 6r + 9 = 0
\]
Le discriminant est $\Delta = 6^2 - 4 \times 1 \times 9 = 36 - 36 = 0$.
Il y a une \textbf{racine réelle double} :
\[
r_0 = \frac{-6}{2} = -3
\]
Comme $\Delta = 0$, le régime est \textbf{critique} (amorti critique).
La solution générale de l'équation homogène est donc :
\[
u_h(t) = (A + B t) e^{-3t} \quad (A, B \in \mathbb{R})
\]

\medskip
\noindent \textbf{2. Solution particulière (Régime permanent)}
Le second membre est la constante $1980$. On cherche une solution particulière constante $u_p = K$.
En remplaçant dans l'équation : $0 + 0 + 9K = 1980 \Rightarrow K = \frac{1980}{9} = 220$.
\[
u_p(t) = 220
\]
\emph{Interprétation :} La tension tend vers $220\,\text{V}$ quand $t \to +\infty$. C'est la tension nominale du réseau. Le régime permanent est atteint.

\medskip
\noindent \textbf{3. Solution générale et détermination des constantes}
La solution générale de l'équation complète est :
\[
u(t) = u_h(t) + u_p(t) = (A + B t) e^{-3t} + 220
\]
Calcul de la dérivée (règle du produit) :
\[
u'(t) = B e^{-3t} + (A + B t)(-3) e^{-3t} = \bigl( B - 3A - 3B t \bigr) e^{-3t}
\]
Application des conditions initiales :
\begin{itemize}
    \item $u(0) = 0$ : $(A + 0) \cdot 1 + 220 = 0 \Rightarrow A = -220$.
    \item $u'(0) = 0$ : $(B - 3A - 0) \cdot 1 = 0 \Rightarrow B - 3A = 0 \Rightarrow B = 3A = -660$.
\end{itemize}
L'expression exacte de la tension est donc :
\[
\boxed{u(t) = 220 - (220 + 660 t) e^{-3t} \qquad (t \ge 0)}
\]

\medskip
\noindent \textbf{4. Étude graphique et numérique : Franchissement de 200 V et vérification surtension}
On étudie la fonction $u(t) = 220 - (220 + 660t)e^{-3t}$ sur $[0\,;\,3]$.
\begin{itemize}
    \item $u(0) = 0$.
    \item $\lim_{t\to+\infty} u(t) = 220$.
    \item $u'(t) = 1980 t e^{-3t}$ (vérification : $u'(t) = -[660 - 3(220+660t)]e^{-3t} = -[660-660-1980t]e^{-3t} = 1980 t e^{-3t}$).
    \item Pour $t > 0$, $u'(t) > 0$. La fonction est \textbf{strictement croissante} sur $[0\,;\,+\infty[$.
    \item Pas d'oscillation, pas de maximum local $> 220$. La tension monte monotone vers 220 V.
\end{itemize}
\emph{Conséquence immédiate :} Il n'y a \textbf{aucun risque de surtension} (la tension ne dépasse jamais 220 V, donc reste bien en dessous de 242 V).

\medskip
\noindent \textbf{Recherche de $t_1$ tel que $u(t_1) = 200$ :}
On résout $220 - (220 + 660 t) e^{-3t} = 200$
$\Leftrightarrow (220 + 660 t) e^{-3t} = 20$
$\Leftrightarrow 220 + 660 t = 20 e^{3t}$
$\Leftrightarrow 11 + 33 t = e^{3t}$ \quad (Équation transcendante, résolution numérique par dichotomie).

On pose $f(t) = e^{3t} - 33t - 11$. On cherche $t$ tel que $f(t)=0$.
\begin{center}
\begin{tabular}{|c|c|c|}\hline
$t$ (s) & $f(t) = e^{3t} - 33t - 11$ & Signe \\ \hline
1.30 & $e^{3.9} - 42.9 - 11 \approx 49.40 - 53.9 = -4.50$ & $-$ \\ \hline
1.33 & $e^{3.99} - 43.89 - 11 \approx 54.00 - 54.89 = -0.89$ & $-$ \\ \hline
1.335 & $e^{4.005} - 44.055 - 11 \approx 54.87 - 55.055 = -0.185$ & $-$ \\ \hline
1.336 & $e^{4.008} - 44.088 - 11 \approx 55.02 - 55.088 = -0.068$ & $-$ \\ \hline
1.337 & $e^{4.011} - 44.121 - 11 \approx 55.17 - 55.121 = +0.049$ & $+$ \\ \hline
\end{tabular}
\end{center}
La solution est encadrée : $t_1 \in [1,336\,;\,1,337]$. On retient $t_1 \approx \mathbf{1,336\,s}$ (soit environ \textbf{1,34 s} arrondi au centième).

\medskip
\noindent \textbf{Représentation graphique (TikZ/pgfplots) :}
\begin{popfigure}
\begin{tikzpicture}
\begin{axis}[
    width=\linewidth, height=8cm,
    axis lines=middle,
    xlabel={$t$ (s)}, ylabel={$u(t)$ (V)},
    xmin=0, xmax=3, ymin=0, ymax=230,
    xtick={0,0.5,1,1.336,1.5,2,2.5,3},
    ytick={0,50,100,150,198,200,220,242},
    grid=major, grid style={dashed, gray!50},
    domain=0:3, samples=200,
    legend style={at={(0.98,0.02)}, anchor=south east, fill=white, fill opacity=0.8},
    clip=false
]
\addplot[popcurve, thick, popblue] {220 - (220 + 660*x)*exp(-3*x)};
\addlegendentry{$u(t) = 220 - (220+660t)e^{-3t}$}
% Ligne régime permanent
\addplot[dashed, popdark, thick] coordinates {(0,220) (3,220)};
\addlegendentry{Régime permanent (220 V)}
% Ligne 200 V
\addplot[dotted, poporange, thick] coordinates {(0,200) (3,200)};
\addlegendentry{Seuil 200 V (90\% $U_n$)}
% Ligne 242 V
\addplot[dotted, popred, thick] coordinates {(0,242) (3,242)};
\addlegendentry{Surtension max (242 V)}
% Point intersection 200V
\addplot[only marks, mark=*, mark size=2.5, poporange] coordinates {(1.336,200)};
\node[above right, poporange, font=\small] at (axis cs:1.336,200) {$t_1 \approx 1,34\,\text{s}$};
% Flèche temps de réponse
\draw[<->, popgreen, thick] (axis cs:0,198) -- (axis cs:1.336,198) node[midway, below, popgreen, font=\small] {$t_{90\%} \approx 1,34\,\text{s}$};
\end{axis}
\end{tikzpicture}
\legende{Évolution de la tension $u(t)$ après rétablissement. Régime critique : montée monotone sans dépassement.}
\end{popfigure}

\medskip
\noindent \textbf{5. Conclusion technique (Note pour le chef d'atelier)}
\begin{quote}
\textbf{Objet :} Analyse du transitoire de tension post-coupure — Atelier Bois \& Précision, Nkongsamba.

Monsieur le Responsable,

Suite à la modélisation du rétablissement de tension sur notre bus principal (équation $u''+6u'+9u=1980$), nous confirmons que le circuit présente un \textbf{régime critique (amorti critique)}.

Les points clés sont :
\begin{itemize}
    \item \textbf{Absence de surtension :} La tension monte de 0 à 220 V de façon strictement croissante, sans jamais osciller ni dépasser la valeur nominale. Les équipements sensibles (variateurs, automates) sont \textbf{protégés} contre les pics de tension.
    \item \textbf{Temps de réponse :} La tension atteint 200 V (90\% de $U_n$) au bout de \textbf{1,34 s}. Ce délai est compatible avec les exigences de redémarrage des automates (généralement $< 2\,\text{s}$).
    \item \textbf{Comparaison régimes :}
    \begin{itemize}
        \item Un régime \emph{pseudo-périodique} (sous-amorti) aurait été plus rapide au début mais aurait généré des oscillations (dépassements de 242 V probables), destructeurs pour l'électronique de puissance.
        \item Un régime \emph{apériodique} (suramorti) n'aurait pas de dépassement mais serait plus lent (ex: $> 2\,\text{s}$ pour 90\%), retardant la reprise de production.
    \end{itemize}
\end{itemize}
\textbf{Avis :} Le dimensionnement actuel (résistances câbles + transformateur) assurant un amortissement critique est \textbf{optimal} pour notre installation. Aucune modification n'est requise. La coupure/rétablissement est gérée en toute sécurité.

Cordialement, \\
L'équipe Maintenance (Groupe \#...)
\end{quote}
\end{exoresolu}

\courssub{Bilan de l'activité}
Cette activité a permis de mettre en évidence les points essentiels suivants :
\begin{itemize}
    \item \textbf{Modélisation :} Une situation technique réelle (transitoire électrique) se traduit naturellement par une EDL du second ordre à coefficients constants avec second membre constant.
    \item \textbf{Discriminant et régimes :} Le signe de $\Delta$ dicte la forme de la solution homogène et le comportement physique (oscillations ou non, rapidité). Le cas $\Delta = 0$ (racine double) impose la forme $(A+Bt)e^{rx}$, point crucial souvent source d'erreur.
    \item \textbf{Conditions initiales :} Elles traduisent l'état physique du système à l'instant $t=0$ (tension, courant) et permettent de figer les constantes $A, B$. La dérivation du produit $(A+Bt)e^{rx}$ est une compétence technique incontournable.
    \item \textbf{Interprétation :} La solution particulière donne le régime permanent (valeur cible). L'étude de la dérivée ($u'(t) = 1980 t e^{-3t} \ge 0$) prouve la monotonie et exclut tout dépassement, validant le choix du régime critique pour la protection des équipements.
    \item \textbf{Compétences transverses :} Résolution numérique d'équation transcendante (dichotomie/solveur), lecture graphique, rédaction technique argumentée (vocabulaire : régime permanent, transitoire, amorti critique, surtension, temps de réponse).
\end{itemize}
L'élève a ainsi consolidé la chaîne complète : \textbf{Problème concret $\to$ Modèle math $\to$ Résolution algébrique $\to$ Analyse numérique/graphique $\to$ Décision technique}.

\newpage

\coursec{Méthodes et exercices résolus}

\coursec{Équations différentielles}

\courssub{Présentation et vocabulaire}

\begin{definition}[Vocabulaire des équations différentielles]
Une \textbf{équation différentielle} (ED) est une équation qui lie une fonction inconnue $y$ (ou $f$), sa variable $x$ (souvent le temps $t$) et ses dérivées successives.
\begin{itemize}
\item \textbf{Ordre :} ordre de la dérivée la plus élevée qui apparaît (1 pour $y'$, 2 pour $y''$).
\item \textbf{Linéaire :} l'inconnue et ses dérivées apparaissent au premier degré seulement (pas de $y^2$, $\sin y$, $y y'$, etc.).
\item \textbf{À coefficients constants :} les coefficients devant $y$, $y'$, $y''$ sont des nombres réels (pas des fonctions de $x$).
\item \textbf{Homogène :} second membre nul (ex. $y'' + 3y' + 2y = 0$).
\item \textbf{Non homogène :} second membre non nul (ex. $y'' + 3y' + 2y = 6$).
\item \textbf{Solution générale :} famille de toutes les solutions, dépendant de constantes arbitraires ($k$ pour l'ordre 1, $k_1, k_2$ pour l'ordre 2).
\item \textbf{Solution particulière :} une solution précise, obtenue en fixant les constantes grâce à des \textbf{conditions initiales} (valeurs de $y$ et $y'$ en un point).
\end{itemize}
\end{definition}

\begin{propriete}[Théorème de structure (admis)]
L'ensemble des solutions d'une équation différentielle linéaire homogène d'ordre $n$ à coefficients constants est un espace vectoriel de dimension $n$.
L'ensemble des solutions d'une équation non homogène est un espace affine : toute solution s'écrit $y = y_h + y_p$ où $y_h$ est la solution générale de l'équation homogène associée et $y_p$ une solution particulière de l'équation non homogène.
\end{propriete}

\courssub{Équations du type $y' = ay$}

\begin{methode}[Résolution de $y' = ay$ (avec ou sans condition initiale)]
Pour résoudre une équation différentielle du premier ordre $y' = ay$ ($a \in \mathbb{R}$) :
\begin{enumerate}
\item \textbf{Identifier} le coefficient $a$ (attention au signe !). Si l'équation est $y' + 3y = 0$, la ramener à $y' = -3y$ donc $a = -3$.
\item \textbf{Écrire la solution générale} : $y(x) = k\,e^{ax}$, où $k \in \mathbb{R}$ est une constante arbitraire. (Démonstration : séparation des variables ou recherche de $y = e^{rx}$).
\item S'il y a une \textbf{condition initiale} $y(x_0) = y_0$, remplacer $x$ par $x_0$ dans la solution générale pour déterminer $k$ : $k\,e^{a x_0} = y_0 \implies k = y_0\,e^{-a x_0}$.
\item \textbf{Vérifier} en dérivant la solution trouvée : on doit retrouver $y' = ay$.
\item Conclure par une phrase donnant la solution de l'équation (ou de l'équation avec condition initiale).
\end{enumerate}
\end{methode}

\begin{experience}[Découverte : pourquoi $e^{ax}$ ?]
On cherche les solutions de $y' = ay$ sous la forme $y = e^{rx}$ ($r$ réel).
\begin{enumerate}
\item Calculer $y'$ et substituer dans l'équation.
\item Simplifier par $e^{rx}$ (jamais nul). Quelle équation obtient-on pour $r$ ?
\item En déduire la forme de la solution générale.
\end{enumerate}
\tcblower
\textbf{1.} $y = e^{rx} \implies y' = r e^{rx}$. L'équation $y' = ay$ devient $r e^{rx} = a e^{rx}$.
\textbf{2.} On simplifie par $e^{rx} \neq 0$ : $r = a$.
\textbf{3.} La solution est $y = e^{ax}$. Comme l'équation est linéaire homogène, toute solution est un multiple de celle-ci : $y = k e^{ax}$.
\end{experience}

\begin{exoresolu}[Refroidissement d'une pièce mécanique]
Une pièce sortant d'un four est plongée dans un atelier. Sa température $\theta$ (en degrés), fonction du temps $t$ (en minutes), vérifie l'équation différentielle $\theta' = -0{,}2\,\theta$. Au départ, $\theta(0) = 100$.
\begin{enumerate}
\item Donner la solution générale de l'équation différentielle.
\item Déterminer la solution vérifiant la condition initiale.
\item Calculer la température de la pièce au bout de $10$ minutes (on donne $e^{-2} \approx 0{,}135$).
\end{enumerate}
\tcblower
\textbf{1.} L'équation est de la forme $y' = ay$ avec $a = -0{,}2$. La solution générale est :
\[ \theta(t) = k\,e^{-0{,}2\,t}, \quad k \in \mathbb{R}. \]

\textbf{2.} Condition initiale $\theta(0) = 100$ :
\[ \theta(0) = k\,e^{0} = k = 100. \]
Solution particulière : $\theta(t) = 100\,e^{-0{,}2\,t}$.
\emph{Vérification :} $\theta'(t) = 100 \times (-0{,}2)\,e^{-0{,}2\,t} = -0{,}2\,\theta(t)$. \checkmark

\textbf{3.} À $t = 10$ min :
\[ \theta(10) = 100\,e^{-2} \approx 100 \times 0{,}135 = 13{,}5. \]
\textbf{Conclusion :} au bout de $10$ minutes, la température est d'environ $13{,}5^\circ$.
\end{exoresolu}

\begin{popfigure}
\begin{tikzpicture}
\begin{axis}[
    width=\linewidth, height=6cm,
    axis lines=middle, xlabel=$t$ (min), ylabel=$\theta(t)$ ($^\circ$),
    domain=0:15, samples=100,
    xmin=0, xmax=15, ymin=0, ymax=105,
    xtick={0,5,10,15}, ytick={0,25,50,75,100},
    grid=major, grid style={popline},
]
\addplot[popcurve, thick] {100*exp(-0.2*x)};
\addplot[poporange, dashed] coordinates {(0,13.5) (10,13.5) (10,0)};
\node[poporange, anchor=north west] at (axis cs:10,13.5) {$\theta(10)\approx 13,5$};
\end{axis}
\end{tikzpicture}
\legende{Refroidissement exponentiel de la pièce : $\theta(t) = 100 e^{-0,2t}$}
\end{popfigure}

\courssub{Équations $ay'' + by' + cy = 0$ : équation caractéristique}

\begin{methode}[Équation caractéristique et discriminant]
Pour résoudre $ay'' + by' + cy = 0$ ($a \neq 0$, $a,b,c \in \mathbb{R}$) :
\begin{enumerate}
\item \textbf{Écrire l'équation caractéristique} : $ar^2 + br + c = 0$, d'inconnue $r$.
\item \textbf{Calculer le discriminant} $\Delta = b^2 - 4ac$.
\item \textbf{Discuter selon le signe de $\Delta$} pour obtenir la forme de la solution générale (voir section suivante).
\end{enumerate}
\textbf{Formules des racines :} $r = \dfrac{-b \pm \sqrt{\Delta}}{2a}$.
Si $\Delta < 0$ : $r = \alpha \pm i\beta$ avec $\alpha = -\dfrac{b}{2a}$ et $\beta = \dfrac{\sqrt{-\Delta}}{2a}$.
\end{methode}

\begin{experience}[Démonstration guidée : recherche de solutions exponentielles]
On cherche des solutions de $ay''+by'+cy=0$ sous la forme $y = e^{rx}$.
\begin{enumerate}
\item Calculer $y'$ et $y''$, substituer dans l'équation et factoriser par $e^{rx}$.
\item Montrer que $r$ doit vérifier $ar^2+br+c=0$.
\item Expliquer pourquoi, si $r_1 \neq r_2$ sont deux racines réelles, $k_1 e^{r_1 x} + k_2 e^{r_2 x}$ est solution générale.
\end{enumerate}
\tcblower
\textbf{1.} $y' = r e^{rx}$, $y'' = r^2 e^{rx}$. Substitution : $a r^2 e^{rx} + b r e^{rx} + c e^{rx} = (ar^2+br+c)e^{rx}=0$.
Comme $e^{rx} \neq 0$, on a $ar^2+br+c=0$ (équation caractéristique).
\textbf{2.} Si $\Delta > 0$, deux racines réelles distinctes $r_1, r_2$. Les fonctions $e^{r_1 x}$ et $e^{r_2 x}$ sont linéairement indépendantes (leur Wronskien n'est pas nul). L'espace des solutions a dimension 2, donc la solution générale est $y = k_1 e^{r_1 x} + k_2 e^{r_2 x}$.
\end{experience}

\begin{exoresolu}[Trois cas de figure pour le discriminant]
Résoudre dans $\mathbb{R}$ les équations différentielles suivantes :
\begin{enumerate}
\item $y'' - 5y' + 6y = 0$ ;
\item $y'' - 4y' + 4y = 0$ ;
\item $y'' + 2y' + 5y = 0$.
\end{enumerate}
\tcblower
\textbf{1.} Équation caractéristique : $r^2 - 5r + 6 = 0$.
$\Delta = 25 - 24 = 1 > 0$. Racines : $r_1 = 2$, $r_2 = 3$.
Solution générale : $y(x) = k_1 e^{2x} + k_2 e^{3x},\ (k_1,k_2)\in\mathbb{R}^2$.

\textbf{2.} Équation caractéristique : $r^2 - 4r + 4 = 0$.
$\Delta = 16 - 16 = 0$. Racine double : $r_0 = 2$.
Solution générale : $y(x) = (k_1 x + k_2)\,e^{2x},\ (k_1,k_2)\in\mathbb{R}^2$.

\textbf{3.} Équation caractéristique : $r^2 + 2r + 5 = 0$.
$\Delta = 4 - 20 = -16 < 0$. Racines : $r = -1 \pm 2i$ ($\alpha=-1,\ \beta=2$).
Solution générale : $y(x) = e^{-x}\big(k_1 \cos(2x) + k_2 \sin(2x)\big),\ (k_1,k_2)\in\mathbb{R}^2$.
\end{exoresolu}

\courssub{Solutions générales de $ay'' + by' + cy = 0$}

\begin{propriete}[Formes des solutions générales selon $\Delta$]
Soit $ay'' + by' + cy = 0$ ($a \neq 0$) et $\Delta = b^2 - 4ac$.
\begin{itemize}
\item \textbf{$\Delta > 0$ :} deux racines réelles distinctes $r_1, r_2$.
  $y(x) = k_1 e^{r_1 x} + k_2 e^{r_2 x}$.
\item \textbf{$\Delta = 0$ :} une racine réelle double $r_0 = -\dfrac{b}{2a}$.
  $y(x) = (k_1 x + k_2)\,e^{r_0 x}$.
  \emph{Attention :} le terme $x e^{r_0 x}$ est indispensable (indépendance linéaire).
\item \textbf{$\Delta < 0$ :} deux racines complexes conjuguées $\alpha \pm i\beta$ ($\alpha = -\frac{b}{2a},\ \beta = \frac{\sqrt{-\Delta}}{2a}$).
  $y(x) = e^{\alpha x}\big(k_1 \cos(\beta x) + k_2 \sin(\beta x)\big)$.
  \emph{Attention :} c'est $e^{\alpha x}$ (partie réelle) qui multiplie les trigonométriques, pas $e^{\beta x}$.
\end{itemize}
Dans tous les cas, $k_1, k_2 \in \mathbb{R}$ sont des constantes arbitraires.
\end{propriete}

\begin{methode}[Exploiter deux conditions initiales]
Pour trouver LA solution de $ay'' + by' + cy = 0$ vérifiant $y(x_0) = y_0$ et $y'(x_0) = y'_0$ :
\begin{enumerate}
\item Écrire la \textbf{solution générale} $y$ avec $k_1, k_2$.
\item \textbf{Calculer la dérivée} $y'(x)$ (règle du produit si $\Delta \le 0$).
\item Traduire les conditions : système de deux équations à deux inconnues $k_1, k_2$.
\item \textbf{Résoudre le système} (souvent simplifié si $x_0 = 0$ : $e^0=1,\ \cos 0=1,\ \sin 0=0$).
\item \textbf{Conclure} en écrivant la solution particulière, puis vérifier.
\end{enumerate}
\end{methode}

\begin{exoresolu}[Oscillations d'une suspension de véhicule]
La position $y(t)$ (en cm) d'un système de suspension vérifie $y'' + 2y' + 2y = 0$, avec $y(0) = 1$ et $y'(0) = 0$.
\begin{enumerate}
\item Donner la solution générale.
\item Déterminer la solution vérifiant les conditions initiales.
\end{enumerate}
\tcblower
\textbf{1.} Équation caractéristique : $r^2 + 2r + 2 = 0$.
$\Delta = 4 - 8 = -4 < 0$. Racines : $r = -1 \pm i$ ($\alpha=-1,\ \beta=1$).
Solution générale : $y(t) = e^{-t}\big(k_1 \cos t + k_2 \sin t\big),\ (k_1,k_2)\in\mathbb{R}^2$.

\textbf{2.} Condition $y(0)=1$ : $y(0) = e^0(k_1 \cdot 1 + k_2 \cdot 0) = k_1 = 1$.

Dérivée (produit $u=v=e^{-t}$, $v=k_1\cos t+k_2\sin t$) :
$y'(t) = -e^{-t}(k_1\cos t+k_2\sin t) + e^{-t}(-k_1\sin t+k_2\cos t)$
$= e^{-t}\big((k_2-k_1)\cos t - (k_1+k_2)\sin t\big)$.

Condition $y'(0)=0$ : $y'(0) = (k_2 - k_1) = 0 \implies k_2 = k_1 = 1$.

\textbf{Solution :} $y(t) = e^{-t}(\cos t + \sin t)$.
L'amortissement $e^{-t}$ traduit l'atténuation progressive des oscillations.
\end{exoresolu}

\begin{popfigure}
\begin{tikzpicture}
\begin{axis}[
    width=\linewidth, height=6cm,
    axis lines=middle, xlabel=$t$ (s), ylabel=$y(t)$ (cm),
    domain=0:10, samples=200,
    xmin=0, xmax=10, ymin=-1.5, ymax=1.5,
    xtick={0,3.14,6.28,9.42}, xticklabels={$0$,$\pi$,$2\pi$,$3\pi$},
    ytick={-1,-0.5,0.5,1},
    grid=major, grid style={popline},
]
\addplot[popcurve, thick, domain=0:10] {exp(-x)*(cos(deg(x))+sin(deg(x)))};
\addplot[poporange, dashed, domain=0:10] {exp(-x)*1.414};
\addplot[poporange, dashed, domain=0:10] {-exp(-x)*1.414};
\node[poporange, anchor=south east] at (axis cs:2,1.2) {Enveloppe $\pm\sqrt{2}e^{-t}$};
\end{axis}
\end{tikzpicture}
\legende{Oscillations amorties de la suspension : $y(t)=e^{-t}(\cos t+\sin t)$}
\end{popfigure}

\courssub{Équations $ay'' + by' + cy = d$ (second membre constant)}

\begin{methode}[Résolution : homogène + particulière constante]
Pour résoudre $ay'' + by' + cy = d$ ($d \in \mathbb{R}$ constant, $a \neq 0$) :
\begin{enumerate}
\item \textbf{Résoudre l'équation homogène} $ay'' + by' + cy = 0$ : obtenir $y_h$ (avec $k_1, k_2$).
\item \textbf{Chercher une solution particulière constante} $y_p = \lambda$ : $y_p'=0,\ y_p''=0 \implies c\lambda = d$.
  \begin{itemize}
  \item Si $c \neq 0$ : $\lambda = \dfrac{d}{c}$, donc $y_p = \dfrac{d}{c}$.
  \item Si $c = 0$ (équation $ay''+by'=d$) : chercher $y_p = \lambda x$ (car $\lambda$ constant donnerait $0=d$ impossible).
  \end{itemize}
\item \textbf{Additionner} : solution générale $y = y_h + y_p$.
\item S'il y a des conditions initiales, déterminer $k_1, k_2$ \textbf{sur $y_h + y_p$} (jamais sur $y_h$ seule).
\end{enumerate}
\end{methode}

\begin{exoresolu}[Charge d'un condensateur en série RC]
La tension $u(t)$ (en volts) aux bornes d'un condensateur en charge vérifie $u'' + 3u' + 2u = 6$, avec $u(0) = 0$ et $u'(0) = 0$.
\begin{enumerate}
\item Résoudre l'équation sans second membre.
\item Déterminer une solution particulière constante, puis la solution générale complète.
\item Déterminer la solution vérifiant les conditions initiales.
\end{enumerate}
\tcblower
\textbf{1.} Homogène : $u'' + 3u' + 2u = 0$. Caractéristique : $r^2+3r+2=0$.
$\Delta = 1 > 0$. Racines : $r_1=-2,\ r_2=-1$.
$u_h(t) = k_1 e^{-2t} + k_2 e^{-t}$.

\textbf{2.} Particulière constante $u_p = \lambda$ : $2\lambda = 6 \implies \lambda = 3$.
Solution générale : $u(t) = k_1 e^{-2t} + k_2 e^{-t} + 3$.

\textbf{3.} Conditions sur la solution \textbf{complète} :
$u(0) = k_1 + k_2 + 3 = 0 \implies k_1 + k_2 = -3$.
$u'(t) = -2k_1 e^{-2t} - k_2 e^{-t} \implies u'(0) = -2k_1 - k_2 = 0$.
Système : $\begin{cases} k_1+k_2=-3 \\ -2k_1-k_2=0 \end{cases} \implies -k_1 = -3 \implies k_1=3,\ k_2=-6$.

\textbf{Solution :} $u(t) = 3e^{-2t} - 6e^{-t} + 3$.
\emph{Interprétation :} $\lim_{t\to+\infty} u(t) = 3$ V. La tension se stabilise à la valeur imposée par le générateur (régime permanent).
\end{exoresolu}

\begin{popfigure}
\begin{tikzpicture}
\begin{axis}[
    width=\linewidth, height=6cm,
    axis lines=middle, xlabel=$t$ (s), ylabel=$u(t)$ (V),
    domain=0:5, samples=100,
    xmin=0, xmax=5, ymin=-0.5, ymax=3.5,
    xtick={0,1,2,3,4,5}, ytick={0,1,2,3},
    grid=major, grid style={popline},
]
\addplot[popcurve, thick] {3*exp(-2*x) - 6*exp(-x) + 3};
\addplot[poporange, dashed] coordinates {(0,3) (5,3)};
\node[poporange, anchor=south west] at (axis cs:0.5,3) {Régime permanent $u=3$ V};
\end{axis}
\end{tikzpicture}
\legende{Charge du condensateur : $u(t) = 3e^{-2t} - 6e^{-t} + 3$}
\end{popfigure}

\courssub{Applications à des situations concrètes et synthèse}

\begin{methode}[Modéliser et résoudre un problème concret]
\begin{enumerate}
\item \textbf{Nommer} la grandeur inconnue $y$ (fonction) et la variable $t$ (temps), avec unités.
\item \textbf{Traduire} l'énoncé en équation différentielle :
  \begin{itemize}
  \item « Taux de variation proportionnel à la quantité » $\to y' = ky$ (ordre 1).
  \item « Accélération / forces / loi de Kirchhoff » $\to$ ordre 2 ($y''$ apparaît).
  \end{itemize}
\item \textbf{Identifier le type} : $y'=ay$, ou $ay''+by'+cy=0$, ou $ay''+by'+cy=d$.
\item \textbf{Résoudre} avec la méthode adaptée + conditions initiales du problème.
\item \textbf{Interpréter} : répondre à la question (valeur, instant, limite, seuil) avec phrase et unité.
\end{enumerate}
\end{methode}

\begin{exoresolu}[Désintégration d'un produit chimique (atelier traitement de surface)]
La quantité $q(t)$ (en grammes) d'un produit se décompose selon $q' = -0{,}1\,q$. Initialement, $q(0) = 500$ g.
\begin{enumerate}
\item Déterminer $q(t)$.
\item Quantité restante à $t = 10$ h ($e^{-1} \approx 0{,}368$).
\item Le produit est inefficace sous $200$ g. Est-il utilisable à $t = 8$ h ($e^{-0{,}8} \approx 0{,}449$) ?
\end{enumerate}
\tcblower
\textbf{1.} $q' = aq$ avec $a=-0{,}1$. Solution générale $q(t)=k e^{-0{,}1t}$.
$q(0)=k=500 \implies q(t) = 500 e^{-0{,}1t}$.

\textbf{2.} $q(10) = 500 e^{-1} \approx 500 \times 0{,}368 = 184$ g.

\textbf{3.} $q(8) = 500 e^{-0{,}8} \approx 500 \times 0{,}449 = 224{,}5$ g.
$224{,}5 > 200$ : le produit est \textbf{encore utilisable} à $8$ h.
\textbf{Conclusion :} le seuil $200$ g est franchi entre $8$ h et $10$ h.
\end{exoresolu}

\begin{exoresolu}[Sujet type Probatoire / Baccalauréat Technique]
On considère l'équation différentielle $(E)$ : $y'' - 2y' + y = 0$.
\begin{enumerate}
\item Résoudre $(E)$.
\item Déterminer la solution $f$ de $(E)$ vérifiant $f(0) = 2$ et $f'(0) = 5$.
\end{enumerate}
\tcblower
\textbf{1.} Équation caractéristique : $r^2 - 2r + 1 = 0$.
$\Delta = 4 - 4 = 0$. Racine double $r_0 = 1$.
Solution générale de $(E)$ : $y(x) = (k_1 x + k_2)\,e^{x},\ (k_1,k_2)\in\mathbb{R}^2$.

\textbf{2.} $f(0) = (0 + k_2)e^0 = k_2 = 2 \implies k_2 = 2$.
Dérivée : $f'(x) = k_1 e^x + (k_1 x + k_2)e^x = (k_1 x + k_1 + k_2)e^x$.
$f'(0) = (k_1 + k_2) = 5 \implies k_1 = 5 - 2 = 3$.

\textbf{Solution :} $f(x) = (3x + 2)e^x$.
\emph{Vérification :} $f(0)=2$, $f'(0)=3+2=5$. \checkmark
\end{exoresolu}

\begin{savaistu}[Histoire et applications au Cameroun]
\begin{itemize}
\item Les équations différentielles naissent avec Newton (1671) pour décrire le mouvement ($F=ma \implies m y'' = F$).
\item \textbf{Électricité :} circuits RLC (alternateurs, transformateurs des barrages d'Édéa, Song Loulou).
\item \textbf{Mécanique :} amortisseurs de véhicules (routes camerounaises), vibrations de machines-outils.
\item \textbf{Gestion / Biologie :} modèle de Malthus ($P' = kP$) pour population ou bactéries ; modèle de Verhulst (logistique) pour ressources limitées.
\item \textbf{Chimie :} cinétique d'ordre 1 (désintégration, radioactivité, dégradation polluants).
\end{itemize}
Au Probatoire, on attend une rédaction soignée : \textbf{énoncé de la méthode}, \textbf{calculs détaillés}, \textbf{constantes explicitées}, \textbf{phrase de conclusion}.
\end{savaistu}

\begin{attention}[Les 5 erreurs qui coûtent des points au Probatoire]
\begin{enumerate}
\item \textbf{Forme canonique oubliée.} $y' + 3y = 0 \implies a = -3$ (pas $+3$). Solution $y = k e^{-3x}$.
\item \textbf{Confusion des 3 cas selon $\Delta$.}
  $\Delta=0 \to (k_1 x + k_2)e^{r_0 x}$ (ne pas oublier le $x$).
  $\Delta<0 \to e^{\alpha x}(k_1\cos\beta x + k_2\sin\beta x)$ (exposant $\alpha$, pas $\beta$).
\item \textbf{Conditions initiales sur $y_h$ au lieu de $y_h+y_p$.} Pour $ay''+by'+cy=d$, les constantes se déterminent \textbf{après} avoir ajouté $y_p = d/c$.
\item \textbf{Dérivation de produit bâclée.} Pour $y=(k_1 x+k_2)e^{r_0 x}$ ou $y=e^{\alpha x}(k_1\cos\beta x+k_2\sin\beta x)$, appliquer $(uv)'=u'v+uv'$ soigneusement.
\item \textbf{Conclusion absente ou incomplète.} Écrire « $y = k_1 e^{2x}+k_2 e^{3x}$ » sans « $(k_1,k_2)\in\mathbb{R}^2$ », ou ne pas répondre à la question (valeur numérique, unité, interprétation physique).
\end{enumerate}
\end{attention}

\newpage

\coursec{Exercices}

\courssub{Je vérifie mes connaissances}

\begin{exercice}[QCM : à chaud]
Pour chaque question, une seule réponse est exacte. Entoure-la et justifie brièvement.
\begin{enumerate}
\item La solution générale de l'équation $y' = 2y$ est :
\begin{enumerate}
\item $y = 2x + C$ \quad \item $y = Ce^{2x}$ \quad \item $y = Ce^{-2x}$ \quad \item $y = 2e^{x}$
\end{enumerate}
\item L'équation caractéristique de $y'' - 5y' + 6y = 0$ est :
\begin{enumerate}
\item $r^2 - 5r + 6 = 0$ \quad \item $r - 5r + 6 = 0$ \quad \item $r^2 + 5r + 6 = 0$ \quad \item $r^2 - 5r - 6 = 0$
\end{enumerate}
\item Si l'équation caractéristique de $ay'' + by' + cy = 0$ admet deux racines réelles distinctes $r_1$ et $r_2$, alors la solution générale est :
\begin{enumerate}
\item $y = (A + Bx)e^{r_1 x}$ \quad \item $y = e^{\alpha x}(A\cos\beta x + B\sin\beta x)$ \quad \item $y = Ae^{r_1 x} + Be^{r_2 x}$
\end{enumerate}
\item Une solution particulière constante de $y'' + 3y' + 2y = 6$ est :
\begin{enumerate}
\item $y = 6$ \quad \item $y = 3$ \quad \item $y = 2$ \quad \item $y = 1$
\end{enumerate}
\end{enumerate}
\end{exercice}

\begin{exercice}[Vrai ou faux, justifié]
Réponds par VRAI ou FAUX en justifiant chaque réponse par un argument ou un calcul.
\begin{enumerate}
\item Toute solution de l'équation $y'' + y = 0$ est une fonction périodique.
\item La fonction $t \longmapsto e^{2t}$ est solution de l'équation $y'' - 4y = 0$.
\item L'équation $y' = -3y$ admet une unique solution $f$ telle que $f(0) = 5$.
\item Si l'équation caractéristique de $ay'' + by' + cy = 0$ a un discriminant strictement négatif, alors ses solutions s'écrivent sans fonctions trigonométriques.
\end{enumerate}
\end{exercice}

\begin{exercice}[Vocabulaire et reconnaissance]
Pour chacune des équations différentielles suivantes, précise son ordre, si elle est linéaire à coefficients constants, et si elle est avec ou sans second membre :
\begin{enumerate}
\item $y' = -0{,}5y$ ;
\item $2y'' + 3y' - y = 0$ ;
\item $y'' + 4y = 8$ ;
\item $y'' + (y')^2 = 0$.
\end{enumerate}
Donne ensuite la définition complète d'une \cle{solution} d'une équation différentielle sur un intervalle $I$.
\end{exercice}

\begin{exercice}[Calculs directs]
\begin{enumerate}
\item Résous les équations caractéristiques suivantes et indique, sans autre calcul, la forme de la solution générale de l'équation différentielle associée :
\[ r^2 - 4r + 3 = 0 \quad ; \quad r^2 + 2r + 1 = 0 \quad ; \quad r^2 + 4 = 0. \]
\item Vérifie que la fonction $f : x \longmapsto 3e^{-2x}$ est solution de $y' = -2y$ et qu'elle vérifie $f(0) = 3$.
\item Détermine une solution particulière constante de l'équation $y'' - y' + 2y = 10$.
\end{enumerate}
\end{exercice}

\courssub{Je m'entraîne}

\begin{exercice}[Premier ordre : décroissance]
On considère l'équation différentielle $y' = -3y$.
\begin{enumerate}
\item Donne la solution générale de cette équation.
\item Détermine la solution $f$ vérifiant la condition initiale $f(0) = 5$.
\item Détermine l'instant $t_0$ (arrondi à $10^{-3}$) où la solution vaut la moitié de sa valeur initiale, c'est-à-dire $f(t_0) = \dfrac{5}{2}$.
\end{enumerate}
\end{exercice}

\begin{exercice}[Second ordre : deux racines distinctes]
On considère l'équation différentielle $(E) : y'' - 4y' + 3y = 0$.
\begin{enumerate}
\item Résous l'équation caractéristique de $(E)$.
\item Donne la solution générale de $(E)$.
\item Détermine la solution $f$ de $(E)$ vérifiant $f(0) = 2$ et $f'(0) = 4$.
\item Étudie la limite de $f(x)$ lorsque $x$ tend vers $+\infty$.
\end{enumerate}
\end{exercice}

\begin{exercice}[Second ordre : racine double]
On considère l'équation différentielle $(E) : y'' + 2y' + y = 0$.
\begin{enumerate}
\item Montre que l'équation caractéristique admet une racine double que l'on précisera.
\item Donne la solution générale de $(E)$.
\item Détermine la solution $f$ vérifiant $f(0) = 1$ et $f'(0) = -1$.
\item Montre que la courbe de $f$ coupe l'axe des abscisses en un unique point, dont on déterminera l'abscisse.
\end{enumerate}
\end{exercice}

\begin{exercice}[Avec second membre constant]
On considère l'équation différentielle $(E) : y'' + 4y = 8$.
\begin{enumerate}
\item Résous l'équation sans second membre $y'' + 4y = 0$.
\item Détermine une solution particulière constante de $(E)$.
\item En déduire la solution générale de $(E)$.
\item Détermine la solution $f$ de $(E)$ vérifiant $f(0) = 0$ et $f'(0) = 0$.
\item Montre que pour tout réel $t$, $f(t) = 2 - 2\cos(2t)$, puis détermine les valeurs minimale et maximale prises par $f(t)$.
\end{enumerate}
\end{exercice}

\courssub{Je me prépare au Probatoire}

\begin{exercice}[Décharge d'un condensateur — circuit RC]
Un technicien en électricité étudie la décharge d'un condensateur de capacité $C$ dans un circuit comportant une résistance $R$. On note $q(t)$ la charge électrique (en coulombs) du condensateur à l'instant $t$ (en secondes). La physique montre que $q$ est solution de l'équation différentielle
\[ q' = -\dfrac{1}{RC}\,q. \]
On donne $R = \num{2000}$ ohms, $C = 5 \times 10^{-4}$ farad, et la charge initiale $q(0) = q_0 = 8 \times 10^{-3}$ coulomb.
\begin{enumerate}
\item Calcule le produit $RC$ (appelé \cle{constante de temps} du circuit, notée $\tau$).
\item Résous l'équation différentielle et donne l'expression de $q(t)$ en fonction de $t$ et de $q_0$.
\item Calcule la charge du condensateur à l'instant $t = 1$ seconde (arrondir à $10^{-4}$).
\item On dit que le condensateur est déchargé à $90\,\%$ lorsque $q(t) = 0{,}1\,q_0$. Détermine, par le calcul, l'instant $t_1$ correspondant (on donnera la valeur exacte puis une valeur arrondie au centième de seconde).
\item Exprime $t_1$ en fonction de $\tau$ et de $\ln 10$. Le technicien affirme : « il suffit d'attendre environ $2{,}3\tau$ pour que le condensateur soit déchargé à $90\,\%$ ». Cette affirmation est-elle exacte ? Justifie.
\end{enumerate}
\end{exercice}

\begin{exercice}[Suspension d'un véhicule — ressort amorti]
Lors d'un essai en atelier, on étudie la suspension d'un véhicule. On note $x(t)$ le déplacement (en centimètres) de la caisse par rapport à sa position d'équilibre à l'instant $t$ (en secondes). On admet que $x$ est solution de l'équation différentielle
\[ (E) : x'' + 4x' + 4x = 0, \]
avec les conditions initiales $x(0) = 10$ (la caisse est écartée de \SI{10}{cm}) et $x'(0) = 0$ (on la lâche sans vitesse initiale).
\begin{enumerate}
\item Résous l'équation caractéristique de $(E)$. Que remarque-t-on ?
\item Donne la solution générale de $(E)$, puis montre que pour tout $t \geq 0$ :
\[ x(t) = 10(1 + 2t)e^{-2t}. \]
\item Calcule la dérivée $x'(t)$ et montre que $x'(t) = -40\,t\,e^{-2t}$.
\item Dresse le tableau de variations de $x$ sur $[0\,;\,+\infty[$ (on admettra que $\displaystyle\lim_{t \to +\infty} x(t) = 0$).
\item Le mécanicien hésite entre deux réglages de l'amortisseur :
\begin{itemize}
\item réglage A : l'équation serait $x'' + 2x' + 4x = 0$ ;
\item réglage B : l'équation $(E)$ ci-dessus.
\end{itemize}
Pour le réglage A, montre que l'équation caractéristique a un discriminant strictement négatif et écris la forme de la solution générale (sans déterminer les constantes).
\item En comparant les formes des solutions des deux réglages, explique pourquoi le réglage A provoque des oscillations de la caisse alors que le réglage B n'en provoque pas. Rédige une conclusion justifiant le choix du réglage B pour le confort des passagers.
\end{enumerate}
\end{exercice}

\begin{exercice}[Oscillations d'un vibrateur industriel]
Une machine d'atelier est équipée d'un vibrateur dont la position $y(t)$ (en millimètres) à l'instant $t$ (en secondes) vérifie l'équation différentielle
\[ (E) : y'' + 4y = 8, \]
avec les conditions initiales $y(0) = 0$ et $y'(0) = 0$.
\begin{enumerate}
\item Résous l'équation sans second membre $(E_0) : y'' + 4y = 0$.
\item Vérifie que la fonction constante $y_p = 2$ est une solution particulière de $(E)$.
\item En déduire la solution générale de $(E)$, puis montrer que l'unique solution vérifiant les conditions initiales est
\[ y(t) = 2 - 2\cos(2t). \]
\item Montre que $y$ est périodique et précise sa période $T$.
\item Détermine les instants $t \in [0\,;\,2\pi]$ pour lesquels le vibrateur atteint sa position maximale, et précise cette position maximale.
\item Trace la courbe représentative de $y$ sur l'intervalle $[0\,;\,2\pi]$ dans un repère orthogonal (unités : \SI{2}{cm} par seconde en abscisse, \SI{1}{cm} par millimètre en ordonnée).
\item Interprète le mouvement du vibrateur : autour de quelle position oscille-t-il ? Quelle est l'amplitude des oscillations ? Rédige ta réponse en quelques phrases.
\end{enumerate}
\end{exercice}

\newpage

\coursec{Corrigés des exercices}

\begin{corrige}[Exercice 1 — QCM : à chaud]
\begin{enumerate}
\item \textbf{Réponse b :} $y = Ce^{2x}$. En effet, si $y = Ce^{2x}$, alors $y' = 2Ce^{2x} = 2y$. La solution générale de $y' = ay$ est $y = Ce^{ax}$, ici avec $a = 2$.
\item \textbf{Réponse a :} $r^2 - 5r + 6 = 0$. L'équation caractéristique de $ay'' + by' + cy = 0$ est $ar^2 + br + c = 0$ ; ici $a = 1$, $b = -5$, $c = 6$.
\item \textbf{Réponse c :} $y = Ae^{r_1 x} + Be^{r_2 x}$. La forme $(A + Bx)e^{r_1 x}$ correspond à une racine double, et $e^{\alpha x}(A\cos\beta x + B\sin\beta x)$ à des racines complexes conjuguées.
\item \textbf{Réponse b :} $y = 3$. On cherche une constante $k$ : alors $y' = y'' = 0$, donc l'équation devient $2k = 6$, soit $k = 3$.
\end{enumerate}
\end{corrige}

\begin{corrige}[Exercice 2 — Vrai ou faux, justifié]
\begin{enumerate}
\item \textbf{VRAI.} L'équation caractéristique $r^2 + 1 = 0$ a pour racines $r = \pm i$ (donc $\alpha = 0$, $\beta = 1$). La solution générale est $y = A\cos x + B\sin x$, somme de fonctions périodiques de période $2\pi$ : toute solution est périodique (de période $2\pi$).
\item \textbf{VRAI.} Soit $f(t) = e^{2t}$. Alors $f'(t) = 2e^{2t}$ et $f''(t) = 4e^{2t}$. Donc $f''(t) - 4f(t) = 4e^{2t} - 4e^{2t} = 0$ : $f$ est bien solution de $y'' - 4y = 0$.
\item \textbf{VRAI.} La solution générale de $y' = -3y$ est $f(x) = Ce^{-3x}$. La condition $f(0) = 5$ donne $C = 5$, d'où l'unique solution $f(x) = 5e^{-3x}$.
\item \textbf{FAUX.} C'est le contraire : si le discriminant est strictement négatif, les racines sont complexes conjuguées $\alpha \pm i\beta$ et les solutions s'écrivent $y = e^{\alpha x}(A\cos\beta x + B\sin\beta x)$, qui font intervenir des fonctions trigonométriques.
\end{enumerate}
\end{corrige}

\begin{corrige}[Exercice 3 — Vocabulaire et reconnaissance]
\begin{enumerate}
\item $y' = -0{,}5y$ : équation d'\cle{ordre 1}, linéaire à coefficients constants, sans second membre.
\item $2y'' + 3y' - y = 0$ : équation d'\cle{ordre 2}, linéaire à coefficients constants, sans second membre.
\item $y'' + 4y = 8$ : équation d'\cle{ordre 2}, linéaire à coefficients constants, avec second membre (constant, égal à $8$).
\item $y'' + (y')^2 = 0$ : équation d'\cle{ordre 2}, \textbf{non linéaire} (à cause du terme $(y')^2$), sans second membre.
\end{enumerate}
\textbf{Définition :} on appelle \cle{solution} d'une équation différentielle sur un intervalle $I$ toute fonction $f$ dérivable (autant de fois que nécessaire) sur $I$, telle que pour tout $x$ de $I$, l'égalité définissant l'équation soit vérifiée lorsqu'on remplace $y$, $y'$, $y''$ par $f(x)$, $f'(x)$, $f''(x)$.
\end{corrige}

\begin{corrige}[Exercice 4 — Calculs directs]
\begin{enumerate}
\item \textbf{Équation $r^2 - 4r + 3 = 0$ :} $\Delta = 16 - 12 = 4 > 0$, racines $r_1 = \dfrac{4 - 2}{2} = 1$ et $r_2 = \dfrac{4 + 2}{2} = 3$. Solution générale : $y = Ae^{x} + Be^{3x}$.

\textbf{Équation $r^2 + 2r + 1 = 0$ :} $\Delta = 4 - 4 = 0$, racine double $r_0 = -1$. Solution générale : $y = (A + Bx)e^{-x}$.

\textbf{Équation $r^2 + 4 = 0$ :} $\Delta = -16 < 0$, racines $r = \pm 2i$ ($\alpha = 0$, $\beta = 2$). Solution générale : $y = A\cos(2x) + B\sin(2x)$.
\item Soit $f(x) = 3e^{-2x}$. Alors $f'(x) = 3 \times (-2)e^{-2x} = -6e^{-2x} = -2 \times 3e^{-2x} = -2f(x)$ : $f$ est bien solution de $y' = -2y$. De plus $f(0) = 3e^{0} = 3$.
\item On cherche une solution constante $y = k$ : alors $y' = y'' = 0$ et l'équation devient $2k = 10$, soit $k = 5$. Une solution particulière est $y_p = 5$.
\end{enumerate}
\end{corrige}

\begin{corrige}[Exercice 5 — Premier ordre : décroissance]
\begin{enumerate}
\item L'équation $y' = -3y$ est du type $y' = ay$ avec $a = -3$. Sa solution générale est
\[ f(t) = Ce^{-3t}, \quad C \in \mathbb{R}. \]
\item $f(0) = Ce^{0} = C = 5$. Donc $f(t) = 5e^{-3t}$.
\item On résout $f(t_0) = \dfrac{5}{2}$ :
\begin{align*}
5e^{-3t_0} &= \dfrac{5}{2} \\
e^{-3t_0} &= \dfrac{1}{2} \\
-3t_0 &= \ln\dfrac{1}{2} = -\ln 2 \\
t_0 &= \dfrac{\ln 2}{3} \approx 0{,}231.
\end{align*}
La moitié de la valeur initiale est atteinte à $t_0 \approx 0{,}231$ (environ $\num{0,231}$ unité de temps).
\end{enumerate}
\end{corrige}

\begin{corrige}[Exercice 6 — Second ordre : deux racines distinctes]
\begin{enumerate}
\item Équation caractéristique : $r^2 - 4r + 3 = 0$. $\Delta = 16 - 12 = 4 > 0$ ; $r_1 = \dfrac{4 - 2}{2} = 1$, $r_2 = \dfrac{4 + 2}{2} = 3$.
\item Solution générale : $f(x) = Ae^{x} + Be^{3x}$, avec $A, B \in \mathbb{R}$.
\item $f(0) = A + B = 2$. De plus $f'(x) = Ae^{x} + 3Be^{3x}$, donc $f'(0) = A + 3B = 4$. On résout le système :
\[ \begin{cases} A + B = 2 \\ A + 3B = 4 \end{cases} \]
Par soustraction : $2B = 2$, donc $B = 1$ puis $A = 1$. La solution cherchée est
\[ f(x) = e^{x} + e^{3x}. \]
\item Quand $x \to +\infty$, $e^{x} \to +\infty$ et $e^{3x} \to +\infty$, donc
\[ \lim_{x \to +\infty} f(x) = +\infty. \]
\end{enumerate}
\end{corrige}

\begin{corrige}[Exercice 7 — Second ordre : racine double]
\begin{enumerate}
\item Équation caractéristique : $r^2 + 2r + 1 = 0$, c'est-à-dire $(r + 1)^2 = 0$. Le discriminant est $\Delta = 4 - 4 = 0$ : l'équation admet une \cle{racine double} $r_0 = -1$.
\item Solution générale : $f(x) = (A + Bx)e^{-x}$, avec $A, B \in \mathbb{R}$.
\item $f(0) = A = 1$. Dérivons : $f'(x) = Be^{-x} - (A + Bx)e^{-x} = (B - A - Bx)e^{-x}$. Alors $f'(0) = B - A = -1$, donc $B = -1 + A = 0$. La solution est
\[ f(x) = e^{-x}. \]
\item $f(x) = e^{-x} > 0$ pour tout réel $x$ : la courbe de $f$ \textbf{ne coupe jamais} l'axe des abscisses. L'affirmation de l'énoncé serait vraie pour d'autres conditions initiales, mais avec $f(x) = e^{-x}$, il n'existe \textbf{aucun} point d'intersection avec l'axe des abscisses, puisqu'une exponentielle est strictement positive. (Point de vigilance : toujours vérifier le signe de la solution obtenue avant de chercher une intersection.)
\end{enumerate}
\end{corrige}

\begin{corrige}[Exercice 8 — Avec second membre constant]
\begin{enumerate}
\item Équation sans second membre : $y'' + 4y = 0$. Équation caractéristique : $r^2 + 4 = 0$, soit $r = \pm 2i$ ($\alpha = 0$, $\beta = 2$). Solution générale :
\[ y_0 = A\cos(2x) + B\sin(2x). \]
\item On cherche $y_p = k$ constante : $y_p' = y_p'' = 0$, donc $4k = 8$, soit $k = 2$. Une solution particulière est $y_p = 2$.
\item La solution générale de $(E)$ est la somme :
\[ y = A\cos(2x) + B\sin(2x) + 2. \]
\item $f(0) = A + 2 = 0$ donne $A = -2$. Dérivée : $f'(x) = -2A\sin(2x) + 2B\cos(2x)$, donc $f'(0) = 2B = 0$, soit $B = 0$. Donc
\[ f(t) = 2 - 2\cos(2t). \]
\item D'après la question précédente, $f(t) = 2 - 2\cos(2t)$ pour tout réel $t$. Comme $-1 \leq \cos(2t) \leq 1$, on a $-2 \leq -2\cos(2t) \leq 2$, donc
\[ 0 \leq f(t) \leq 4. \]
La valeur \textbf{minimale} est $0$ (atteinte quand $\cos(2t) = 1$, par exemple $t = 0$) et la valeur \textbf{maximale} est $4$ (atteinte quand $\cos(2t) = -1$, par exemple $t = \dfrac{\pi}{2}$).
\end{enumerate}
\end{corrige}

\begin{corrige}[Exercice 9 — Décharge d'un condensateur — circuit RC]
\begin{enumerate}
\item $\tau = RC = \num{2000} \times \num{5e-4} = 1$ seconde.
\item L'équation $q' = -\dfrac{1}{RC}q$ est du type $y' = ay$ avec $a = -\dfrac{1}{RC} = -1$. Solution générale : $q(t) = Ke^{-t/RC}$. La condition $q(0) = q_0$ donne $K = q_0$, d'où
\[ q(t) = q_0\,e^{-t/RC} = 8 \times 10^{-3}\,e^{-t}. \]
\item $q(1) = 8 \times 10^{-3}\,e^{-1} = \dfrac{8 \times 10^{-3}}{e} \approx 2{,}943 \times 10^{-3} \approx 2{,}9 \times 10^{-3}$ coulomb (arrondi à $10^{-4}$ : $q(1) \approx \num{0,0029}$ C).
\item On résout $q(t_1) = 0{,}1\,q_0$ :
\begin{align*}
q_0\,e^{-t_1} &= 0{,}1\,q_0 \\
e^{-t_1} &= 0{,}1 \\
-t_1 &= \ln(0{,}1) = -\ln 10 \\
t_1 &= \ln 10 \approx 2{,}30 \text{ s}.
\end{align*}
\item Comme $\tau = RC = 1$, on a $t_1 = \tau\ln 10$. Or $\ln 10 \approx 2{,}303$, donc $t_1 \approx 2{,}3\tau$. \textbf{L'affirmation du technicien est exacte} : la décharge à $90\,\%$ est atteinte au bout d'environ $2{,}3$ fois la constante de temps.
\end{enumerate}
\textbf{Barème indicatif (10 points) :} calcul de $\tau$ : 1 pt ; solution générale : 2 pts ; prise en compte de $q(0)$ : 1 pt ; $q(1)$ : 2 pts ; résolution de l'équation avec logarithme : 2 pts ; conclusion argumentée : 2 pts.

\textbf{Points de vigilance :} citer la forme $y' = ay \Rightarrow y = Ce^{ax}$ ; simplifier par $q_0 \neq 0$ avant de prendre le logarithme ; ne pas oublier que $\ln(0{,}1) = -\ln 10$ ; arrondir correctement ($\ln 10 \approx 2{,}3026$, donc $2{,}30$ au centième).
\end{corrige}

\begin{corrige}[Exercice 10 — Suspension d'un véhicule — ressort amorti]
\begin{enumerate}
\item Équation caractéristique : $r^2 + 4r + 4 = 0$, soit $(r + 2)^2 = 0$. $\Delta = 16 - 16 = 0$ : on remarque qu'il y a une \cle{racine double} $r_0 = -2$ (régime dit \emph{critique}).
\item Solution générale : $x(t) = (A + Bt)e^{-2t}$. Conditions initiales : $x(0) = A = 10$. Dérivée : $x'(t) = Be^{-2t} - 2(A + Bt)e^{-2t}$, donc $x'(0) = B - 2A = 0$, soit $B = 2A = 20$. Donc
\[ x(t) = (10 + 20t)e^{-2t} = 10(1 + 2t)e^{-2t}. \]
\item Dérivons $x(t) = 10(1 + 2t)e^{-2t}$ (produit $uv$ avec $u = 1 + 2t$, $u' = 2$, $v = e^{-2t}$, $v' = -2e^{-2t}$) :
\begin{align*}
x'(t) &= 10\left[2e^{-2t} - 2(1 + 2t)e^{-2t}\right] \\
&= 10e^{-2t}\left[2 - 2 - 4t\right] \\
&= -40\,t\,e^{-2t}.
\end{align*}
\item Pour $t \geq 0$ : $e^{-2t} > 0$ et $-40t \leq 0$, donc $x'(t) \leq 0$ (nul seulement en $t = 0$). La fonction $x$ est donc \textbf{décroissante} sur $[0\,;\,+\infty[$, de $x(0) = 10$ vers $0$ (limite admise).
\begin{center}
\begin{tabular}{|c|ccc|}\hline
$t$ & $0$ & & $+\infty$ \\ \hline
$x'(t)$ & $0$ & $-$ & \\ \hline
$x$ & $10$ & $\searrow$ & $0$ \\ \hline
\end{tabular}
\end{center}
\item Réglage A : équation caractéristique $r^2 + 2r + 4 = 0$. Discriminant : $\Delta = 4 - 16 = -12 < 0$. Racines complexes : $r = \dfrac{-2 \pm i\sqrt{12}}{2} = -1 \pm i\sqrt{3}$ ($\alpha = -1$, $\beta = \sqrt{3}$). Solution générale :
\[ x(t) = e^{-t}\left(A\cos(\sqrt{3}\,t) + B\sin(\sqrt{3}\,t)\right). \]
\item Avec le réglage A, la solution contient des fonctions $\cos$ et $\sin$ : la caisse \textbf{oscille} autour de sa position d'équilibre (avec une amplitude qui décroît à cause du facteur $e^{-t}$). Avec le réglage B, la solution $10(1 + 2t)e^{-2t}$ est positive et strictement décroissante : la caisse \textbf{revient directement} vers sa position d'équilibre sans osciller.

\textbf{Conclusion :} pour le confort des passagers, on choisit le \textbf{réglage B} : il évite les oscillations (à-coups) et ramène la caisse en douceur vers l'équilibre, ce qui correspond à un amortissement critique.
\end{enumerate}
\textbf{Barème indicatif (12 points) :} équation caractéristique et racine double : 2 pts ; solution générale : 1 pt ; détermination de $A$ et $B$ : 2 pts ; dérivée : 2 pts ; tableau de variations : 2 pts ; discriminant du réglage A et forme de la solution : 2 pts ; conclusion rédigée et justifiée : 1 pt.

\textbf{Points de vigilance :} bien dériver $(A + Bt)e^{-2t}$ comme un produit ; justifier le signe de $x'$ en citant $e^{-2t} > 0$ ; pour le réglage A, expliciter $\alpha$ et $\beta$ avant d'écrire la forme trigonométrique ; la conclusion doit s'appuyer sur la comparaison des deux formes de solutions.
\end{corrige}

\begin{corrige}[Exercice 11 — Oscillations d'un vibrateur industriel]
\begin{enumerate}
\item $(E_0) : y'' + 4y = 0$. Équation caractéristique : $r^2 + 4 = 0$, racines $r = \pm 2i$ ($\alpha = 0$, $\beta = 2$). Solution générale :
\[ y_0(t) = A\cos(2t) + B\sin(2t). \]
\item Si $y_p = 2$, alors $y_p' = 0$ et $y_p'' = 0$, donc $y_p'' + 4y_p = 0 + 8 = 8$ : $y_p$ est bien solution de $(E)$.
\item Solution générale de $(E)$ : $y(t) = A\cos(2t) + B\sin(2t) + 2$. Conditions initiales : $y(0) = A + 2 = 0$ donne $A = -2$ ; $y'(t) = -2A\sin(2t) + 2B\cos(2t)$, donc $y'(0) = 2B = 0$, soit $B = 0$. L'unique solution est
\[ y(t) = 2 - 2\cos(2t). \]
\item Pour tout réel $t$ : $y(t + \pi) = 2 - 2\cos(2t + 2\pi) = 2 - 2\cos(2t) = y(t)$, car le cosinus est $2\pi$-périodique. Donc $y$ est \cle{périodique} de période $T = \pi$.
\item $y$ est maximale lorsque $\cos(2t) = -1$, c'est-à-dire $2t = \pi + 2k\pi$, soit $t = \dfrac{\pi}{2} + k\pi$. Sur $[0\,;\,2\pi]$ : $t = \dfrac{\pi}{2}$ et $t = \dfrac{3\pi}{2}$. La position maximale est $y = 2 - 2 \times (-1) = 4$ mm.
\item Courbe de $y(t) = 2 - 2\cos(2t)$ sur $[0\,;\,2\pi]$ :
\begin{popfigure}
\begin{tikzpicture}
\begin{axis}[width=\linewidth,height=6cm,axis lines=middle,xlabel={$t$ (s)},ylabel={$y$ (mm)},xmin=0,xmax=6.6,ymin=-0.4,ymax=4.6,xtick={1.571,3.142,4.712,6.283},xticklabels={$\frac{\pi}{2}$,$\pi$,$\frac{3\pi}{2}$,$2\pi$},ytick={1,2,3,4},samples=200]
\addplot[popcurve,domain=0:6.283]{2-2*cos(deg(2*x))};
\end{axis}
\end{tikzpicture}
\legende{Courbe de $y(t) = 2 - 2\cos(2t)$ sur $[0\,;\,2\pi]$ : deux périodes complètes.}
\end{popfigure}
\item Le vibrateur oscille autour de la position d'équilibre $y = 2$ mm (la solution particulière constante). L'écart maximal à cette position est $4 - 2 = 2$ mm : l'\cle{amplitude} des oscillations est de $2$ mm. Le mouvement est périodique, de période $\pi$ secondes, sans amortissement : le vibrateur monte et descend indéfiniment entre $0$ et $4$ mm.
\end{enumerate}
\textbf{Barème indicatif (12 points) :} résolution de $(E_0)$ : 2 pts ; vérification de $y_p$ : 1 pt ; solution générale : 1 pt ; conditions initiales : 2 pts ; périodicité démontrée : 2 pts ; instants et maximum : 2 pts ; courbe correcte (échelles, allure) : 1 pt ; interprétation rédigée : 1 pt.

\textbf{Points de vigilance :} pour la périodicité, écrire explicitement $y(t + T) = y(t)$ en utilisant $\cos(\theta + 2\pi) = \cos\theta$ ; la période de $\cos(2t)$ est $\dfrac{2\pi}{2} = \pi$ et non $2\pi$ ; sur le graphique, respecter les unités demandées et marquer les extrema ; l'interprétation doit citer la position d'équilibre et l'amplitude.
\end{corrige}

\newpage

\coursec{Fiche bilan}

\begin{popfigure}
\centering
\begin{tikzpicture}[
    mindmap,
    concept color=popblue,
    level 1/.style={level distance=3.5cm, sibling angle=60},
    level 2/.style={level distance=2.5cm},
    every node/.style={font=\small, text=popdark},
    branch1/.style={concept color=poporange, grow=90},
    branch2/.style={concept color=popgreen, grow=30},
    branch3/.style={concept color=poppurple, grow=-30},
    branch4/.style={concept color=poppink, grow=-90},
    branch5/.style={concept color=popgold, grow=-150},
    branch6/.style={concept color=popblue, grow=150}
]
\node[concept, minimum size=2.2cm, font=\normalsize\bfseries, align=center]{Équations\\Différentielles}
    child[branch1] { node[concept, minimum size=1.8cm, align=center]{Vocabulaire\\Ordre, CI, SG, SP} }
    child[branch2] { node[concept, minimum size=1.8cm, align=center]{Ordre 1\\$y' = ay$\\$y = Ce^{ax}$} }
    child[branch3] { node[concept, minimum size=1.8cm, align=center]{Ordre 2 Homogène\\$ay''+by'+cy=0$\\Éq. caractéristique} }
    child[branch4] { node[concept, minimum size=1.8cm, align=center]{Cas $\Delta$\\$>0$: $e^{r_1x}, e^{r_2x}$\\$=0$: $e^{rx}, xe^{rx}$\\$<0$: $e^{\alpha x}\cos,\sin$} }
    child[branch5] { node[concept, minimum size=1.8cm, align=center]{Ordre 2 Non homogène\\$ay''+by'+cy=d$\\$y = y_h + y_p$\\$y_p = \frac{d}{c}$} }
    child[branch6] { node[concept, minimum size=1.8cm, align=center]{Applications\\RC, RL, Mécanique\\Croissance, Chute} };
\end{tikzpicture}
\legende{Carte mentale : synthèse de la leçon sur les équations différentielles}
\end{popfigure}

\begin{bilan}

\courssub{Définitions et vocabulaire clés}

\begin{itemize}
    \item \textbf{Équation différentielle (ED)} : Relation liant une fonction inconnue $y$, sa variable $x$ et ses dérivées successives.
    \item \textbf{Ordre} : Ordre de la dérivée la plus élevée qui apparaît.
    \item \textbf{Solution générale (SG)} : Famille de toutes les solutions, dépendant de constantes arbitraires (autant que l'ordre).
    \item \textbf{Solution particulière (SP)} : Solution obtenue en fixant les constantes (via des \textbf{conditions initiales} : $y(x_0), y'(x_0), \dots$).
    \item \textbf{Équation homogène} : Second membre nul ($=0$). \textbf{Non homogène} : Second membre non nul.
\end{itemize}

\courssub{Formules et théorèmes à connaître par cœur}

\begin{center}
\begin{tabularx}{\linewidth}{|>{\bfseries}l|X|X|}
\hline
\rowcolor{popblueL}
\textbf{Type d'équation} & \textbf{Résolution} & \textbf{Solution Générale} \\
\hline
$y' = ay$ \newline ($a \in \mathbb{R}^*$) & Intégration directe ou recherche exponentielle & $y(x) = C e^{ax}$ \newline ($C \in \mathbb{R}$) \\
\hline
$ay'' + by' + cy = 0$ \newline ($a \neq 0$) & Équation caractéristique : $ar^2 + br + c = 0$ \newline $\Delta = b^2 - 4ac$ & 
\begin{itemize}[nosep, leftmargin=*]
\item $\Delta > 0$ : $r_1 \neq r_2$ réels $\Rightarrow y = C_1 e^{r_1 x} + C_2 e^{r_2 x}$
\item $\Delta = 0$ : $r$ double $\Rightarrow y = (C_1 + C_2 x) e^{rx}$
\item $\Delta < 0$ : $r = \alpha \pm i\beta$ $\Rightarrow y = e^{\alpha x}(C_1 \cos \beta x + C_2 \sin \beta x)$
\end{itemize} \\
\hline
$ay'' + by' + cy = d$ \newline ($d \in \mathbb{R}^*$, $c \neq 0$) & 1. Résoudre homogène $\to y_h$ \newline 2. Chercher $y_p = k$ (cste) $\Rightarrow ck = d$ & $y(x) = y_h(x) + \dfrac{d}{c}$ \\
\hline
\end{tabularx}
\end{center}

\courssub{Méthodes express (étapes)}

\begin{methode}[Équation $y' = ay$]
\begin{enumerate}
    \item Écrire sous la forme $y' - ay = 0$.
    \item Poser $y = e^{ax}$ (ou intégrer $\frac{y'}{y} = a$).
    \item Écrire SG : $y = C e^{ax}$.
    \item Utiliser condition initiale (si donnée) pour trouver $C$.
\end{enumerate}
\end{methode}

\begin{methode}[Équation $ay'' + by' + cy = 0$]
\begin{enumerate}
    \item Écrire l'équation caractéristique : $ar^2 + br + c = 0$.
    \item Calculer $\Delta = b^2 - 4ac$.
    \item Selon le signe de $\Delta$, écrire la SG correspondante (cf. tableau).
    \item Appliquer les conditions initiales ($y(x_0), y'(x_0)$) pour déterminer $C_1, C_2$ (système 2x2).
\end{enumerate}
\end{methode}

\begin{methode}[Équation $ay'' + by' + cy = d$ ($c \neq 0$)]
\begin{enumerate}
    \item Résoudre l'équation homogène associée $\to y_h$.
    \item Trouver une solution particulière constante : $y_p = \frac{d}{c}$.
    \item SG = $y_h + y_p$.
    \item Appliquer conditions initiales sur $y$ et $y'$ (attention : $y_p' = 0$).
\end{enumerate}
\end{methode}

\begin{attention}[Pièges fréquents]
\begin{itemize}
    \item Oublier le cas $\Delta = 0$ (solution $xe^{rx}$) ou $\Delta < 0$ (forme trigonométrique).
    \item Confondre $y_p$ pour $ay''+by'+cy=d$ : c'est $\frac{d}{c}$ \textbf{uniquement si $c \neq 0$}. Si $c=0$, chercher $y_p = kx$.
    \item Ne pas dériver correctement $y_h$ pour appliquer la condition sur $y'(x_0)$ (produit, chaîne).
    \item Oublier de vérifier que la solution trouvée vérifie bien l'équation de départ.
\end{itemize}
\end{attention}

\end{bilan}

\begin{aretenir}[Checklist avant le Probatoire]
\begin{itemize}
    \item $\square$ Je sais définir l'ordre d'une ED et donner sa solution générale / particulière.
    \item $\square$ Je résous automatiquement $y' = ay \Rightarrow y = Ce^{ax}$.
    \item $\square$ J'écris l'équation caractéristique $ar^2+br+c=0$ sans erreur pour $ay''+by'+cy=0$.
    \item $\square$ Je distingue les 3 cas selon $\Delta$ ($>0, =0, <0$) et j'écris la SG correspondante.
    \item $\square$ Je sais dériver une solution du type $e^{\alpha x}(C_1\cos\beta x + C_2\sin\beta x)$ (produit + chaîne).
    \item $\square$ Pour $ay''+by'+cy=d$ ($c\neq0$), je pose $y_p = \frac{d}{c}$ et j'ajoute $y_h$.
    \item $\square$ Je sais monter le système 2x2 pour trouver $C_1, C_2$ à partir de $y(x_0)$ et $y'(x_0)$.
    \item $\square$ Je vérifie que ma solution particulière convient bien (cas $c=0 \Rightarrow y_p=kx$).
    \item $\square$ Je modélise un problème concret (RC, RL, chute, croissance) en ED d'ordre 1 ou 2.
    \item $\square$ J'interprète physiquement les constantes (amortissement, pulsation, temps caractéristique).
    \item $\square$ Je rédige une conclusion claire : « La solution cherchée est $y(x) = \dots$ ».
\end{itemize}
\end{aretenir}

\findecours

\end{document}
